\documentclass[JEP,XML,SOM,Unicode,NoFloatCountersInSection]{cedram}
\datereceived{2017-04-28}
\dateaccepted{2019-04-17}
\dateepreuves{2019-04-29}

\usepackage{verbatim}
\RequirePackage{eucal}
\usepackage{mathrsfs}
\let\mathcal\mathscr

\newenvironment{enumeratea}
{\bgroup\def\theenumi{\alph{enumi}}\begin{enumerate}}
{\end{enumerate}\egroup}
\newenvironment{enumerateap}
{\bgroup\def\theenumi{\alph{enumi}$'$}\begin{enumerate}}
{\end{enumerate}\egroup}
\newcommand{\RedefinitSymbole}[1]{%
\expandafter\let\csname old\string#1\endcsname=#1
\let#1=\relax
\newcommand{#1}{\csname old\string#1\endcsname\,}%
}
\RedefinitSymbole{\forall} \RedefinitSymbole{\exists}

\newcommand{\mto}{\mathchoice{\longmapsto}{\mapsto}{\mapsto}{\mapsto}}
\newcommand{\hto}{\mathrel{\lhook\joinrel\to}}
\newcommand{\To}[1]{\mathchoice{\xrightarrow{\textstyle\kern4pt#1\kern3pt}}{\stackrel{#1}{\longrightarrow}}{}{}}
\newcommand{\Hto}[1]{\mathrel{\lhook\joinrel\To{#1}}}

\let\oldvee\vee
\renewcommand{\vee}{{\scriptscriptstyle\oldvee}}
\newcommand{\spfrac}[2]{\sfrac{#1}{(#2)}}
\newcommand{\psfrac}[2]{\sfrac{(#1)}{#2}}
\newcommand{\Psfrac}[2]{(\sfrac{#1}{#2})}
\newcommand{\sbullet}{{\scriptscriptstyle\bullet}}
\newcommand{\scbullet}{{\raisebox{1pt}{$\sbullet$}}}
\newcommand{\Changel}{\mathcode`l="7160}
\newcommand{\NoChangel}[1]{%
\expandafter\let\csname old\string#1\endcsname=#1
\let#1=\relax
\newcommand{#1}{\mathcode`l="716C\csname old\string#1\endcsname\mathcode`l="7160 }%
}

\usepackage{multirow}

\let\cites\cite
\newcommand{\citestar}[3]{\cite[#3]{#1}}

\usepackage{pdflscape}
\usepackage{afterpage}

\usepackage{dynkin-diagrams}

\usepackage{mathtools}

\newcommand{\Leaf}{\EuScript{L}}

\theoremstyle{plain}
\newtheorem{theo}{Theorem}[section]
\newtheorem{lemm}[theo]{Lemma}
\newtheorem{prop}[theo]{Proposition}
\newtheorem{coro}[theo]{Corollary}

\newtheorem{theoA}{Theorem}

\renewcommand*{\thetheoA}{\Alph{theoA}}

\newtheorem{coroA}[theoA]{Corollary}
\newtheorem{propA}[theoA]{Proposition}

\newtheorem*{theorem*}{Theorem }

\theoremstyle{definition}
\newtheorem{defi}[theo]{Definition}
\newtheorem{fact}[theo]{Fact}
\newtheorem{nota}[theo]{Notation}
\newtheorem{exem}[theo]{Example}
\newtheorem{rema}[theo]{Remark}
\newtheorem*{rema*}{Remark}

\let\cst c
\let\ov\overline
\let\la\langle
\let\ra\rangle
\let\rgy r

\DeclareMathOperator{\ad}{ad}
\DeclareMathOperator{\diag}{diag}
\DeclareMathOperator{\End}{End}
\DeclareMathOperator{\rF}{F}
\DeclareMathOperator{\Fix}{Fix}
\DeclareMathOperator{\GL}{GL}
\DeclareMathOperator{\Gr}{Gr}
\DeclareMathOperator{\Hom}{Hom}
\DeclareMathOperator{\im}{Im}\let\Im\im
\DeclareMathOperator{\Isom}{Isom}
\DeclareMathOperator{\Ker}{Ker}
\DeclareMathOperator{\Orb}{Orb}
\DeclareMathOperator{\Pic}{Pic}
\DeclareMathOperator{\rk}{rk}
\DeclareMathOperator{\rS}{S}
\DeclareMathOperator{\SL}{SL}
\DeclareMathOperator{\SO}{SO}
\DeclareMathOperator{\SU}{SU}
\DeclareMathOperator{\Sp}{Sp}
\DeclareMathOperator{\Spin}{Spin}
\DeclareMathOperator{\Stab}{Stab}
\DeclareMathOperator{\supp}{supp}
\DeclareMathOperator{\rU}{U}
\DeclareMathOperator{\vol}{vol}

\newcommand{\rtf}{\mathrm{tf}}

\newcommand{\cF}{\mathcal{F}}
\newcommand{\cG}{\mathcal{G}}
\newcommand{\cC}{\mathcal{C}}
\newcommand{\cE}{\mathcal{E}}
\newcommand{\cL}{\mathcal{L}}
\newcommand{\cM}{\mathcal{M}}
\newcommand{\cO}{\mathcal{O}}
\newcommand{\cR}{\mathcal{R}}
\newcommand{\cS}{\mathcal{S}}
\newcommand{\cT}{\mathcal{T}}
\newcommand{\cU}{\mathcal{U}}
\newcommand{\cV}{\mathcal{V}}
\newcommand{\cX}{\mathcal{X}}
\newcommand{\cW}{\mathcal{W}}

\newcommand{\bA}{\mathbb{A}}
\newcommand{\bB}{\mathbb{B}}
\newcommand{\bC}{\mathbb{C}}
\newcommand{\bE}{\mathbb{E}}
\newcommand{\bF}{\mathbb{F}}
\newcommand{\bH}{\mathbb{H}}
\newcommand{\bO}{\mathbb{O}}
\newcommand{\bP}{\mathbb{P}}
\newcommand{\bS}{\mathbb{S}}
\newcommand{\bU}{\mathbb{U}}
\newcommand{\bV}{\mathbb{V}}
\newcommand{\bCP}{\mathbb{CP}}
\newcommand{\bZ}{\mathbb{Z}}
\newcommand{\bR}{\mathbb{R}}
\newcommand{\bQ}{\mathbb{Q}}

\newcommand{\HC}{{\mathbb{H}^n_\mathbb{C}}}
\newcommand{\THC}{T_{\HC}}

\newcommand{\PTX}{{\mathbb{P}T_X}}
\newcommand{\PTH}{\mathbb{P}T_{\mathbb{H}^n_\mathbb{C}}}
\newcommand{\bPTM}{{\mathbb{P}T_M}}
\newcommand{\OX}{\mathcal{O}_{\PTX}(-1)}
\newcommand{\Vr}{\mathbb{V}^\rgy}
\newcommand{\Wr}{\mathbb{W}^\rgy}
\newcommand{\Vb}{\mathbb{V}^{\rgb}}
\newcommand{\Vkr}{\mathbb{V}^{kr}}

\newcommand{\by}{\boldsymbol{y}}
\newcommand{\bfG}{\boldsymbol{G}}

\newcommand{\fe}{\mathfrak{e}}
\newcommand{\fgg}{\mathfrak{g}}
\newcommand{\fh}{\mathfrak{h}}
\newcommand{\fJ}{\mathfrak{J}}
\newcommand{\fk}{\mathfrak{k}}
\newcommand{\fl}{\mathfrak{l}}
\newcommand{\fm}{\mathfrak{m}}
\newcommand{\fn}{\mathfrak{n}}
\newcommand{\fp}{\mathfrak{p}}
\newcommand{\fq}{\mathfrak{q}}
\newcommand{\ft}{\mathfrak{t}}
\newcommand{\fz}{\mathfrak{z}}
\newcommand{\fsl}{\mathfrak{sl}}
\newcommand{\fso}{\mathfrak{so}}
\newcommand{\fsp}{\mathfrak{sp}}
\newcommand{\fsu}{\mathfrak{su}}
\newcommand{\fstab}{\mathfrak{stab}}
\newcommand{\fu}{\mathfrak{u}}

\newcommand{\suma}[1]{\alpha_1^\vee + \cdots + \alpha_{#1}^\vee}
\newcommand{\av}{\alpha^\vee}

\newcommand{\zmax}{z_{\mathtt{max}}}
\newcommand{\Zmax}{Z_{\mathtt{max}}}
\newcommand{\Ek}{\bE^{\otimes k}}
\newcommand{\cEk}{\cE^{\otimes k}}

\newcommand{\rkm}{{r_M}}
\newcommand{\rky}{{r(y)}}
\newcommand{\rgb}{{\rk\beta}}

\newcommand{\ovLrR}{\ov{L}_{r,\bR}}
\newcommand{\rhoc}{\rho_\mathrm{cpt}}
\newcommand{\flb}{{\ov \fl_r}}
\newcommand{\LB}{\ov L_r}
\newcommand{\ovE}{\ov{E}}
\newcommand{\ovb}{\ov{\beta}}
\newcommand{\ovg}{\ov{\gamma}}
\newcommand{\ovt}{\ov{\theta}}

\renewcommand{\iff}{if and only if }

\newcommand{\scal}[1]{\langle #1 \rangle}

\NoChangel{\length}
\NoChangel{\log}
\NoChangel{\ln}
\NoChangel{\lim}
\NoChangel{\limsup}
\NoChangel{\liminf}
\NoChangel{\varlimsup}
\NoChangel{\varliminf}
\NoChangel{\varinjlim}
\NoChangel{\varprojlim}
\NoChangel{\vol}
\NoChangel{\fl}
\NoChangel{\fsl}

\begin{document}
\frontmatter

\title[Maximal representations, a uniform approach]{Maximal representations\\ of cocompact
complex hyperbolic lattices,\\ a uniform approach}

\author[\initial{P.-E.} \lastname{Chaput}]{\firstname{Pierre-Emmanuel} \lastname{Chaput}}
\address{Institut Élie Cartan de Lorraine,
Université de Lorraine\\
Site de Nancy, B.P. 70239, F-54506 Vandoeuvre-lès-Nancy Cedex}
\email{pierre-emmanuel.chaput@univ-lorraine.fr}
\urladdr{http://www.iecl.univ-lorraine.fr/~Pierre-Emmanuel.Chaput/}

\author[\initial{J.} \lastname{Maubon}]{\firstname{Julien} \lastname{Maubon}}
\address{Institut Élie Cartan de Lorraine,
Université de Lorraine\\
Site de Nancy, B.P. 70239, F-54506 Vandoeuvre-lès-Nancy Cedex}
\email{julien.maubon@univ-lorraine.fr}

\begin{abstract}
We complete the classification of maximal representations of cocompact complex hyperbolic lattices in Hermitian Lie groups by dealing with the exceptional groups $\mathrm{E}_{6(-14)}$ and $\mathrm{E}_{7(-25)}$. We prove that if $\rho$ is a maximal representation of a cocompact complex hyperbolic lattice $\Gamma\subset\mathrm{SU}(1,n)$, $n>1$, in an exceptional Hermitian group $G_\mathbb{R}$, then $n=2$ and $G_\mathbb{R}=\mathrm{E}_{6(-14)}$, and we describe completely the representation $\rho$. The case of classical Hermitian target groups was treated by Vincent Koziarz and the second named author \cite{KM}. However we do not focus immediately on the exceptional cases and instead we provide a more unified perspective, as independent as possible of the classification of the simple Hermitian Lie groups. This relies on the study of the cominuscule representation of the complexification $G$ of the target group $G_\mathbb{R}$. As a by-product of our methods, when the target Hermitian group $G_\mathbb{R}$ has tube type, we obtain an inequality on the Toledo invariant of the representation $\rho:\Gamma\to G_\mathbb{R}$ which is stronger than the Milnor-Wood inequality (thereby excluding maximal representations in such groups).
\end{abstract}

\alttitle{\selectlanguage{french}Représentations maximales des réseaux hyperboliques complexes cocompacts: une approche unifiée}

\begin{altabstract}
Nous complétons la classification des représentations maximales des réseaux hyperboliques complexes dans les groupes de Lie hermitiens en traitant le cas des groupes exceptionnels $\mathrm{E}_{6(-14)}$ et $\mathrm{E}_{7(-25)}$. Nous montrons que si $\rho$ est une représentation maximale d'un réseau hyper\-bo\-lique complexe cocompact $\Gamma \subset \mathrm{SU}(1,n)$, avec $n>1$, dans un groupe hermitien $G_\mathbb{R}$ de type exceptionnel, alors $n=2$ et $G_\mathbb{R} = \mathrm{E}_{6(-14)}$, et nous décrivons complètement la représentation $\rho$. Le cas des groupes hermitiens classiques avait été traité par Vincent Koziarz et le deuxième auteur cité \cite{KM}. Cependant, nous ne nous restreignons pas immédiatement aux groupes exceptionnels: nous proposons au contraire une approche unifiée, aussi indépendante que possible de la classification des groupes de Lie hermitiens simples. Cette approche repose sur une étude de la représentation cominuscule de la complexification du groupe d'arrivée $G_\mathbb{R}$. Dans le cas où $G_\mathbb{R}$ est de type tube, nos méthodes permettent en particulier d'établir une inégalité sur l'invariant de Toledo de la représentation $\rho:\Gamma\rightarrow G_\mathbb{R}$ qui est plus forte que l'inégalité de Milnor-Wood et qui exclut donc la possibilité d'une représentation maximale pour de tels groupes.\end{altabstract}

\subjclass{53C35, 22E40, 32L05, 32Q15, 17B10, 20G05}

\keywords{Complex hyperbolic lattices, Milnor-Wood inequality, maximal representations, cominuscule representations, exceptional Lie groups, harmonic Higgs bundles, holomorphic foliations}

\altkeywords{Réseaux hyperboliques complexes, inégalité de Milnor-Wood, représentations maximales, représentations cominuscules, groupes de Lie exceptionnels, fibrés de Higgs harmoniques, feuilletages holomorphes}
\maketitle

\newpage
\tableofcontents
\mainmatter

\bgroup
\Changel

\section{Introduction}

This paper deals with maximal representations of complex hyperbolic lattices in semisimple Hermitian
Lie groups with no compact factors.

A \emph{complex hyperbolic lattice} $\Gamma$ is a lattice in the Lie group $\SU(1,n)$, a finite
cover of the group of biholomorphisms of the $n$-dimensional complex hyperbolic space
$\HC=\SU(1,n)/\rU(n)$.
We shall always
assume that our lattice $\Gamma$ is \emph{cocompact}, or \emph{uniform}, meaning that the quotient
$X:=\Gamma\backslash\HC$ is
compact. To simplify matters, we also assume in this introduction (except in the statements of
our main results) that $\Gamma$ is torsion free, so that $X$ is also a
manifold. The
$\SU(1,n)$-invariant Kähler form on $\HC$ with constant holomorphic sectional curvature $-1$ and the
Kähler form it induces on~$X$ will both be
denoted by $\omega$.

A Lie group $G_\bR$ is said to be a \emph{real algebraic Hermitian Lie group} (or a \emph{Hermitian
group} for short) if it is the connected component $\bfG(\bR)^\circ$ of the
group of real point of an algebraic group $\bfG$ defined over $\bR$, if it is semisimple
with no compact factors, and if its
associated symmetric space $M=G_\bR/K_\bR$ is a Hermitian symmetric space
(of the noncompact type).
This means that the noncompact symmetric space $M$ admits a
$G_\bR$-invariant complex structure, with respect to which the $G_\bR$-invariant Riemannian metrics
are (necessarily) Kähler. The real rank $\rk_\bR G_\bR$ of
$G_\bR$ coincides with the rank~$\rkm$ of $M$ as a symmetric space, namely the maximal dimension of
a flat subspace in $M$. Simple Hermitian groups are classified: there are four
infinite families of \emph{classical} groups, which up to isogeny are $\SU(p,q)$ with $1\leq p\leq
q$, $\SO_0(2,p)$ with $p\geq 3$, $\Sp(2m,\bR)$ with $m\geq 2$ and
$\SO^\star(2m)$ with $m\geq 4$; and two \emph{exceptional} groups $\mathrm{E}_{6(-14)}$ and
$\mathrm{E}_{7(-25)}$. The real ranks of these groups are respectively
$p$, $2$, $m$, $\lfloor \sfrac m2 \rfloor$, $2$ and $3$.

Let $\omega_M$ denote the $G_\bR$-invariant Kähler form of
$M$, uniquely normalized so that its holomorphic sectional curvatures lie between $-1$ and
$-\sfrac 1{\rkm}$.

If $\rho$ is a representation (a group homomorphism) from $\Gamma$ to $G_\bR$, we
define its \emph{Toledo invariant} $\tau(\rho)$ as follows:
$$
\tau(\rho)=\frac 1{n!}\int_X f^\star\omega_M\wedge\omega^{n-1},
$$
where $f:\HC\to M$ is a $\cC^\infty$ and $\rho$-equivariant map and $f^\star\omega_M$ is seen as a 2-form on $X$ by
$\Gamma$-invariance. The Toledo invariant does not depend on the choice of the map $f$, it depends only
on $\rho$, and in fact, only on the connected component of $\rho$ in $\Hom(\Gamma,G_\bR)$. Moreover, it
satisfies the following Milnor-Wood type inequality:
$$
|\tau(\rho)|\leq \rkm\vol(X),
$$
a fundamental property established in full generality in~\cite{BI07}.

\emph{The maximal representations} $\rho:\Gamma\to G_\bR$ are those representations
for which the Milnor-Wood inequality is an equality.

In \cite{KM}, to the introduction of which we refer for a history of the subject, Koziarz and the
second named author classified maximal representations when $n\geq 2$, \ie $\Gamma$ is not a surface
group, and when $G_\bR$ is a classical group.
Their proof uses the standard representations of the classical groups $G_\bR$ and their
complexifications $G:=\bfG(\bC)$ and is therefore quite dependent on the classification of simple
Hermitian groups.
In the present work, using in a uniform way the \emph{cominuscule} representation of $G$
we give a new proof of the Milnor-Wood inequality relying on
algebraic properties of this cominuscule representation shared by all Hermitian groups,
see the proof of Theorem~\ref{theo:MW}.
This proof provides information which,
in the case of maximal representations, allows us to extend
the classification to all target Hermitian groups.

\begin{theoA}
\label{maintheo}
Let $\Gamma$ be a uniform lattice in $\SU(1,n)$, $n\geq 2$.
Let $G_\bR$ be a real algebraic Hermitian Lie group and let $M$ be the Hermitian symmetric space of
the noncompact type associated with $G_\bR$. If $\rho$ is a maximal
representation of $\Gamma$ in $G_\bR$, then:
\begin{itemize}
\item each simple factor of $G_\bR$ is isogenous either to
$\SU(p,q)$ for some $(p,q)$ with $q\geq n\, p$,
or to the exceptional group $\mathrm{E}_{6(-14)}$, the latter being possible if and only if $n=2$;
\item if $\tau(\rho)>0$, there exists a
$\rho$-equivariant holomorphic map $f:\HC\to M$. Moreover it satisfies
$f^\star\omega_M=\rkm\,\omega$;
\item if $\tau(\rho)<0$, there exists a $\rho$-equivariant antiholomorphic map $f:\HC\to
M$. Moreover it satisfies $f^\star\omega_M=-\rkm\,\omega$;
\item in both cases, the map $f$ is unique as it is the unique $\rho$-equivariant harmonic map $\HC\to
M$. It
is a totally geodesic embedding. It is also equivariant w.r.t.\
a uniquely defined homomorphism of Lie groups $\varphi: \SU(1,n)\to G_\bR$.
\end{itemize}
\end{theoA}

We can also deduce a complete structure result
for the maximal representation~$\rho$, namely that $\rho$ is essentially the restriction to $\Gamma$ of
the morphism of Lie
groups \hbox{$\varphi: \SU(1,n) \to G_\bR$} given by Theorem~\ref{maintheo}.

\begin{coroA}
\label{maincoro}
Under the assumptions of Theorem~\ref{maintheo}, the representation $\rho$ is reduc\-tive, discrete, faithful, and
acts cocompactly on the image of $f$ in $M$. The centralizer~$Z_\bR$ of the image of $\varphi$ in
$G_\bR$ is compact and
there exists a group morphism $\rhoc: \Gamma \to\nobreak Z_\bR$ such that
$$
\forall \gamma \in \Gamma, \quad\rho(\gamma) = \varphi(\gamma)\rhoc(\gamma) =
\rhoc(\gamma)\varphi(\gamma).
$$
\end{coroA}

The group $Z_\bR$ is described in Lemma~\ref{lemm:Z}. Moreover, Lemma~\ref{lemm:fixator} says that $Z_\bR$ is exactly the subgroup of $G_\bR$ fixing $f(\HC)$
pointwise. Just
before Lemma~\ref{lemm:Z}, we also describe the
morphism $\varphi:\SU(1,n)\to G_\bR$. If $G_\bR$ is simple, and up to conjugating $\rho$ by an
element in $G_\bR$, it is always given by a diagonal embedding
$\SU(1,n)\hto\SU(\rkm,n\rkm)\hto G_\bR$, where
$\SU(\rkm,n\rkm)$ is a subgroup of $G_\bR$ stable under conjugation by elements of the maximal torus.

The global strategy we
adopt here is the same as in~\cite{KM}.
First, we use known results (see~\ref{subsec:Higgs}) to get a harmonic $\rho$-equivariant
map $f:\HC\to M$ and a Higgs bundle $(\bar E,\bar\theta)$ on the quotient
$X=\Gamma\backslash\HC$ associated
to a (reductive) representation \hbox{$\rho:\Gamma\to G_\bR$}. Then,
we translate the
Milnor-Wood inequality into an inequality involving degrees of subbundles of $\bar E$
(see~\ref{subsec:reword}). This inequality is then proved (in~\ref{subsec:proofMW}) using the
Higgs-stability properties of
$(\bar E,\bar\theta)$,
or rather the \emph{leafwise} Higgs-stability properties of the pull-back $(E,\theta)$ of
$(\bar E,\bar\theta)$ to the
projectivized tangent bundle $\PTX$ of $X$ with respect to the \emph{tautological foliation} on $\PTX$
(see Sections~\ref{subsec:HiggsFol} and~\ref{subsec:toto}).

Although classical target groups were already treated in \cite{KM}, we decided not to focus
immediately on the exceptional cases and instead to provide a
more unified perspective, as independent as possible of the
classification of the simple Hermitian
Lie groups, in the spirit of \cite{biquard}. To achieve this, instead of considering the Higgs
vector bundle associated with the standard representation of the complexification $G$ of~$G_\bR$
(which is only defined in the classical cases), we
work with the Higgs
vector bundle $(E,\theta)$ associated with the \emph{cominuscule} representation
$\bE$ of $G$ defined by $M$ (and a choice of invariant complex structure),
see Section~\ref{subsec:minuscule-rep}. This representation is such that
$M$ is holomorphically embedded as a locally closed subset of
the projectivization $\bP\bE$ of $\bE$: this is sometimes called the \emph{first canonical
embedding} of the compact dual $\check{M}$, see \citestar{nakagawa}*{p.\,651}.

On the algebraic side, we present in
Section~\ref{subsec:higgs-submodule} a general
construction of a self-dual graded subspace $\Vr$ of $\bE$ associated with an element of
$\fm^\pm$ of rank $r$ (here,
$\fgg = \fk \oplus (\fm^+ \oplus \fm^-)$ is the Cartan decomposition of the Lie algebra of
$G$). The subspace~$\Vr$ is in fact a module under the action of a complex subgroup $\ov{L}_r$
of $G$, and the self-duality of $\Vr$ corresponds to the fact that $\ov{L}_r$ admits a noncompact
real form $\ov{L}_{r,\bR}$ which is a simple (if $G_\bR$ is simple) Hermitian Lie group whose
associated symmetric space
$M_r$ is of \emph{tube type} and has rank $r$, see Section~\ref{subsubsec:tube}.
On the geometric side, this construction produces a leafwise
Higgs subsheaf $\cV$ of $(E,\theta)\to\PTX$ associated with the holomorphic component
of the Higgs field $\theta$ (see Section~\ref{subsec:HiggsSubsheaf}), whose existence is then used to
prove the
Milnor-Wood inequality. To be a bit more precise, on a \emph{generic} fiber of $(E,\theta)\to\PTX$, the
leafwise Higgs subsheaf we define admits a
purely representation theoretic description. The algebraic counterparts of the generic objects
are first introduced and studied in
Section~\ref{sec:subnil}. This is then used in Section~\ref{subsec:HiggsSubsheaf} to define the subsheaf
and prove that it has the desired properties.

One interesting by-product of this unified approach is that it allows to exclude a priori the
possibility of maximal representations in any tube type
real algebraic Hermitian Lie group, and in particular in $\mathrm{E}_{7(-25)}$. Recall
that up to isogeny the simple tube type
Hermitian groups are $\SU(p,p)$, $\SO_0(2,p)$, $\Sp(2m,\bR)$, $\SO^\star(2m)$ with~$m$ even,
and $\mathrm{E}_{7(-25)}$. See also Section~\ref{subsubsec:tube}. Indeed,
we prove in Section~\ref{subsec:tube} that for tube type targets
the representation $\rho$ satisfies a stronger inequality than the Milnor-Wood inequality.

\begin{propA}
\label{mainprop}
Let $\Gamma$ be a uniform lattice in $\SU(1,n)$, with $n\geq 2$, and let $X=\Gamma\backslash\HC$.
Assume that the real algebraic Hermitian Lie group $G_\bR$ has tube type and let~$\rkm$ be the real rank of $G_\bR$. Let $\rho$ be a
representation $\Gamma\to G_\bR$. Then
$$ |\tau(\rho)|\,\leq\, \max\Bigl\{{\rkm-1,\frac \rkm{2}\cdot \frac{n+1}{n}}\Bigr\}\,
\vol(X)<\rkm \vol(X).$$
\end{propA}

When the representation $\rho$ is maximal, the results of \cite{KM} imply that the restric\-tion of the
leafwise Higgs subsheaf $\cV$ corresponding to the module $\bV^\rkm$
to almost all leaves $\Leaf$ of the tautological foliation on $\PTX$ is a Higgs subbundle of
$(E,\theta)_{|\Leaf}$. Moreover, the self-duality of $\bV^\rkm$ means that the
holomorphic subbundles $\cV_i$ composing $\cV=\bigoplus_{0\leq i\leq\rkm}\cV_i$ have the
symmetries of a weight $\rkm$ real variation of Hodge structure: $\ov{\cV}_{\rkm-i}\simeq\cV_i$.
This is the Higgs bundle analogue, in the higher dimensional setting we are in, of the fact that
maximal representations of surface groups stabilize maximal tube type subdomains in $M$ (see the
remark below). We also refer to Section~\ref{subsec:hodge} for a discussion of the cominuscule
representation $\bE$ and the submodule $\Vr$ from the Hodge theoretic point of view developed in \cites{Gross, ShengZuo}.

To prove Theorem~\ref{maintheo}, one then needs to deduce that the $\rho$-equivariant
harmonic map $f$ is (anti-)holomorphic from the facts we just mentioned. To this end, thanks to
Proposition~\ref{mainprop} and \cite{KM}, we only need to
deal with the case of $\mathrm{E}_{6 (-14)}$.
Maximal representations in this exceptional group are treated
in Sections~\ref{subsec:E6} and~\ref{subsec:proof} where we prove that they exist if and
only if $n=2$, in which case they are
essentially induced by a homomorphism $ \varphi: \SU(1,2)\to\SU(2,4)\to \mathrm{E}_{6 (-14)}$.

\begin{rema*}\label{ngeq2}
The assumption $n\geq 2$ is of course essential in all these results. Maximal representations of
lattices of $\SU(1,1)$, \ie surface groups, form a very interesting
and intensely studied subclass
of representations, see \eg \cites{Hernandez,Xia00,MarkmanXia, BGPG1, BGPG2,
BIW10,GuichardWienhard}. They are also reductive, faithful and discrete, and they
stabilize a tube type subdomain of rank $\rkm$ in $M$, but
they exist for all target Hermitian groups and they are not in
general ``induced'' by a homomorphism $\SU(1,1)\to G_\bR$. On the contrary they define rich moduli
spaces (in the case $G_\bR=\SU(1,1)$, this is the Teichmüller space).

Accordingly, if $n=1$, the first inequality of Proposition~\ref{mainprop} is nothing but the
Milnor-Wood inequality.
One might ask whether this inequality is sharp when $n\geq 2$ and what could be said about representations
achieving the optimal bound (whatever it may be). Note that this question is not restricted to tube
type targets. For instance if $G_\bR=\SU(p,q)$, Theorem~\ref{maintheo} says that the Milnor-Wood
inequality is sharp when $q \geq pn$, but says nothing about the
other cases (when $p=2$
and $q<2n$, a better bound is known \cite{KMrank2}).

The interesting (and probably difficult) question of finding the sharp bound for a given pair
$(n,G_\bR)$
for which no maximal representation exists, and of characterizing the
representations achieving this bound remains open.
\end{rema*}

\begin{rema*}
As we mentioned, our main results are true for lattices with torsion, see Remark~\ref{rema:torsion}.
The assumption that $\Gamma$ is uniform should not be necessary,
and we believe that our method can be adapted to cover the case of non uniform lattices, but we
decided to leave that aside for a future work.
\end{rema*}

\subsubsection*{Acknowledgements}
We thank the anonymous referees for their detailed and helpful
comments, which led to substantial improvements in the presentation of the article. We thank them in
particular for bringing the references \cites{Gross, ShengZuo} to our attention.

\section{Submodule of a cominuscule representation associated with a~nilpotent element}
\label{sec:subnil}

Here we develop the algebraic material that we will need in Section~\ref{sec:MW} to give a new and
unified proof of the Milnor-Wood inequality.

\subsection{Hermitian symmetric spaces}
\label{subsec:hermitian}

We begin by recalling useful facts and definitions concerning Hermitian symmetric spaces and setting
up our notation. Good references for what follows are \cites{WolfFineStructure,
AshMumfordRapoportTai, Helgason, Knapp}.

Let $G_\bR$ be a simple real algebraic Hermitian Lie group and $K_\bR$ a maximal compact subgroup of
$G_\bR$. By our convention, this means that
$G_\bR$ is a simple noncompact Lie group whose associated symmetric space $M=G_\bR/K_\bR$ is a Hermitian
symmetric space of the noncompact type ($M$ is irreducible
since $G_\bR$ is simple), and that moreover $G_\bR$ is the connected component of the group
of real points $\bfG(\bR)$ of an algebraic group $\bfG$ defined over $\bR$. Observe that for example the connected component
$\Isom^\circ(M)$ of the group of isometries of a Hermitian symmetric space of the noncompact type
$M$ is a real algebraic Hermitian Lie group, see \eg \citestar{AshMumfordRapoportTai}*{p.\,106}. The group of complex points $\bfG(\bC)$ of $\bfG$ will be denoted by $G$.

We shall first assume that the complex Lie group $G$ is simply connected. This assumption simplifies the
exposition of the paper and the
arguments of our proofs. We shall see in Remarks~\ref{rema:1-connected}
and~\ref{rema:1-connected-bis} how to deal with the
general case.

The \emph{rank} of the symmetric space $M$, or equivalently the \emph{real rank} of $G_\bR$, will be
denoted by $\rkm$. Recall that a $k$-dimensional \emph{flat} in $M$ is the image of a totally geodesic
isometric embedding of $\bR^k$ in $M$. By definition the rank $\rkm$ of $M$ is the maximal dimension of
a flat in $M$. Maximal flats are conjugated under $G_\bR$ and maximal flats through the origin
$eK_\bR$ are conjugated under $K_\bR$. Equivalently, in the Hermitian setting,
one can consider \emph{polydiscs} in $M$, that is,
images of totally geodesic homomorphic embeddings of a polydisc $\Delta^k$ in $M$. Then $\rkm$ is also
the maximal (complex) dimension of a polydisc in $M$ (maximal polydiscs
are complexifications of maximal flats).

The fact that $M$ is an irreducible Hermitian symmetric space is equivalent to the fact that
the Lie algebra $\fk_\bR$ of $K_\bR$ has a 1-dimensional center $\fz_\bR$ and that the centralizer of
$\fz_\bR$ in the Lie algebra $\fgg_\bR$ of $G_\bR$ is $\fk_\bR$, see \eg \cite{Knapp}. If
$\fgg_\bR=\fk_\bR\oplus\fm_\bR$ is a Cartan decomposition of $\fgg_\bR$, the complex structure at the
origin is given by the adjoint action on $\fm_\bR$ of a suitably normalized element of $\fz_\bR$. It
follows that any maximal Abelian subspace $\ft_\bR$ of $\fk_\bR$ contains $\fz_\bR$ and is a Cartan
subalgebra of $\fgg_\bR$. The corresponding subgroup $T_\bR\subset K_\bR\subset G_\bR$ is a torus.

The Lie algebra $\fgg$ of $G$ is the complexification of the Lie algebra $\fgg_\bR$ of $G_\bR$. We
denote by $\fz\subset\ft\subset\fk \subset\fgg$ the complexifications of the subalgebras
$\fz_\bR\subset\ft_\bR\subset\fk_\bR$ and by $Z\subset T\subset K$ the corresponding complex algebraic
subgroups of $G$. We also let $\fm$ be the complexification of $\fm_\bR$, so that $\fgg=\fk\oplus\fm$
is a Cartan decomposition of $\fgg$. In particular, $[\fk,\fk]\subset\fk$, $[\fk,\fm]\subset\fm$ and $[\fm,\fm]
\subset\fk$.

Let $z\neq 0$ be an element in the center $\fz$ of $\fk$. Since the adjoint action of $\fz_\bR$ gives the
complex structure of $\fm_\bR$, $\ad(z)_{|\fm}$ has exactly two opposite eigenvalues. We
choose $z$ so that these eigenvalues are $2$ and $-2$
and we let
$$
\fm^+ = \{x\in\fm\mid \ad(z)x=2\,x\},\quad
\fm^- = \{x\in\fm\mid \ad(z)x=-2\,x\}
$$
be the corresponding eigenspaces. These are abelian subspaces of $\fm$.

This in turn implies that $\fp=\fk\oplus\fm^+$ is a maximal parabolic subalgebra of $\fgg$. Let~$P$ be
the corresponding parabolic subgroup of $G$ and $\check{M}$ the projective variety $G/P$. Then
$\check{M}$ is a symmetric space of the compact type called the \emph{compact dual} of $M$. Indeed if
we let $\check{\fgg}=\fk_\bR\oplus i\fm_\bR$, $\check{\fgg}$ is the Lie algebra of a real compact form
$\check{G}_\bR$ of $G$, and $\check{M}=\check{G}_\bR/K_\bR$. The symmetric space $M$ is then an open
$G_\bR$-orbit in $\check{M}=G/P$.

\begin{exem}
\label{exem:su}
For $p\leq q$ let $G_\bR=\SU(p,q)$ be the special unitary group of the Hermitian form of
signature $(p,q)$ on $\bC^{p+q}$ whose matrix in the canonical basis is
$I_{p,q}=\diag(-1,\ldots,-1,1,\ldots,1)$. Then one can choose
$K_\bR=\rS(\rU(p)\times \rU(q))$ and the symmetric space $M$ identifies with the set of
$p$-planes in $\bC^{p+q}$ on which the Hermitian form is negative definite. It has rank $p$. As a
bounded symmetric domain, it identifies with $\{Y\in\cM_{p,q}(\bC)\mid I_p-Y^\star Y\mbox{
is positive definite} \}$. We have
$G=\SL(p+q,\bC)$, $K=\rS(\GL(p)\times
\GL(q))$, $T$ is the torus of diagonal matrices in $G$, $Z$ is the 1-dimensional subgroup of
diagonal matrices of the form
$\diag(\lambda^q,\ldots,\lambda^q,\lambda^{-p},\ldots,\lambda^{-p})$
for $\lambda\in\bC^\star$,
$\fm\simeq \bC^{pq}\times\bC^{pq}$ is the subspace of block anti-diagonal matrices in
$\fsl(p+q,\bC)$. The compact real form of $G$ is $\SU(p+q)$ and the compact dual
$\check{M}$ identifies with the Grassmannian of $p$-planes in $\bC^{p+q}$, of which $M$ is an open
subset.
\end{exem}

Let $R$ be the set of roots of $G$. For $\alpha\in R$, $\fgg_\alpha\subset\fgg$ is the root space of $\alpha$,
$\fsl(\alpha)$ is the Lie subalgebra of $\fgg$ generated by $\fgg_\alpha$ and $\fgg_{-\alpha}$ and $\SL(\alpha)$ the
corresponding subgroup of $G$.

A root space $\fgg_\alpha$, $\alpha\in R$, is either a subspace of $\fk$ or a subspace of $\fm$, and the root~$\alpha$ is
said to be \emph{compact}
or \emph{noncompact} accordingly. Equivalently, a root $\alpha$ is compact if and only if
$\la\alpha,z\ra=0$. We have:
$$
\fk = \ft \oplus \bigoplus_{\alpha: \la\alpha,z\ra=0} \fgg_\alpha,\quad
\fm = \bigoplus_{\alpha: \la\alpha,z\ra\neq 0} \fgg_\alpha,\quad
\fm^+ = \bigoplus_{\alpha: \scal{\alpha,z}=2} \fgg_\alpha,\quad
\fm^- = \bigoplus_{\alpha: \scal{\alpha,z}=-2} \fgg_\alpha.
$$

We choose a basis $\Pi$ of $R$ so that the roots in $R(\fm^+)$ are positive roots. If $\alpha\in R$, we
write $\alpha=\sum_{\beta\in\Pi}n_\beta(\alpha)\beta$ the expression of $\alpha$ in terms of the simple roots. The
support $\supp(\alpha)$ of $\alpha\in R$ is the set $\{\beta\in\Pi\mid n_\beta(\alpha)\neq 0\}$.

The Weyl group of $R$ is denoted by $W$. If $\alpha\in R$, $\alpha^\vee\in\ft$ denotes the corresponding
coroot: it is defined by the relation $s_\alpha(\beta)=\beta-\scal{\beta,\alpha^\vee}\alpha$ for all $\beta\in \ft^*$, where
$s_\alpha\in W$ is the reflexion corresponding to $\alpha$ (in particular, $s_\alpha(\alpha)=-\alpha$ and hence
$\scal{\alpha,\alpha^\vee}=2$). If $\alpha=\sum_{\beta\in\Pi}n_\beta(\alpha)\beta$ is the
expression of a root $\alpha$ in the simple roots, then $\alpha^\vee= \sum_{\beta\in\Pi}n_\beta(\alpha)\,
\frac{\scal{\beta,\alpha^\vee}}{\scal{\alpha,\beta^\vee}}\,\beta^\vee$.
If $\beta\in\Pi$ is a simple root, $\varpi_\beta$ denotes the
fundamental weight corresponding to $\beta$.
The set of simple roots $\Pi$ is a basis of $\ft^\star$,
$\Pi^\vee=\{\beta^\vee\mid\beta\in\Pi\}$ is a basis of $\ft$ and $\{\varpi_\beta\mid \beta\in\Pi\}$ is by
definition the basis of $\ft^\star$ dual to $\Pi^\vee$.

In this paper, we use the convention that if the root system $R$ of $\fgg$ is simply laced (equivalently,
of type A, D or E), then all
the roots are long. Therefore short roots exist only if $R$ is not simply laced. There can be at most
two root lengths in~$R$. Moreover if there are two root lengths the ratio
$\frac{\text{long
root length}}{\text{short root length}}$ equals $\sqrt{2}$ because Hermitian symmetric spaces exist
only when the type of $G$ is $A$, $B$, $C$, $D$, $E_6$ or $E_7$
(in~particular, $G_2$ does not appear). We recall that for $\alpha,\beta\in R$, the ratio
$\sfrac{\scal{\alpha,\beta^\vee}}{\scal{\beta,\alpha^\vee}}$ equals
$\sfrac{(\text{root length of }\alpha)^2}{(\text{root length of }\beta)^2}$.

Linearly independent positive
roots $\delta_1,\ldots, \delta_s$ are said to be \emph{strongly orthogonal} if for
all $i\neq j$, $\delta_i\pm \delta_j$ is not a root. By \cite{Harish}, or
\citestar{Helgason}*{Ch.\,VIII, \textsection 7}, we have

\begin{fact}\label{fact:maxstrong}
All maximal sets of noncompact strongly orthogonal roots have cardinality $\rkm$.
\end{fact}

If $(\delta_1,\ldots, \delta_\rkm)$ is such a set then
$\fgg_{-\delta_1}\oplus\cdots\oplus\fgg_{-\delta_\rkm}$ is the tangent space to a maximal polydisc through the
origin in $M$.

\Subsection{The cominuscule representation of $G$ associated with $M$}
\label{subsec:minuscule-rep}

\Subsubsection{The representation}

\begin{defi}
A simple root $\beta\in \Pi$, is \emph{cominuscule} w.r.t.\ the root system $R$ if
$n_\beta(\alpha)\in\{-1,0,1\}$ for all $\alpha\in R$, or equivalently if $n_\beta(\Theta)=1$, where $\Theta$ is
the highest root of $R$.

An irreducible representation of $G$ whose highest weight is equal to $\varpi_\beta$ for $\beta$
a cominuscule root of $R$ is called a \emph{cominuscule representation}.
\end{defi}

The following is well-known, but we include a proof.

\begin{prop}
\label{prop:zeta}
There is a unique simple noncompact root. This root is long and cominuscule.
\end{prop}

\begin{nota}\label{nota:zeta}
This unique simple noncompact root will be denoted by $\zeta$.
\end{nota}

\begin{proof}
Since $R$ is
irreducible, the highest root $\Theta$ of $R$ is unique and $\Theta$ has a positive
coefficient on every simple root. Now, if there are more than one simple noncompact root, or if
$n_\zeta(\Theta)>1$ for some
noncompact simple root $\zeta$, then
$\la\Theta,z\ra>2$. Contradiction. Hence $\zeta$ is unique and cominuscule.
Assume that $\zeta$ is short (so that $R$ is not simply laced). Then, since $\Theta$ is long,
$$\scal{\Theta,\zeta^\vee}=\scal{\zeta,\Theta^\vee} \, \frac{(\text{root length of }\Theta)^2}{(\text{root length of }\zeta)^2}
\geq 2.$$
Since $s_\zeta(\Theta)$ is a positive root, this implies that $n_\zeta(\Theta) \geq 2$, contradicting the
fact that~$\zeta$ is cominuscule.
\end{proof}

Note that the maximal parabolic subalgebra $\fp=\fk\oplus\fm^+$ of $\fgg$ is the standard parabolic defined by the root
$\zeta$:
$$\fp=\ft\oplus\bigoplus_{\alpha:n_\zeta(\alpha)\geq 0}\fgg_\alpha.$$
Recall that $P$ is the corresponding parabolic subgroup of $G$ and that the compact
dual of $M$ is $\check{M}=G/P$.
By \cite{Murakami}, since $G$ is simply connected, the Picard group $\Pic(G/P)$ is
isomorphic to
the group of characters $\cX(P)$ of $P$. Since $\fp=\fz\oplus[\fp,\fp]$, $\cX(P)$ is isomorphic to
$\bZ$ and thus it is generated by the smallest positive character of $P$, namely $\varpi:=\varpi_\zeta$.
Here we embed $\cX(P)$ in $\ft^*$ via $\lambda \mto d\lambda$, the differential of $\lambda$ at the unit
element of $T$.
Moreover, the isomorphism $\iota:\cX(P) \simeq \Pic(G/P)$ is given by
$\lambda \mto (G \times \bC_\lambda)/P$,
where $\bC_\lambda$ denotes the $1$-dimensional $P$-module defined by $\lambda$.
In~particular, $\cL:=\iota(\varpi)$
is a generator of $\Pic(\check{M})$.

Let now $\bE$ be the cominuscule irreducible representation of $G$ whose highest weight is
$\varpi=\varpi_\zeta$. Let $\bE_\varpi$
be the $\varpi$-eigenspace of $\bE$. Then $\bE_\varpi$ is 1-dimensional and its
stabilizer in $G$ is $P$. This gives a $G$-equivariant holomorphic and isometric embedding of
$\check{M}=G/P$ in the projective space $\bP\bE$. It is called the \emph{first canonical embedding} of~$\check{M}$, and since $\bE_\varpi\simeq \bC_{\varpi}$ as $P$-modules,
we have $\cL \simeq \cO_{\bP
\bE}(-1)$.
See \eg \cite{nakagawa} for more details.

\begin{rema}
The complex group $G$ may have several cominuscule simple roots or representations. In
fact, cominuscule simple roots are in one to one correspondence with Hermitian real forms
$G_\bR$ (or $\check{G}_\bR$) of $G$.
When we will need a case by case argument, we will rather use the
classification of the cominuscule roots than the classification of the real groups $G_\bR$. This correspondence
is shown in the first two columns of Table~\ref{equa:tableau}.
\end{rema}

We now begin our study of the cominuscule representation $\bE$ of $G$.

\begin{nota}
We denote by $\mu_0$ the lowest weight of $\bE$ and by $X(\bE)$ the set
of weights of $\bE$. For $\chi \in X(\bE)$, we write
$\bE_\chi$ for the corresponding weight space. Recall that $\varpi$ is the highest weight of $\bE$.
\end{nota}

The fact that $\bE$ is cominuscule has the following consequence on the weights of $\bE$.

\begin{lemm}
\label{lemm:weight2}
For any weight $\chi$ of $\bE$, and any root $\alpha\in R$, $\scal{\chi,\alpha^\vee}
\in\{-2,-1,0,1,2\}$, and the equality
$\scal{\chi,\alpha^\vee} = \pm 2$ implies that $\alpha$ is short (hence that $R$ is not simply laced).
\end{lemm}

\begin{proof}
For the highest weight $\varpi$ of $\bE$, the results follows from the facts that
$\scal{\varpi,\alpha^\vee}=n_\zeta(\alpha)\,\sfrac{\scal{\zeta,\alpha^\vee}}{\scal{\alpha,\zeta^\vee}}$
and $\zeta$ is a long root. The
result still holds if
$\varpi$ is replaced by $w \cdot \varpi$, where $w \in W$ is arbitrary, and since any weight of
$\bE$ is in the convex hull of $W \cdot \varpi$, it holds for any weight.
\end{proof}

We deduce that the structure of $\bE$ with respect to the action of $\fgg_\alpha$ for $\alpha$ a long
root is particularly simple.

\begin{lemm}
\label{lemm:gV}
Let $\alpha$ be a long root and let $\chi$ be a weight of $\bE$.
We have
$$
\fgg_{-\alpha} \cdot \bE_\chi =
\begin{cases}
\bE_{\chi-\alpha}& \mbox{if } \scal{\chi,\alpha^\vee} = 1, \\
\{ 0 \}& \mbox{otherwise}.
\end{cases}
$$
\end{lemm}
\begin{proof}
Let $\alpha$ be long and let $\fsl(\alpha)$ be the Lie subalgebra of $\fgg$ isomorphic to $\fsl_2$
corresponding to $\alpha$.
Let $S \subset \bE$ be the $\fsl(\alpha)$-submodule generated by $\bE_\chi$. By
Lemma~\ref{lemm:weight2} and the representation theory of $\fsl_2$,
any irreducible component $V$ of $S$ is an $\fsl(\alpha)$-module of dimension $1$ or $2$.

We therefore have
only three possibilities. The first case is when $V=V_\chi$, $\fgg_\alpha \cdot V_\chi = \{0\}$ and
$\fgg_{-\alpha} \cdot V_\chi = \{0\}$. In this case, $\scal{\chi,\alpha^\vee}=0$.
The second case is when $V=V_\chi \oplus V_{\chi - \alpha}$, $\fgg_\alpha \cdot V_\chi = \{0\}$ and
$\fgg_{-\alpha} \cdot V_\chi = V_{\chi-\alpha}$. In this case, $\scal{\chi,\alpha^\vee}=1$ (and
$\scal{\chi-\alpha,\alpha^\vee}=-1$).
The third (symmetric) case is when $V=V_\chi \oplus V_{\chi + \alpha}$, $\fgg_\alpha \cdot V_\chi =
V_{\chi + \alpha}$ and
$\fgg_{-\alpha} \cdot V_\chi = \{0\}$. In this case, $\scal{\chi,\alpha^\vee}=-1$ (and
$\scal{\chi+\alpha,\alpha^\vee}=1$).

If $\scal{\chi,\alpha^\vee} = 1$, we deduce that $S=S_\chi \oplus S_{\chi-\alpha}$ and that
$\fgg_{-\alpha} \cdot S_\chi = S_{\chi-\alpha}$. We~have $s_\alpha(\chi)=\chi-\alpha$ so
$\dim(\bE_\chi) = \dim(\bE_{\chi-\alpha})$. The lemma is proved in this case.
If~$\scal{\chi,\alpha^\vee} \leq\nobreak 0$, we see that $\fgg_{-\alpha} \cdot \bE_\chi = \{0\}$.
\end{proof}

\Subsubsection{The grading}

\begin{nota}
Let $\zmax$ denote the number $\scal{\varpi,z}$.
\end{nota}

\begin{rema}
The precise value of $\zmax$ will not play any role in the following, however it has already
been computed in \citestar{KMgeneraltype}*{p.\,214-216}): with our choice of $z\in\fz$ we have
$\zmax=2\,\sfrac{\dim \check{M}}{c_1(\check{M})}$, where $c_1(\check{M})$ is the first Chern number
of $\check{M}$. For example if $G_\bR=\SU(p,q)$, $\zmax=\spfrac{2pq}{p+q}$.
\end{rema}

\begin{prop}
\label{prop:rang}
The set $\{\scal{\chi,z}\mid \chi \in X(\bE) \}$ is the set $$\{\zmax,\zmax-2,\ldots,\zmax-2\rkm\}.$$
\end{prop}
\begin{proof}
It follows from \citestar{richardson}*{Th.\,2.1} that the $W$-orbit of the
weight $\varpi$ is
exactly the set of weights of the form
$\varpi - \sum_{i=1}^k \delta_i$, where $(\delta_i)_{1 \leq i \leq k}$ is a family of noncompact
long strongly orthogonal roots.
In fact, the orbit $W \cdot \varpi$ is in bijection with $W/W_P$, so the equivalence between items
(c) and (e) in the cited
theorem proves the claim.
For any $i$, we have $\scal{\delta_i,z}=2 $, thus we have the equality of sets
$$\{\scal{\mu,z} \, | \, \mu \in W \cdot \varpi \} = \{ \zmax, \zmax-2, \ldots, \zmax-2\rkm \} .$$
In particular, $\scal{\mu_0,z}\!=\!\zmax-2 \rkm$ and for $\chi\!\in\!X(\bE)$, we have
\hbox{$\zmax-2\rkm\!\leq\!\scal{\chi,z} \leq \zmax$}. The result of the proposition now follows from the fact that
$2 $ is a divisor of $\scal{\alpha,z}$ for any root $\alpha$.
\end{proof}

Now we can introduce the grading of $\bE$.

\begin{defi}
\label{defi:Ei}
For a relative integer $i$, let $$\bE_i:= \bigoplus_{\chi: \scal{\chi,z} = \zmax-2 i} \bE_\chi.$$
\end{defi}

This grading corresponds to the decomposition of $\bE$ into irreducible $K$-modules.

\begin{prop}
\label{prop:Ei-irr}
The $K$-modules $\bE_i$ are irreducible.
\end{prop}

\begin{proof}
This might be well-known to experts, but we include a proof for completeness.
We give a case by case argument and use Table~\ref{equa:tableau}.

In type $A_{n-1}$, we have $\bE = \wedge^\rkm (\bC^\rkm \oplus \bC^{n-\rkm})$
and thus $\bE_i = \wedge^{\rkm-i} \bC^\rkm \otimes \wedge^i \bC^{n-\rkm}$: this is an irreducible $\rS(\GL_\rkm
\times \GL_{n-\rkm})$-module.

In type $B_n$, we have $\bE = \bC^{2n+1} = \bC \oplus \bC^{2n-1} \oplus \bC$,
and each summand is an irreducible $\Spin_{2n-1}$-module, hence an irreducible $K$-module.
In type $(D_n,\varpi_1)$ the situation is similar.

In type $(D_n,\varpi_n)$, $\bE$ is the spinor
representation of the spin group and, according to \cite{chevalley}, we have $\bE = \wedge^0 \bC^n
\oplus \wedge^2 \bC^n \oplus \cdots \oplus \wedge^{2p} \bC^n$ (where $p=\lfloor n/2\rfloor $). Thus $\bE_i =
\wedge^{2i} \bC^n$, and this is an irreducible $\GL_n$-module.

For the types $E_6$ and $E_7$, we use
models of these exceptional Lie algebras and their minuscule representations, as given for example
in \cite{manivel}. In type $E_6$, we have $\bE = \bC \oplus V_{16} \oplus V_{10}$, where $V_{16}$ is
a spinor representation and $V_{10}$ the vector representation of the spin group $\Spin_{10}$. In
type $E_7$, we have $\bE = \bC \oplus V_{27} \oplus V'_{27} \oplus \bC$, where $V_{27}$ and $V'_{27}$
are the two minuscule representations of a group of type $E_6$. In both cases, the $K$-modules
$\bE_i$ are irreducible.

We now deal with the case of type $C_n$. The following representation-theoretic
argument has been suggested to us by an anonymous referee that we would like to thank.
We denote by $\bC^{2n} = \bC^n_a \oplus \bC^n_b$ a symplectic $2n$-dimensional space, with $\bC^n_a$ and
$\bC^n_b$ supplementary isotropic subspaces. We then have
$\bE=\left (\wedge^n \bC^{2n} \right)_\omega$,
where the symbol $\omega$ means that we take in $\wedge^n \bC^{2n}$ the irreducible
$\Sp_{2n}$-submodule containing the highest weight line $\wedge^n \bC^n_a$.
More precisely, by \citestar{fulton-harris}*{Th.\,17.5}, we have
a decomposition of $\wedge^n \bC^{2n}$ as an $\Sp_{2n}$-module as follows:
$\wedge^n \bC^{2n} = \bE \oplus \omega \wedge (\wedge^{n-2} \bC^{2n})$.

Thus,
$\wedge^{n-i} \bC^n_a \otimes \wedge^i \bC^n_b =
\bE_i \oplus \omega \wedge (\wedge^{n-i-1} \bC^n_a) \otimes \wedge^{i-1} \bC^n_b.$
We now consider this as a representation of $K=\GL_n$. As $K$-modules, we have $\bC^n_b \simeq {(\bC^n_a)}^\star$.
We thus get
$\wedge^{n-i} \bC^n_a \otimes \wedge^{n-i} \bC^n_a \simeq
\bE_i \oplus \wedge^{n-i-1} \bC^n_a \otimes \wedge^{n-i+1} \bC^n_a.$
From this last equation, it follows that $\bE_i$ is the Cartan square of $\wedge^{n-i} \bC^n_a$ and therefore
it is an irreducible $\GL_n$-module.
\end{proof}

\begin{prop}
\label{prop:Ei}
We have the following properties:
\begin{enumeratea}
\item\label{prop:Eia} $\bE = \bE_0 \oplus \bE_{1} \oplus \cdots \oplus \bE_{\rkm}$.
\item\label{prop:Eib} $\bE_0 = \bE_{\varpi}$.
\item\label{prop:Eic} $\bE_{i+1} = \fm^- \cdot \bE_i$.
\item\label{prop:Eid} The map $\bE_0 \otimes \fm^- \to \bE_1$ is an isomorphism.
\end{enumeratea}
\end{prop}
\begin{proof}
Only the last two points need a proof. Let $U(\fm^-)$ denote the enveloping algebra of $\fm^-$. The
third point
follows from the fact that $\bE = U(\fm^-) \cdot \bE_{\varpi}$ and the fact that for $\alpha$ a root of $\fm^-$,
we have $\scal{\alpha,z} = -2 $.
The last point follows by Schur's lemma since $\fm^-$ and $\bE_1$ are irreducible $\fk$-modules and
$\bE_0$ is 1-dimensional.
\end{proof}

\begin{rema}
Those statements are well-known. In the case where $\bE$ is of tube type, they are proved in \citestar{Gross}*{Prop.\,5.2}.
\end{rema}

\subsection{Dominant orthogonal sequences}
\label{subsec:DOS}

In \cite{kostant}, Kostant introduced his so-called ``chain cascade'' of orthogonal roots.
Here we will need a version of his algorithm where we impose that all the roots of the chain cascade
have a positive coefficient on~$\zeta$. Note that a similar algorithm is used in \cite{buch}.

We define an integer $q$ and, for any integer $i$
such that $1 \leq i \leq q$,
a root $\alpha_i$ together with the subset $\Pi_i \subset \Pi$, by the following inductive process:
\begin{itemize}
\item We let $\Pi_1=\Pi$.
\item Assuming that $\alpha_1,\ldots,\alpha_{i-1}$ and $\Pi_1,\ldots,\Pi_i$ have been defined,
we let $\alpha_i$ be the highest root of the root system $R(\Pi_i)$ generated by $\Pi_i$.
\item We let $\Sigma_i \subset \Pi_i$ be the set of simple roots $\beta$ such that
$\scal{\alpha_i^\vee,\beta} \neq 0$.
\item If $\zeta \in \Sigma_i$, then $q=i$ and the algorithm terminates. Otherwise, $\Pi_{i+1}$
is the connected component of $\Pi_i \setminus \Sigma_i$ containing $\zeta$.
\end{itemize}

\begin{defi}
\label{def:cascade}
If $(\alpha_i)_{1 \leq i \leq q}$ is the sequence defined by this process,
we say that it is the \emph{maximal dominant orthogonal sequence} for $\varpi$. More generally,
the sequences $(\alpha_i)_{1 \leq i \leq r}$ for $1\leq
r\leq q$ are called the \emph{dominant orthogonal sequences}.
\end{defi}

\begin{exem}
\label{exem:cascade}
In type $A_{n-1}$ with the standard base $(\beta_1,\ldots,\beta_{n-1})$,
and for
$\zeta=\beta_p$ with $2p\leq n$, we get $\Pi_i=\{\beta_i,\ldots,\beta_{n-i}\}$ and the maximal
dominant orthogonal sequence is
$(\, \sum_{i=1}^{n-1}\beta_i, \, \sum_{i=2}^{n-2}\beta_i,\, \ldots, \, \sum_{i=p}^{n-p}\beta_i
\,),$ so $q=p$.
\end{exem}

For later use, we record in Table~\ref{equa:tableau} below what are
the dominant orthogonal sequences $(\alpha_1,\ldots,\alpha_r)$ in all cases. The number $r$ varies from
$1$ to $r_M$, see Proposition~\ref{prop:cascade}\,\eqref{prop:cascade2}.

\afterpage{%
\clearpage
\centering
\smaller\smaller
\begin{table}[ht]
\centering
$$
\begin{array}{|cc|rcc|}
\hline
\multirow{2}{*}{$(G,\varpi)$}
& \multirow{2}{*}{$\fgg_\bR$, $\rkm$}
& \multirow{2}{*}{}
&\multirow{2}{*}{$\alpha_1^\vee+\cdots+\alpha_r^\vee$}
& \multirow{2}{*}{$\alpha_r$}
\\
&&&&\\
\hline
\multirow{2}{*}{ $(A_{n-1},\varpi_{p})$, $2p\leq n$}
& \multirow{2}{*}{$\fsu(p,n-p)$}
& \multirow{2}{*}{$\mbox{if } 2r < n:$}
& \multirow{2}{*}{\begin{tikzpicture}
\dynkin{A}{o.ooo.ooo.o}
\dynkinLabelRoot{1}{0}
\dynkinLabelRoot{2}{0}
\dynkinLabelRoot{3}{1}
\dynkinLabelRoot*{3}{r}
\dynkinLabelRoot{4}{0}
\dynkinLabelRoot{5}{0}
\dynkinLabelRoot{6}{1}
\dynkinLabelRoot{7}{0}
\dynkinLabelRoot{8}{0}
\dynkinLabelRoot*{6}{n-r}
\end{tikzpicture}}
&\multirow{2}{*}{\begin{tikzpicture}
\dynkin{A}{o.oo.oo.o}
\dynkinLabelRoot{1}{0}
\dynkinLabelRoot{2}{0}
\dynkinLabelRoot{3}{1}
\dynkinLabelRoot*{3}{r}
\dynkinLabelRoot{4}{1}
\dynkinLabelRoot{5}{0}
\dynkinLabelRoot{6}{0}
\dynkinLabelRoot*{4}{n-r}
\end{tikzpicture}}
\\
\multirow{2}{*}{\vspace{-8pt}\begin{tikzpicture}
\dynkin{A}{o.o*o.o}
\dynkinLabelRoot*{3}{p}
\dynkinLabelRoot*{1}{1}
\dynkinLabelRoot*{5}{n-1}
\end{tikzpicture}}
& \multirow{2}{*}{\vspace{-8pt}$r_M=p$}&&&\\
\cline{3-5}
&
&\multirow{2}{*}{$\mbox{if } 2r=n:$}
& \multirow{2}{*}{\begin{tikzpicture}
\dynkin{A}{o.ooo.o}
\dynkinLabelRoot{1}{0}
\dynkinLabelRoot{2}{0}
\dynkinLabelRoot{3}{2}
\dynkinLabelRoot*{3}{p}
\dynkinLabelRoot{4}{0}
\dynkinLabelRoot{5}{0}
\end{tikzpicture}}
&\multirow{2}{*}{\begin{tikzpicture}
\dynkin{A}{o.ooo.o}
\dynkinLabelRoot{1}{0}
\dynkinLabelRoot{2}{0}
\dynkinLabelRoot{3}{1}
\dynkinLabelRoot*{3}{p}
\dynkinLabelRoot{4}{0}
\dynkinLabelRoot{5}{0}
\end{tikzpicture}}
\\
&&&&\\
\hline
\multirow{2}{*}{$(B_n,\varpi_1)$}
& \multirow{2}{*}{$\fso(2,2n-1)$}
& \multirow{2}{*}{$\mbox{if } r=1:$}
& \multirow{2}{*}{
\vspace{-6pt}\begin{tikzpicture}
\dynkin{B}{ooo.ooo}
\dynkinLabelRoot{1}{0}
\dynkinLabelRoot{2}{1}
\dynkinLabelRoot{3}{0}
\dynkinLabelRoot{4}{0}
\dynkinLabelRoot{5}{0}
\dynkinLabelRoot{6}{0}
\end{tikzpicture}
}
& \multirow{2}{*}{
\vspace{-6pt}\begin{tikzpicture}
\dynkin{B}{oo.ooo}
\dynkinLabelRoot{1}{1}
\dynkinLabelRoot{2}{2}
\dynkinLabelRoot{3}{2}
\dynkinLabelRoot{4}{2}
\dynkinLabelRoot{5}{2}
\end{tikzpicture}
}
\\
\multirow{2}{*}{\vspace{-8pt}\begin{tikzpicture}
\dynkin{B}{*o.ooo}
\dynkinLabelRoot*{1}{1}
\dynkinLabelRoot*{5}{n}
\end{tikzpicture}}
& \multirow{2}{*}{\vspace{-8pt}$\rkm=2$}&&&
\\
\cline{3-5}
&
& \multirow{2}{*}{$\mbox{if } r=2:$}
& \multirow{2}{*}{
\vspace{-6pt}\begin{tikzpicture}
\dynkin{B}{oo.ooo}
\dynkinLabelRoot{1}{2}
\dynkinLabelRoot{2}{0}
\dynkinLabelRoot{3}{0}
\dynkinLabelRoot{4}{0}
\dynkinLabelRoot{5}{0}
\end{tikzpicture}
}
& \multirow{2}{*}{
\vspace{-6pt}\begin{tikzpicture}
\dynkin{B}{oo.ooo}
\dynkinLabelRoot{1}{1}
\dynkinLabelRoot{2}{0}
\dynkinLabelRoot{3}{0}
\dynkinLabelRoot{4}{0}
\dynkinLabelRoot{5}{0}
\end{tikzpicture}}
\\
&&&&
\\
\hline
\multirow{2}{*}{$(C_n,\varpi_n)$}
&\multirow{2}{*}{$\fsp(2n,\bR)$}
&\multirow{2}{*}{$\mbox{if } r < n:$}
& \multirow{2}{*}{
\begin{tikzpicture}
\dynkin{C}{o.ooo.ooo}
\dynkinLabelRoot*{3}{r}
\dynkinLabelRoot{1}{0}
\dynkinLabelRoot{2}{0}
\dynkinLabelRoot{3}{2}
\dynkinLabelRoot{4}{0}
\dynkinLabelRoot{5}{0}
\dynkinLabelRoot{6}{0}
\dynkinLabelRoot{7}{0}
\end{tikzpicture}
}
& \multirow{2}{*}{
\begin{tikzpicture}
\dynkin{C}{o.oo.ooo}
\dynkinLabelRoot*{3}{r}
\dynkinLabelRoot{1}{0}
\dynkinLabelRoot{2}{0}
\dynkinLabelRoot{3}{2}
\dynkinLabelRoot{4}{2}
\dynkinLabelRoot{5}{2}
\dynkinLabelRoot{6}{1}
\end{tikzpicture}}
\\
\multirow{2}{*}{\vspace{-8pt}\begin{tikzpicture}
\dynkin{C}{o.oo*}
\dynkinLabelRoot*{1}{1}
\dynkinLabelRoot*{4}{n}
\end{tikzpicture}}
& \multirow{2}{*}{\vspace{-8pt}$\rkm=n$}&&&
\\
\cline{3-5}
&
& \multirow{2}{*}{$\mbox{if } r = n:$}
& \multirow{2}{*}{
\vspace{-6pt}\begin{tikzpicture}
\dynkin{C}{o.ooo}
\dynkinLabelRoot{1}{0}
\dynkinLabelRoot{2}{0}
\dynkinLabelRoot{3}{0}
\dynkinLabelRoot{4}{2}
\end{tikzpicture}
}
& \multirow{2}{*}{
\vspace{-6pt}\begin{tikzpicture}
\dynkin{C}{o.ooo}
\dynkinLabelRoot{1}{0}
\dynkinLabelRoot{2}{0}
\dynkinLabelRoot{3}{0}
\dynkinLabelRoot{4}{1}
\end{tikzpicture}
}
\\
&&&& \\
\hline
\multirow{3}{*}{ $(D_n,\varpi_1)$}
& \multirow{3}{*}{$\fso(2,2n-2)$}
& \multirow{3}{*}{$\mbox{if } r = 1:$}
& \multirow{3}{*}{\begin{tikzpicture}
\dynkin{D}{ooo.oooo}
\dynkinLabelRoot{1}{0}
\dynkinLabelRoot{2}{1}
\dynkinLabelRoot{3}{0}
\dynkinLabelRoot{4}{0}
\dynkinLabelRoot{5}{0}
\dynkinLabelRoot{6}{0}
\dynkinLabelRoot{7}{0}
\end{tikzpicture}}
& \multirow{3}{*}{ \begin{tikzpicture}
\dynkin{D}{oo.oooo}
\dynkinLabelRoot{1}{1}
\dynkinLabelRoot{2}{2}
\dynkinLabelRoot{3}{2}
\dynkinLabelRoot{4}{2}
\dynkinLabelRoot{5}{1}
\dynkinLabelRoot{6}{1}
\end{tikzpicture}}
\\
\multirow{3}{*}{\vspace{-16pt}\begin{tikzpicture}
\dynkin{D}{*o.ooo}
\dynkinLabelRoot*{1}{1}
\dynkinLabelRoot{4}{n-1}
\dynkinLabelRoot{5}{n}
\end{tikzpicture}}
& \multirow{3}{*}{\vspace{-16pt} $\rkm = 2$}&&&\\
&&&&\\
\cline{3-5}
&
& \multirow{3}{*}{ $\mbox{if } r = 2:$}
& \multirow{3}{*}{\begin{tikzpicture}
\dynkin{D}{oo.oooo}
\dynkinLabelRoot{1}{2}
\dynkinLabelRoot{2}{0}
\dynkinLabelRoot{3}{0}
\dynkinLabelRoot{4}{0}
\dynkinLabelRoot{5}{0}
\dynkinLabelRoot{6}{0}
\end{tikzpicture} }
& \multirow{3}{*}{ \begin{tikzpicture}
\dynkin{D}{oo.oooo}
\dynkinLabelRoot{1}{1}
\dynkinLabelRoot{2}{0}
\dynkinLabelRoot{3}{0}
\dynkinLabelRoot{4}{0}
\dynkinLabelRoot{5}{0}
\dynkinLabelRoot{6}{0}
\end{tikzpicture}}
\\
&&&&\\
&&&&\\
\hline
\multirow{3}{*}{ $(D_n,\varpi_n)$}
& \multirow{3}{*}{$\fso^\star(2n)$}
& \multirow{3}{*}{$\mbox{if } 2r \leq n-2:$}
& \multirow{3}{*}{
\begin{tikzpicture}
\dynkin{D}{o.ooo.oooo}
\dynkinLabelRoot{1}{0}
\dynkinLabelRoot{2}{0}
\dynkinLabelRoot{3}{1}
\dynkinLabelRoot*{3}{2r}
\dynkinLabelRoot{4}{0}
\dynkinLabelRoot{5}{0}
\dynkinLabelRoot{6}{0}
\dynkinLabelRoot{7}{0}
\dynkinLabelRoot{8}{0}
\end{tikzpicture}
}
& \multirow{3}{*}{
\begin{tikzpicture}
\dynkin{D}{o.ooo.oooo}
\dynkinLabelRoot{1}{0}
\dynkinLabelRoot{2}{0}
\dynkinLabelRoot{3}{1}
\dynkinLabelRoot*{4}{2r}
\dynkinLabelRoot{4}{2}
\dynkinLabelRoot{5}{2}
\dynkinLabelRoot{6}{2}
\dynkinLabelRoot{7}{1}
\dynkinLabelRoot{8}{1}
\end{tikzpicture}
}
\\
\multirow{3}{*}{\vspace{-16pt}\begin{tikzpicture}
\dynkin{D}{oo.oo*}
\dynkinLabelRoot*{1}{1}
\dynkinLabelRoot{4}{n-1}
\dynkinLabelRoot{5}{n}
\end{tikzpicture}}
& \multirow{3}{*}{\vspace{-16pt}$\rkm=\lfloor\sfrac n2\rfloor$}&&&\\
&&&&\\
\cline{3-5}
&
& \multirow{3}{*}{$\mbox{if } 2r = n-1:$}
& \multirow{3}{*}{
\begin{tikzpicture}
\dynkin{D}{o.oooo}
\dynkinLabelRoot{1}{0}
\dynkinLabelRoot{2}{0}
\dynkinLabelRoot{3}{0}
\dynkinLabelRoot{4}{1}
\dynkinLabelRoot{5}{1}
\end{tikzpicture}
}
& \multirow{3}{*}{
\begin{tikzpicture}
\dynkin{D}{o.oooo}
\dynkinLabelRoot{1}{0}
\dynkinLabelRoot{2}{0}
\dynkinLabelRoot{3}{1}
\dynkinLabelRoot{4}{1}
\dynkinLabelRoot{5}{1}
\end{tikzpicture}
}
\\
&&&&\\
&&&&\\
\cline{3-5}
&
& \multirow{3}{*}{$\mbox{if } 2r = n:$}
& \multirow{3}{*}{
\begin{tikzpicture}
\dynkin{D}{o.oooo}
\dynkinLabelRoot{1}{0}
\dynkinLabelRoot{2}{0}
\dynkinLabelRoot{3}{0}
\dynkinLabelRoot{4}{0}
\dynkinLabelRoot{5}{2}
\end{tikzpicture}}
& \multirow{3}{*}{
\begin{tikzpicture}
\dynkin{D}{o.oooo}
\dynkinLabelRoot{1}{0}
\dynkinLabelRoot{2}{0}
\dynkinLabelRoot{3}{0}
\dynkinLabelRoot{4}{0}
\dynkinLabelRoot{5}{1}
\end{tikzpicture}}
\\
&&&&\\
&&&&\\
\hline
\multirow{3}{*}{$ (E_6,\varpi_1)$}
& \multirow{3}{*}{$\fe_{6(-14)}$}
& \multirow{3}{*}{$\mbox{if } r = 1:$}
& \multirow{3}{*}{
\begin{tikzpicture}
\dynkin{E}{oooooo}
\dynkinLabelRoot{1}{0}
\dynkinLabelRoot{2}{1}
\dynkinLabelRoot{3}{0}
\dynkinLabelRoot{4}{0}
\dynkinLabelRoot{5}{0}
\dynkinLabelRoot{6}{0}
\end{tikzpicture}}
& \multirow{3}{*}{
\begin{tikzpicture}
\dynkin{E}{oooooo}
\dynkinLabelRoot{1}{1}
\dynkinLabelRoot{2}{2}
\dynkinLabelRoot{3}{2}
\dynkinLabelRoot{4}{3}
\dynkinLabelRoot{5}{2}
\dynkinLabelRoot{6}{1}
\end{tikzpicture}}
\\
\multirow{3}{*}{\vspace{-16pt}
\begin{tikzpicture}
\dynkin{E}{*ooooo}
\dynkinLabelRoot*{1}{1}
\end{tikzpicture}
}
&\multirow{3}{*}{\vspace{-16pt}$\rkm=2$}&&&\\
&&&&\\
\cline{3-5}
&
& \multirow{3}{*}{$\mbox{if } r = 2:$}
& \multirow{3}{*}{
\begin{tikzpicture}
\dynkin{E}{oooooo}
\dynkinLabelRoot{1}{1}
\dynkinLabelRoot{2}{0}
\dynkinLabelRoot{3}{0}
\dynkinLabelRoot{4}{0}
\dynkinLabelRoot{5}{0}
\dynkinLabelRoot{6}{1}
\end{tikzpicture}}
& \multirow{3}{*}{
\begin{tikzpicture}
\dynkin{E}{oooooo}
\dynkinLabelRoot{1}{1}
\dynkinLabelRoot{2}{0}
\dynkinLabelRoot{3}{1}
\dynkinLabelRoot{4}{1}
\dynkinLabelRoot{5}{1}
\dynkinLabelRoot{6}{1}
\end{tikzpicture}}
\\
&&&&\\
&&&&\\
\hline
\multirow{3}{*}{$(E_7,\varpi_7)$}
& \multirow{3}{*}{$\fe_{7(-25)}$}
& \multirow{3}{*}{$\mbox{if } r = 1:$}
& \multirow{3}{*}{
\begin{tikzpicture}
\dynkin{E}{ooooooo}
\dynkinLabelRoot{1}{1}
\dynkinLabelRoot{2}{0}
\dynkinLabelRoot{3}{0}
\dynkinLabelRoot{4}{0}
\dynkinLabelRoot{5}{0}
\dynkinLabelRoot{6}{0}
\dynkinLabelRoot{7}{0}
\end{tikzpicture}
}
& \multirow{3}{*}{
\begin{tikzpicture}
\dynkin{E}{ooooooo}
\dynkinLabelRoot{1}{2}
\dynkinLabelRoot{2}{2}
\dynkinLabelRoot{3}{3}
\dynkinLabelRoot{4}{4}
\dynkinLabelRoot{5}{3}
\dynkinLabelRoot{6}{2}
\dynkinLabelRoot{7}{1}
\end{tikzpicture}
}
\\
\multirow{3}{*}{\vspace{-16pt}\begin{tikzpicture}
\dynkin{E}{oooooo*}
\dynkinLabelRoot*{7}{7}
\end{tikzpicture}}
&\multirow{3}{*}{\vspace{-16pt}$\rkm=3$}&&&\\
&&&&\\
\cline{3-5}
&
& \multirow{3}{*}{$\mbox{if } r = 2:$}
& \multirow{3}{*}{
\begin{tikzpicture}
\dynkin{E}{ooooooo}
\dynkinLabelRoot{1}{0}
\dynkinLabelRoot{2}{0}
\dynkinLabelRoot{3}{0}
\dynkinLabelRoot{4}{0}
\dynkinLabelRoot{5}{0}
\dynkinLabelRoot{6}{1}
\dynkinLabelRoot{7}{0}
\end{tikzpicture}
}
& \multirow{3}{*}{
\begin{tikzpicture}
\dynkin{E}{ooooooo}
\dynkinLabelRoot{1}{0}
\dynkinLabelRoot{2}{1}
\dynkinLabelRoot{3}{1}
\dynkinLabelRoot{4}{2}
\dynkinLabelRoot{5}{2}
\dynkinLabelRoot{6}{2}
\dynkinLabelRoot{7}{1}
\end{tikzpicture}
}
\\
&&&&\\
&&&&\\
\cline{3-5}
&& \multirow{3}{*}{$\mbox{if } r = 3:$}
& \multirow{3}{*}{
\begin{tikzpicture}
\dynkin{E}{ooooooo}
\dynkinLabelRoot{1}{0}
\dynkinLabelRoot{2}{0}
\dynkinLabelRoot{3}{0}
\dynkinLabelRoot{4}{0}
\dynkinLabelRoot{5}{0}
\dynkinLabelRoot{6}{0}
\dynkinLabelRoot{7}{2}
\end{tikzpicture}
}
& \multirow{3}{*}{
\begin{tikzpicture}
\dynkin{E}{ooooooo}
\dynkinLabelRoot{1}{0}
\dynkinLabelRoot{2}{0}
\dynkinLabelRoot{3}{0}
\dynkinLabelRoot{4}{0}
\dynkinLabelRoot{5}{0}
\dynkinLabelRoot{6}{0}
\dynkinLabelRoot{7}{1}
\end{tikzpicture}
}\\
&&&&\\
&&&&\\
\hline
\end{array}
$$
\caption{\smaller Dominant orthogonal sequences $(\alpha_1,\ldots,\alpha_r)$ for $1\leq r\leq r_M$
\smaller (We~do not indicate
all the roots $\alpha_i$, but rather the sum of the corresponding coroots
$\alpha_1^\vee + \cdots + \alpha_r^\vee$, by indicating
on a Dynkin diagram the values
\hbox{$\scal{\alpha_1^\vee + \cdots + \alpha_r^\vee,\beta}$} for all simple roots $\beta$.
In the last column, we express the root $\alpha_r$ which is also the smallest root $\alpha$ such that
$\scal{\alpha_1^\vee + \cdots + \alpha_r^\vee,\alpha}=2$.)}
\label{equa:tableau}
\end{table}
\clearpage
}

The following proposition essentially adapts the results of \cite{kostant} to our context and
explains our terminology.

\begin{prop}
\label{prop:cascade}
Let $(\alpha_i)_{1 \leq i \leq q}$ be the maximal dominant orthogonal sequence for~$\varpi$. Then
\begin{enumerate}
\item\label{prop:cascade1}
The roots $\alpha_i$ are long and strongly orthogonal.
\item\label{prop:cascade2}
$q=\rkm$.
\item\label{prop:cascade3}
For any integer $i \leq \rkm$, $\alpha_1^\vee + \cdots + \alpha_i^\vee$ is a dominant coweight:
for all $\beta$ in $\Pi$, we have $\la \beta, \alpha_1^\vee + \cdots + \alpha_i^\vee\ra\geq 0$.
\item\label{prop:cascade4}
$\varpi - \alpha_1 - \cdots - \alpha_\rkm$ is the lowest weight $\mu_0$ of the irreducible
$G$-module $\bE$.
\end{enumerate}
\end{prop}
\begin{proof}
Kostant \citestar{kostant}*{Lem.\,1.6} proved the strong orthogonality. For the convenience of the reader, we
recall his arguments in our context.
The root $\alpha_i$ is the highest root of $R(\Pi_i)$, and $\Pi_i$ contains the long root $\zeta$,
so $\alpha_i$
is long. By construction, $\scal{\alpha_i^\vee,\alpha_j}=0$ if $j>i$, so $\alpha_i$ and $\alpha_j$ are
orthogonal. Since they are both long, they are strongly orthogonal, otherwise $\alpha_i\pm\alpha_j$ would
be longer than $\alpha_i$. This proves \eqref{prop:cascade1}.

For \eqref{prop:cascade2}, by Fact~\ref{fact:maxstrong}, it is enough to prove that $(\alpha_1,\ldots,\alpha_q)$ is a
maximal sequence of orthogonal roots (see also \citestar{kostant}*{Th.\,1.8}). Let
$\alpha \in R$ be such that $\scal{\alpha_i^\vee,\alpha}=0$ holds for all $i$. We can and will assume that
$\alpha >0$. Let $i$ be the greatest integer such that $\alpha \in R(\Pi_i)$. By maximality of $i$
there exists
a simple root $\beta$ in $\supp(\alpha) \cap \Sigma_i$. Since $\alpha_i$ is dominant on
$R(\Pi_i)$, we have
$\scal{\alpha_i^\vee,\alpha} \geq \scal{\alpha_i^\vee,\beta} > 0$, a contradiction to the assumption made on
$\alpha$.

For the third point, let $i \leq \rkm$ and let $\beta \in \Pi$.
\begin{itemize}
\item If $\beta \in \Pi_i$, by construction, $\scal{\alpha_j^\vee,\beta}=0$ for $j<i$. Since
$\alpha_i$ is the
highest root of~$\Pi_i$, $\scal{\alpha_i^\vee,\beta} \geq 0$. Thus
$\scal{\suma{i},\beta} = \scal{\av_i,\beta} \geq 0$.

\item If $\beta \in \Sigma_{i-1}$, then $\scal{\av_{i-1},\beta} \geq 1$ and $\scal{\av_j,\beta}=0$
for $j<i-1$.
Since $\alpha_i$ is long, $\scal{\av_i,\beta} \geq -1$. Thus, $\scal{\suma{i},\beta} \geq 0$.
\item If $\beta \in \Sigma_j$ for $j<i-1$, then, by construction of $\Pi_i$,
$\scal{\av_i,\beta}=0$. By induction on $i$,
$\scal{\suma{i-1},\beta} \geq 0$, so $\scal{\suma{i},\beta} \geq 0$.
\end{itemize}

For \eqref{prop:cascade4}, by Lemma~\ref{lemm:weight2} and since
$\alpha_i$ is long, we have $\scal{\varpi,\alpha_i^\vee}=1$. Thus $s_{\alpha_i}(\varpi)=\varpi -
\scal{\varpi,\alpha_i^\vee}\alpha_i = \varpi-\alpha_i$.
Therefore
$s_{\alpha_1} \cdots s_{\alpha_\rkm}(\varpi)=\varpi-\alpha_1 - \cdots - \alpha_\rkm$ is a weight of $\bE$.
We prove that it is a lowest weight. Let $\beta \in \Pi$. If $\beta \neq \zeta$, then
$$\scal{\varpi-\alpha_1 - \cdots - \alpha_\rkm,\beta^\vee} = \scal{-\alpha_1 - \cdots -
\alpha_\rkm,\beta^\vee},$$ and
this is non-positive by (3) and because all the roots $\alpha_i$ have the same length. If $\beta=\zeta$, then
we compute that $\scal{\varpi-\alpha_1 - \cdots - \alpha_\rkm,\zeta^\vee} = 1 -
\scal{\alpha_\rkm,\zeta^\vee} \leq 0$
since $\zeta \in \Sigma_\rkm$.
\end{proof}

We make the following observations.

\begin{prop}\label{prop:h}
Let $1\leq r\leq \rkm$, $(\alpha_1,\ldots,\alpha_r)$ be the corresponding dominant orthogonal
sequence, and $h_r=\alpha_1^\vee + \cdots + \alpha_r^\vee$. Then:
\begin{enumerate}
\item \label{item:dominant} $h_r$ is dominant: for any positive root $\alpha$, we have $\scal{\alpha,h_r} \geq 0$.
\item \label{item:leq2} For any root $\alpha$, we have $\scal{\alpha,h_r} \leq 2$.
\item \label{item:smallest} $\alpha_r$ is the smallest root $\alpha$ such that $\scal{\alpha,h_r}=2$.
\item \label{item:scal} If a root $\alpha$ satisfies $\scal{\alpha,h_r} = 2$, then $\scal{\alpha,z}=2$.
\item \label{item:varpi} We have $\scal{\varpi,h_r}=r$.
\end{enumerate}
\end{prop}

\begin{proof}
Recall Table~\ref{equa:tableau}. The first point has been proved in Proposition~\ref{prop:cascade}\,\eqref{prop:cascade3}.
Let $\Theta$ be the highest root of the root system of $G$,
which can be found for example in~\cite{bourbaki}.
Since $h_r$ is dominant, the second item follows from the
fact that $\scal{\Theta,h_r}=2$ in all cases.

For the third item, we have indicated the root $\alpha_r$ in the last column of the table. It
readily follows that it is the smallest root $\alpha$ such that $\scal{\alpha,h_r}=2$.
Since it has coefficient $1$ on $\zeta$, we have $\scal{\alpha_r,z}=2$. Any root $\alpha$ as in
the fourth item will be bigger than $\alpha_r$, so the claim follows.

To prove that $\scal{\varpi,h_r}=r$, recall that if $(\alpha_1,\ldots,\alpha_\rkm)$ is the maximal dominant
orthogonal sequence, the weight
$\varpi- \sum_1^\rkm \alpha_i$ is the lowest weight $\mu_0$ by Proposition~\ref{prop:cascade}\,\eqref{prop:cascade4}. Thus,
$\scal{\varpi-\mu_0,h_r} = \scal{\alpha_1+\cdots+\alpha_\rkm,\alpha_1^\vee+\cdots+\alpha_r^\vee}=2r$,
since $\scal{\alpha_i,\alpha_j^\vee}=2\delta_{i,j}$ by orthogonality of the roots $\alpha_i$.
Moreover, $h_r$ is part of some $\fsl_2$\nobreakdash-triple $(x, h_r,y)$, so
$\scal{\mu_0, h_r}=-\scal{\varpi, h_r}$, by the representation theory of $\fsl_2$.
This proves that $\scal{\varpi, h_r}=r$.
\end{proof}

\Subsection{The self-dual Higgs submodule associated with a nilpotent element}
\label{subsec:higgs-submodule}

We explain in this section that an element $y\in\fm^-$ defines a special graded subspace in
$\bE$. This algebraic construction will be used in Section~\ref{subsec:HiggsSubsheaf} to produce
the leafwise Higgs subsheaves needed to prove the Milnor-Wood inequality.

\subsubsection{$K$-orbits in $\fm^-$: ranks of nilpotent elements}\label{subsubsec:rank}

We begin by describing the $K$\nobreakdash-orbits in $\fm^-$, introducing the notion of \emph{rank}, and then we
choose nice representatives in the orbits to simplify our work. Let us recall the following result, which is proved \eg in
\cite{Harish}, \hbox{\citestar{Helgason}*{Ch.\,VIII,\,\textsection 7}}, or \cite{WolfFineStructure}.

\begin{prop}
\label{prop:rank}
The orbits of the complex group $K$ in $\fm^-$ are parametrized by integers $r \in \{0,\ldots,\rkm\}$.
A representative of each orbit is $y_{\alpha_1} + \cdots + y_{\alpha_r}$, where $(\alpha_1,\ldots,\alpha_{r_M})$ is
the maximal dominant orthogonal sequence.
\end{prop}

\begin{defi}\label{defi:rank}
Let $y\in\fm^-$. The \emph{rank} $\rky$ of $y$ is the integer $r$ such that $y$ is in the $K$-orbit
of $y_{\alpha_1} + \cdots + y_{\alpha_r}$.
\end{defi}

\begin{exem}
\label{exem:matrices}
In case $\fgg$ has type $A$, $\fm^-$ identifies with a space of matrices, and
the rank as defined above of an element in $\fm^-$ coincides with its rank as a matrix.
\end{exem}

The following proposition is a characterization of the rank $\rky$ of $y\in\fm^-$ in terms
of its action on $\bE$. If $y\in\fgg$, we will denote by $\by\in\End(\bE)$ the image of $y$ by the
representation of $\fgg$ on $\bE$.

\begin{prop}
\label{prop:rang-y}
Let $y = y_{\alpha_1} + \cdots + y_{\alpha_\rky}$ with $(\alpha_1,\ldots,\alpha_\rky)$ a
dominant orthogonal sequence.
We have $\by^{\rky+1}=0$ and $\by^\rky(\bE_{0}) = \bE_{\varpi-\alpha_1-\cdots-\alpha_\rky}$, so
in particular $\by^\rky(\bE_{0}) \neq \{ 0 \}$.
In other words, the rank of $y$ is the order of nilpotency of $\by$.
\end{prop}

\begin{proof}
We have $y_{\alpha_i} \in \fgg_{-\alpha_i}$. For any pair $(\chi,\delta)$ with $\chi \in X(\bE)$ and
$\delta$ a long root,
$\chi-2\delta$ cannot be a weight of $\bE$ by Lemma~\ref{lemm:gV}. We deduce from
Proposition~\ref{prop:cascade}\,\eqref{prop:cascade1}
that
$\by_{\alpha_i}^2=0$. On the other hand, for any $i,j$ with $i \neq j$,
the maps $\by_{\alpha_i}$ and $\by_{\alpha_j}$ commute since $\alpha_i + \alpha_j$ is not
a root by Proposition~\ref{prop:cascade}\,\eqref{prop:cascade1} also. Thus we get
$$
\by^k = k!\ \sum \by_{\alpha_{j_1}} \circ \cdots \circ \by_{\alpha_{j_k}},
$$
where the sum is over the increasing sequences $1\leq j_1 < \cdots < j_k \leq \rky$.

In particular $\by^{\rky+1} = 0$, and
$\by^\rky = \rky!\ \by_{\alpha_1} \circ \cdots \circ \by_{\alpha_\rky}$.
By Lemma~\ref{lemm:gV}, we have
$\by_{\alpha_1} \circ \cdots \circ \by_{\alpha_\rky} \cdot \bE_{\varpi} = \bE_{\varpi-\alpha_1-\cdots-\alpha_\rky}$,
so the proposition is proved.
\end{proof}

For the remaining of this section, we let $1\leq \rgy \leq \rkm$ be an integer and $y=y_r:=y_{\alpha_1}+
\cdots + y_{\alpha_\rgy}\in\fm^-$, where
$(\alpha_1,\ldots,\alpha_\rgy)$ is the dominant orthogonal sequence of length $\rgy$ obtained by the
algorithm of Definition~\ref{def:cascade}. By Proposition~\ref{prop:rank}, there is no loss of
generality in considering this particular element $y$.
We denote by
$h_\rgy$ the element
$\alpha_1^\vee + \cdots + \alpha_\rgy^\vee$ spelled out in Table~\ref{equa:tableau}. Recall that $\by$
is $y$ acting on $\bE$ via the cominuscule representation of $\fgg$, and that $r$ is the order
of nilpotency of $\by$.

\subsubsection{Definition of $\,\Vr$ and the Higgs property}

It is a very well-known idea that the nilpotent element $\by$ of $\End(\bE)$ uniquely defines an increasing
filtration $\Wr_\sbullet$ of~$\bE$, called the \emph{Deligne weight filtration of $\by$}, or
\emph{Jacobson-Morozov filtration}. This idea appears in \citestar{deligne}*{\textsection 1.6.13} in the context of
Hodge theory and in \citestar{mcgovern}*{\S\S3.2\,\&\,3.4} in the context of nilpotent orbits in Lie algebras.

The defining properties of the filtration $(\Wr_k)_{k\in\bZ}$ are that for all $k\in\bZ$,
$$
\by(\Wr_k)\subset\Wr_{k-2}
$$
and for all $k\geq 1$,
\begin{equation}
\label{equa:deligne}
\by^k:\Wr_k/\Wr_{k-1}\to\Wr_{-k}/\Wr_{-k-1}\,\mbox{ is an isomorphism.}
\end{equation}
The filtration $(\Wr_k)_{k\in\bZ}$ can be defined as follows (see \eg \cite{SteenbrinkZucker}):
\begin{equation}
\label{equa:filtration-y}
\Wr_k=\sum_{\substack{\ell \geq 0 \\ k+\ell+1 \geq 0}}
\Ker \by^{k+\ell+1} \cap \Im \by^\ell = \sum_{\substack{i-j=k \\
i\geq 0,\, j\geq 0}}
\Ker \by^{i+1} \cap \Im \by^j.
\end{equation}
Since $\rgy$ is the rank of nilpotency of $\by$ acting on $\bE$, we have $\Wr_{k}=\{0\}$ for $k<-\rgy$ and $\Wr_k=\bE$ for $k\geq \rgy$.

The elements $y$ and $h_\rgy$ fit in an $\fsl_2$-triple $(x,h_\rgy,y)$ for some $x\in\fm^+$, and the subspaces $\Wr_k$ can be alternatively defined by means of this triple:
\begin{equation}
\label{equa:filtration-poids}
\Wr_k=\bigoplus_{\chi \in X(\bE) : \scal{\chi,h_\rgy}\leq k}\bE_\chi.
\end{equation}

We combine the gradation $\bE=\bigoplus_{i=0}^\rkm \bE_i$ and the weight filtration $(\Wr_k)_{k\in\bZ}$
to define the subspaces of $\bE$ that will be the center of our attention for the remaining part of
Section~\ref{sec:subnil}.

\begin{nota}\label{defi:V}
Let $\Vr_i=\bE_i\cap\Wr_{r-2i}$ for $0\leq i\leq \rkm$,
and let $\Vr=\bigoplus_i\Vr_i\subset\bE$.
\end{nota}

Observe that $\Vr_0=\bE_0$ and that $\Vr_i\neq\{0\}$ only if $i\leq r$. The subspaces $\Vr_i$ can also be described as sums of weight spaces $\bE_\chi$.

\begin{prop}\label{prop:chi-Ei}
For $\chi\!\in\!X(\bE_i)$, we have $\scal{\chi, h_\rgy}\!\geq\!\rgy - 2i$. Hence, for all $i$ such that $0\leq i\leq r$,
$$
\Vr_i=\bigoplus_{\substack{\chi\in X(\bE_i):\\\scal{\chi,h_\rgy}=r-2i}}\bE_\chi=\bigoplus_{\substack{\chi\in X(\bE):\\ \scal{\chi,z}=\zmax-2i,\\\scal{\chi,h_\rgy}=r-2i}}\bE_\chi.
$$
\end{prop}

\begin{proof}
Proposition~\ref{prop:Ei} and the fact that, for any root $\alpha$,
we have $\scal{\alpha, h_\rgy} \leq 2$ (see Proposition~\ref{prop:h}) imply the first assertion. It is plain from the description of the subspaces $\Wr_k$ given in~\eqref{equa:filtration-poids}
that the second assertion follows from the first.
\end{proof}

The following lemma is reminiscent of
Equation~\eqref{equa:deligne}. We will prove a stronger statement in
Proposition~\ref{prop:V-self-dual}. More Hodge-theoretic considerations will be made in
Section~\ref{subsec:hodge}.

\begin{lemm}
\label{lemm:dim-Vi}
The linear map $\by^{\rgy-2i}$ is an isomorphism between $\Vr_i$ and $\Vr_{\rgy-i}$. In~particular,
$\dim \Vr_{\rgy-i} = \dim \Vr_i$.
\end{lemm}

\begin{proof}
We know that $$\Vr_i=\bE_i\cap\Wr_{\rgy-2i},\quad
\Vr_{\rgy-i}=\bE_{\rgy-i}\cap\Wr_{2i-\rgy},$$ and that for $i$ satisfying \hbox{$0\leq i\leq \rgy/2$}, $\by^{\rgy-2i}$ is
an isomorphism between
$\Wr_{\rgy-2i}/\Wr_{\rgy-2i-1}$ and $\Wr_{2i-\rgy}/\Wr_{2i-\rgy-1}$. By Proposition~\ref{prop:Ei},
$\by^{\rgy-2i}$ maps $\bE_i$ to $\bE_{\rgy-i}$. Since we have \hbox{$\bE_k\!\cap\!\Wr_{\rgy-2k-1}\!=\!\{0\}$} for all $k$ by
Proposition~\ref{prop:chi-Ei} and Equation~\eqref{equa:filtration-poids}, we get that~$\by^{\rgy-2i}$ is an isomorphism between $\Vr_i$ and~$\Vr_{\rgy-i}$.
\end{proof}

\begin{defi}
\label{def:yHiggs}
A subspace $\bS$ of $\bE$ is a \emph{$y$-Higgs subspace} if $y\cdot\bS\subset\bS$ and
$\fm_+\cdot\bS\subset\bS$.
\end{defi}

The following algebraic fact is at the heart of the construction of leafwise Higgs subsheaves
in Section~\ref{sec:MW}.

\begin{prop}
\label{prop:algHiggs}
The subspace $\Vr$ is a $y$-Higgs subspace of $\bE$. More precisely, the following inclusions hold:
\begin{itemize}
\item $y \cdot \Vr_i \subset \Vr_{i+1}$.
\item $\fm^+ \cdot \Vr_i \subset \Vr_{i-1}$.
\end{itemize}
\end{prop}
\begin{proof}
Using the description of $\Vr_i$ given in Proposition~\ref{prop:chi-Ei}, the first inclusion
holds because $y \in \bigoplus_{i} \fgg_{-\alpha_i}$ and each $\alpha_i$ satisfies
$\scal{-\alpha_i,z}=\scal{-\alpha_i,h_\rgy}=-2$. The second one also holds since for $\alpha$ a root of
$\fm^+$, we have $\scal{\alpha,z}=2$ (and $\scal{\alpha,h_\rgy} \leq 2$ by Proposition~\ref{prop:h}).
\end{proof}

\subsubsection{The module structure of $\,\Vr$}

We now show that the subspace $\Vr$ of $\bE$ has the structure of an irreducible $L_\rgy$-module, for
$L_\rgy$ a complex reductive group
determined by $y$, whereas the $\Vr_i$ are irreducible $H_\rgy$-modules, where $H_\rgy=L_\rgy\cap K$.

We begin by defining the relevant subalgebras of $\fgg$.

\begin{nota}\label{defi:qlh}
Let $\fq_\rgy \subset \fgg$ be the parabolic subalgebra defined by the coweight $z-h_\rgy$:
$$ \fq_\rgy = \ft \oplus \bigoplus_{\alpha:\scal{\alpha,h_\rgy} \leq \scal{\alpha,z}} \fgg_\alpha . $$
A Levi factor of $\fq_\rgy$ is
$$\fl_\rgy=\ft \oplus \bigoplus_{\alpha:\scal{\alpha,h_\rgy} = \scal{\alpha,z}} \fgg_\alpha . $$
Let $\fl_\rgy^\pm \subset \fl_\rgy$ be the nilpotent subalgebras of $\fl_\rgy$ defined by
$$\fl_\rgy^\pm = \bigoplus_{\alpha:\scal{\alpha,h_\rgy} = \scal{\alpha,z} = \pm 2} \fgg_\alpha . $$
The intersection $\fk\cap\fq_\rgy$ is a parabolic subalgebra of $\fk$ and a Levi factor of it is
$$\fh_\rgy = \ft \oplus \bigoplus_{\alpha:\scal{\alpha,h_\rgy} = \scal{\alpha,z}=0} \fgg_\alpha .$$

We denote by $Q_\rgy$, $L_\rgy$ and $H_\rgy$, the connected subgroups of $G$ with
Lie algebra $\fq_\rgy$,~$\fl_\rgy$ and~$\fh_\rgy$ respectively. Such groups exist because $\fq_\rgy$ is a parabolic subalgebra, and
$\fh_\rgy$ and $\fl_\rgy$ are Levi subalgebras.
We observe that
$H_\rgy$ is a Levi subgroup of $K \cap Q_\rgy$: in fact, we have $K\cap Q_\rgy = H_\rgy \cdot R(K \cap Q_\rgy)$, where
$R(K \cap Q_\rgy)$ denotes the unipotent
radical of $K \cap Q_\rgy$. Since Levi subgroups correspond to Levi subalgebras,
$L_\rgy$ is a Levi factor of $Q_\rgy$.
\end{nota}

For the convenience of the reader, the conditions that define the different Lie subalgebras of $\fgg$
we are
considering are displayed in Table~\ref{equa:tableau-groupes}.

\begin{table}[ht]
$$\arraycolsep5pt\renewcommand{\arraystretch}{1.3}
\begin{array}{|l|c|c|c|c|c|c|}
\hline \mbox{Lie algebra} & \fk & \fm^\pm & \fq_\rgy & \fl_\rgy & \fl_\rgy^\pm & \fh_\rgy \\
\hline \mbox{Condition} & z=0 & z = \pm 2 & h_\rgy \leq z & h_\rgy=z & h_\rgy=z=\pm 2 & h_\rgy=z=0 \\
\hline
\end{array}
$$
\caption{\smaller Lie subalgebras under consideration (We abbreviate the condition
$\scal{\alpha,z}=0$,
\resp $\scal{\alpha,h_\rgy}=0$, on a root $\alpha$ as $z=0$, \resp $h_\rgy=0$.)}\label{equa:tableau-groupes}
\end{table}

We have the following lemmas concerning the $\Vr_i$'s.
\begin{lemm}
\label{lemm:Viplus1}
We have $\Vr_{i+1} = \fl_\rgy^- \cdot \Vr_i$ and $\Vr_{i-1} = \fl_\rgy^+ \cdot \Vr_i$.
\end{lemm}
\begin{proof}
Let us denote by $X(\bE_i)$ the set of weights
of $\bE_i$. We know by Proposition~\ref{prop:Ei} that $\bE_{i+1} = \fm^-\cdot \bE_i$. Thus,
$$
\bE_{i+1} = \biggl (\bigoplus_{\alpha:\scal{\alpha,z} = -2} \fgg_\alpha \biggr) \cdot
\biggl (\bigoplus_{\chi \in X(\bE_i)} \bE_\chi \biggr) =
\bigoplus_{\substack{\chi \in X(\bE_i)\\ \alpha \in R(\fm^-)}} \bE_{\chi+\alpha} .
$$
Given $\chi \in X(\bE_i)$ and $\alpha \in R(\fm^-)$, we can make two observations:
\begin{itemize}
\item If $\scal{\chi,h_\rgy} > r - 2i$ then we have $\scal{\chi+\alpha,h_\rgy} > \rgy -2(i+1)$ and thus
$\bE_{\chi + \alpha} \not \subset \Vr_{i+1}$.
\item If $\scal{\alpha,h_\rgy} > -2$ then $\scal{\chi+\alpha,h_\rgy} > \rgy -2(i+1)$ and thus
$\bE_{\chi + \alpha} \not \subset \Vr_{i+1}$.
\end{itemize}
The first equality of the lemma follows from these observations.
The proof of the second equality is similar.
\end{proof}

Note that $\fl_\rgy^-$ is an $\fh_\rgy$-module. More precisely, we have

\begin{lemm}
\label{lemm:Vi-irr}
The modules $\Vr_i$ are irreducible $\fh_\rgy$-modules.
The Lie algebra $\fl_\rgy^-$ is an irreducible $\fh_\rgy$-module.
\end{lemm}
\begin{proof}
Denote by $\fk^+$, \resp $\fh_\rgy^+$, the sum of the root spaces for positive roots in~$\fk$, \resp $\fh_\rgy$.
Let $i$ be fixed and such that $\Vr_i \neq \{0\}$. By Proposition~\ref{prop:Ei}, $\bE_i$ is an
irreducible $\fk$-module.
Let $\mu_i$ be the lowest weight of $\bE_i$. We have $\bE_{\mu_i} \subset \Vr_i$.
Since $\bE_i$ is irreducible, we have $\bE_i = U(\fk^+) \cdot \bE_{\mu_i}$ (here $U$ denotes the
universal enveloping algebra). Now, arguing as in the proof
of the Lemma~\ref{lemm:Viplus1}, we see that this implies that $\Vr_i = U(\fh_\rgy^+) \cdot
\bE_{\mu_i}$. Now, $\bE_{\mu_i}$ has dimension 1, hence $\Vr_i$ is indecomposable. Since $\fh_\rgy$ is reductive, this proves that $\Vr_i$ is irreducible.
Now, by Lemma~\ref{lemm:Viplus1} again, we have $\Vr_1 = \fl_\rgy^- \cdot \Vr_0 \simeq \fl_\rgy^- \cdot \bE_0$: thus
$\fl_\rgy^-$ is also irreducible.
\end{proof}

Concerning the submodule $\Vr$, we have

\skpt
\begin{prop}
\label{prop:stabV}
\begin{enumerate}
\item \label{item:parabolic}
The subspace $\Vr \subset \bE$ is an irreducible $Q_\rgy$-module.
\item \label{item:reductive-Q}
As a consequence, the nilpotent radical of $Q_\rgy$ acts trivially on $\Vr$.
\item $\Vr_i$ is a $K\cap Q_\rgy$-module.
\item \label{item:stabilizer-V}
The subgroup of elements in $G$ which preserve $\Vr$ is exactly $Q_\rgy$.
\end{enumerate}
\end{prop}

\begin{proof}
In characteristic $0$, given a representation of a connected group, a subspace of this representation
is stable under this group if and only if it is stable for the induced action of the Lie
algebra. Thus, to prove the proposition, it is enough to deal with the corresponding Lie algebras.

Let us prove that $\Vr$ is stable under $\fq_\rgy$. It is clearly stable under $\ft$.
Let $\alpha$ be such that $\fgg_\alpha\subset\fq_\rgy$, \ie $\scal{\alpha,h_\rgy} \leq \scal{\alpha,z}$.
Let $v \in \bE_\chi \subset \Vr_i$, so that we have $\scal{\chi,h_\rgy} = \rgy-2i$.
For $x \in \fgg_\alpha$, we have $x \cdot v \in \bE_{\chi+\alpha} \subset \bE_{i-\scal{\alpha,z}/2}$. Since
$\scal{\alpha,h_\rgy} \leq \scal{\alpha,z}$, we have $\scal{\chi + \alpha, h_\rgy} \leq \rgy-2i+\scal{\alpha,z}$,
thus either $x \cdot v = 0$ or $\scal{\chi + \alpha, h_\rgy} = \rgy-2i+\scal{\alpha,z}$, by
Proposition~\ref{prop:chi-Ei}. In the second case, we get $x \cdot v \in \Vr_{i-\scal{\alpha,z}/2}$.
Combining Lemmas~\ref{lemm:Viplus1} and~\ref{lemm:Vi-irr}, we see that $\Vr$ is an irreducible
$\fl_\rgy$-module, hence also an irreducible $\fq_\rgy$\nobreakdash-module.

It follows from Engel's theorem that there exists a non-zero vector in $\Vr$ which is annihilated by the
nilpotent radical of $\fq_\rgy$. Since this radical is an ideal and $\Vr$ is an irreducible $\fq_\rgy$-module,
this radical acts trivially on $\Vr$.

The fact that $\Vr_i$ is a $K\cap Q_\rgy$-module follows because $\Vr_i = \bE_i \cap \Vr$, $\bE_i$ is $K$-stable, and
$\Vr$ is $Q_\rgy$-stable.

Finally, to prove that the stabilizer of $\Vr$ is exactly $Q_\rgy$, let us denote by
\hbox{$\fstab(\Vr)\!\subset\!\fgg$} the Lie subalgebra
preserving the subspace $\Vr$ in $\bE$.
We know by \eqref{item:parabolic} that \hbox{$\fstab(\Vr)\!\supset\!\fq_\rgy$}.
On the other hand, let $\alpha$ be a root
such that $\fgg_\alpha \cdot \Vr \subset \Vr$.
Let us assume as a first case that $\scal{\alpha,z}=-2$, \ie $\fgg_\alpha\subset\fm^-$. Since, by
Proposition~\ref{prop:Ei}\eqref{prop:Eid}, the action of $\fgg$ on~$\bE$ induces
an isomorphism $\bE_1 \simeq \bE_\varpi \otimes \fm^-$, we have $\fgg_\alpha \cdot \Vr_0=\fgg_\alpha \cdot \bE_{\varpi} = \bE_{\varpi+\alpha}$.
Then $\bE_{\varpi+\alpha} \subset \Vr\cap\bE_1=\Vr_1$, so $\scal{\alpha,h_\rgy}=-2$, and $\fgg_a\subset\fq_\rgy$.
Assume now that $\scal{\alpha,z}=0$, \ie $\fgg_\alpha\subset\fk$, and by contradiction that $\scal{\alpha,h_\rgy}>0$.
Proposition~\ref{prop:h} then implies that $\scal{\alpha,h_\rgy}=1$.
For any integer $1\leq i\leq r$,
we cannot have $\scal{\alpha,\alpha_i^\vee}=2$ because this would imply $\alpha=\alpha_i$ and
$\scal{\alpha,h_\rgy}=2$. Let $i$ be such that $\scal{\alpha,\alpha_i^\vee}>0$: then
$\scal{\alpha,\alpha_i^\vee}=1$. Therefore, $\alpha-\alpha_i$ is a root.
Since $\fgg_\alpha\subset\fk$, $\fgg_\alpha\cdot \bE_i\subset\bE_i$ and hence
$\fgg_\alpha\cdot\Vr_i\subset\Vr_i$. Furthermore $\fgg_\alpha(\Vr_0)=0$ since $\Vr_0=\bE_\varpi$.
It follows that
$$
\fgg_{\alpha-\alpha_i} \cdot \bE_{\varpi} = [ \fgg_\alpha, \fgg_{-\alpha_i} ] \cdot \bE_{\varpi}
= \fgg_\alpha \cdot (\fgg_{-\alpha_i} \cdot \bE_{\varpi}).
$$
Again by Proposition~\ref{prop:Ei}\eqref{prop:Eid}, $\fgg_{\alpha-\alpha_i} \cdot \bE_{\varpi} = \bE_{\varpi + \alpha -
\alpha_i}$ and $\fgg_{-\alpha_i}\cdot\bE_\varpi= \bE_{\varpi-\alpha_i}$. But $\bE_{\varpi-\alpha_i}\subset \Vr_1$
whereas $\bE_{\varpi + \alpha - \alpha_i}\not\subset\Vr_1$ since we assumed $\scal{\alpha,h_\rgy}>0$.
This contradicts $\fgg_\alpha \cdot \Vr \subset \Vr$.

Let now $\Stab(\Vr) \subset G$ be the subgroup stabilizing $\Vr$. We have $Q_\rgy \subset \Stab
(\Vr)$, thus $\Stab(\Vr)$ is parabolic and therefore connected \citestar{humphreys}*{Cor.\,23.1.B}. It
follows that $Q=\Stab(\Vr)$.
\end{proof}

\subsubsection{The self-duality property. Tube type Hermitian symmetric spaces}
\label{subsubsec:tube}

By definition, a Hermitian symmetric space $M=G_\bR/K_\bR$ is of \emph{tube type} if it is biholomorphically equivalent to a domain $F\oplus i C$, where
$C \subset F$ is a proper open cone in the real vector space $F$. The simplest example is the upper
half-plane, where $F=\bR$ and $C=\bR_+^*$. For irreducible symmetric spaces,
this happens when $G_\bR$ is isogenous to $\SU(p,p)$, $\Spin(2,m)$, $\Sp(2m,\bR)$,
$\SO^\star(2m)$ for $m$ even, and $\mathrm{E}_{7(-25)}$.

We introduce the following definition.

\begin{defi}
\label{defi:admissible-sequence}
A sequence $(\gamma_1,\ldots,\gamma_r)$ of roots will be called \emph{admissible} if all the roots $\gamma_i$ are long
and bigger than $\zeta$, and if they are pairwise strongly orthogonal.
\end{defi}

Note that the dominant orthogonal sequences $(\alpha_1,\ldots,\alpha_r)$ are admissible.
We can now state:

\begin{prop}
\label{prop:self-dual}
The following assertions are equivalent:
\begin{enumerate}
\item \label{item:tube} $M$ has tube type;
\item \label{item:varpi-forall}
For all admissible sequences $\gamma_1,\ldots,\gamma_\rkm$, we have $\gamma_1+\cdots+\gamma_\rkm = 2 \varpi$;
\item \label{item:varpi-a}
$\alpha_1+\cdots+\alpha_\rkm = 2 \varpi$;
\item \label{item:varpi-exists}
There exists an admissible sequence $\gamma_1,\ldots,\gamma_\rkm$ such that $\gamma_1+\cdots+\gamma_\rkm = 2 \varpi$;
\item \label{item:z}
$\alpha_1^\vee + \cdots + \alpha_\rkm^\vee = z$;
\item \label{item:self-dual}
$\bE$ is self-dual: $\bE \simeq \bE^\star$ as $G$-modules;
\item \label{item:Ei}
For all $i$, $\bE_i \simeq \bE_{\rkm-i}^\star$ as $K$-modules;
\item \label{item:Er}
$\dim \bE_\rkm=1$;
\item \label{item:alpha_r}
$\alpha_\rkm=\zeta$.
\end{enumerate}
\end{prop}
\begin{proof}
Glancing at Table~\ref{equa:tableau}, one can readily
check that $G_\bR / K_\bR$ is of tube type if and only if $z=\alpha_1^\vee + \cdots + \alpha_\rkm^\vee$,
and that this occurs also exactly when $\alpha_\rkm=\zeta$. Hence \eqref{item:tube}, \eqref{item:z} and
\eqref{item:alpha_r} are equivalent.

The equality $z=a_1^\vee+\cdots+\alpha_\rkm^\vee$ is equivalent to $\scal{\suma{\rkm},\beta} = 2
\delta_{\beta,\zeta}$ for all
$\beta\in\Pi$ because $\zeta$ is the only noncompact simple root and it has the same length as each
$\alpha_i$ (they are all long).
Since all the roots $\alpha_i$ have the same length, for any root~$\beta$,
the equation $\scal{\suma{\rkm},\beta}=0$ is equivalent to
$\scal{\alpha_1 + \cdots + \alpha_\rkm,\beta^\vee}=0$. It follows that \eqref{item:z} and \eqref{item:varpi-a}
are equivalent.

We now prove that items \eqref{item:varpi-forall}, \eqref{item:varpi-a}, and
\eqref{item:varpi-exists} are equivalent. In fact, we have the trivial implications
$\eqref{item:varpi-forall} \Rightarrow \eqref{item:varpi-a} \Rightarrow \eqref{item:varpi-exists}$. Moreover, if
\eqref{item:varpi-exists} holds, then let $\gamma_1,\ldots,\gamma_\rkm$ be an admissible sequence
such that $\gamma_1+\cdots+\gamma_\rkm = 2\varpi$. By \citestar{richardson}*{Th.\,2.1}, any admissible
sequence can be obtained applying to this sequence an element $w$ of the subgroup of the Weyl group
that fixes $\varpi$.
Condition \eqref{item:varpi-forall} follows from the equality $w(\gamma_1)+\cdots+w(\gamma_\rkm)=w(2\varpi)=2\varpi$.

Since the weights of $\bE^\star$ are the opposite of the weights of $\bE$, the highest
weight of~$\bE^\star$ is $-\mu_0$. Therefore, \eqref{item:self-dual} holds if and only if
$\mu_0=-\varpi$, which is equivalent to~\eqref{item:varpi-a} by Proposition~\ref{prop:cascade}\,\eqref{prop:cascade4}.

Assuming~\eqref{item:self-dual} and using the decomposition into $K$-modules $\bE =
\bigoplus_{i=0}^{\rkm} \bE_i$, we deduce that
$\bE_{\rkm-i} \simeq \bE_{i}^\star$ as $K$-modules. Conversely, given \eqref{item:Ei}, we obviously
have \eqref{item:Er}.
Assuming now \eqref{item:Er}, we deduce that $\bE_\rkm=\bE_{\mu_0}$. Since $\bE_{\rkm}$ is a
$K$-module of dimension~$1$ in this
case, it follows that $\mu_0$ vanishes on all the roots of $\fk$, so $\mu_0$ is an integral
multiple of $\varpi$. It follows
that $\mu_0=-\varpi$ and we get \eqref{item:varpi-a} as above.
\end{proof}

As we just said, the $G$-module $\bE$ is not always self-dual. However, the subspace $\Vr$ is
always a self-dual module under a subgroup $\LB$ of $G$ that we now define.

Recall that the Lie algebra $\fl_r$ is by definition generated by the root spaces of the roots $\alpha$
such that $\scal{h_r,\alpha}=\scal{z,\alpha}$, and that $\fl_r^\pm$ is the subalgebra generated by the root spaces of the roots $\alpha$
such that $\scal{h_r,\alpha}=\scal{z,\alpha}=\pm 2$. We define $\flb$ as the Lie subalgebra of $\fgg$
generated by the spaces $\fl_r^+$ and $\fl_r^-$. Two particular cases are easy:

\begin{fact}
\label{fact:flb-easy}
If $r=1$, then $\fsl_2 \simeq \flb \subset \fgg$ is the Lie subalgebra corresponding to the highest root $\Theta$. If $M$ has tube
type and $r=\rkm$, then $\flb = \fgg$.
\end{fact}
\begin{proof}
In case $r=1$, then $h_r=\Theta^\vee$, so $\scal{\alpha,h_r}=2$ implies $\alpha = \Theta$. Thus $\fl_r^+ = \fgg_\Theta$ in this case.
When $M$ has tube type and $r=\rkm$, we have $h_r=z$ (Proposition~\ref{prop:self-dual}\,\eqref{item:z}), so $\fl_r^+ = \fm^+$ in this case.
\end{proof}

We now describe the Lie algebra $\flb$ in general. To this end, we introduce a definition.

\begin{defi}
\label{defi:connected-to}
Let $R$ be a (possibly reducible) root system, let $\Pi$ be a basis of~$R$, and $\varpi$ be a dominant weight. Given a simple root $\beta$,
we say that $\beta$ is connected to~$\varpi$ in $R$ if there is a simple root $\gamma$ in the same connected component as
$\beta$ in the Dynkin diagram of $R$ and such that $\scal{\varpi,\gamma^\vee}>0$.
\end{defi}

We denote by $\Pi_{h_r=0}$ the set of simple roots $\beta \in \Pi$ such that $\scal{\beta,h_r}=0$.

\begin{lemm}
\label{lemm:flb}
The Lie algebra $\flb$ is simple.
The lowest root $-\Theta$ together with the simple roots which are connected to $\Theta$ in $\Pi_{h_r=0}$
form a basis of the root system of $\flb$.
\end{lemm}
\begin{proof}
Recall that $\zeta$ denotes the only noncompact simple root (see Notation~\ref{nota:zeta}).
Looking at Table~\ref{equa:tableau}, we see that the root $\zeta$
cannot be connected to $\Theta$ in $\Pi_{h_r=0}$.
Thus, if $\gamma$ is connected to $\Theta$ in $\Pi_{h_r=0}$, it satisfies $\scal{\gamma,z}=0$ and,
of course, $\scal{\gamma,h_r}=0$.
Let us assume that $\scal{\Theta,\gamma^\vee}>0$. Then $\Theta-\gamma$ is a root and it satisfies
$\scal{h_r,\Theta-\gamma}=\scal{z,\Theta-\gamma} =2$, so $\fgg_{\Theta-\gamma} \subset \flb$, and so
$\fgg_{\gamma} \subset \flb$. By induction, the root spaces of all simple roots connected to $\Theta$ and their
opposite are in $\flb$,
and, together with $\gamma_{-\Theta}$ and~$\gamma_{\Theta}$, they generate $\flb$.

In general, the Dynkin diagram of a Lie algebra has vertices corresponding to the simple roots of
the Lie algebra.
In our case, given the above description of a basis of $\flb$, the set of vertices of its Dynkin diagram
is the union of some connected components of $\Pi_{h_r=0}$ and the set $\{-\Theta\}$. Moreover, all the
connected components of $\Pi_{h_r=0}$ which occur are connected to $-\Theta$.
Thus the Dynkin diagram of $\flb$ is connected
and $\flb$ is simple.
\end{proof}

In Table~\ref{equa:flb}, we indicate in each case which root system for $\flb$ is obtained (the
cases where $r=1$ or $M$ has
tube type and $r=\rkm$ are not represented, see Fact~\ref{fact:flb-easy}).

\begin{table}[ht]
\pgfkeys{/Dynkin diagram}%,edgeLength=6mm}
$$
\begin{array}{|c|c|c|}
\hline
\multirow{2}{*}{$(G,\varpi)$} & \multirow{2}{*}{$\flb$} & \multirow{2}{*}{$R(\flb)$}
\\ & & \\ \hline
\multirow{3}{*}{$(A_{n-1},\varpi_l)$} &
\multirow{3}{*}{\begin{tikzpicture}
\dynkin{A}{o.oo*oo.o}
\dynkinLabelRoot{1}{\beta_{r-1}}
\dynkinLabelRoot{2}{\beta_2}
\dynkinLabelRoot{3}{\beta_1}
\dynkinLabelRoot{4}{-\Theta}
\dynkinLabelRoot*{4}{~}
\dynkinLabelRoot{5}{\beta_n}
\dynkinLabelRoot{6}{\beta_{n-1}}
\dynkinLabelRoot{7}{\beta_{n-r+1}}
\end{tikzpicture}}
&\multirow{3}{*}{ $A_{2r-1}$}
\\ &&\\ &&\\ \hline
\multirow{3}{*}{$(C_n,\varpi_n)$} &
\multirow{3}{*}{\begin{tikzpicture}
\dynkin{C}{o.o*}
\dynkinLabelRoot{1}{\beta_{r-1}}
\dynkinLabelRoot{2}{\beta_1}
\dynkinLabelRoot{3}{-\Theta}
\dynkinLabelRoot*{3}{~}
\end{tikzpicture}}
& \multirow{3}{*}{$C_r$}
\\ &&\\ &&\\ \hline
\multirow{4}{*}{$(D_n,\varpi_n)$} &
\multirow{4}{*}{\begin{tikzpicture}
\dynkin{D}{o.ooo*}
\dynkinLabelRoot{1}{\beta_{2(r-1)}}
\dynkinLabelRoot{2}{\beta_3}
\dynkinLabelRoot{3}{\beta_2}
\dynkinLabelRoot{4}{\beta_1}
\dynkinLabelRoot{5}{-\Theta}
\end{tikzpicture}}
& \multirow{4}{*}{$D_{2r-1}$}
\\ &&\\ &&\\ &&\\\hline
\multirow{4}{*}{$(E_6,\varpi_{1})$, $r=2$} &
\multirow{4}{*}{\begin{tikzpicture}
\dynkin{D}{*oooo}
\dynkinLabelRoot{1}{-\Theta}
\dynkinLabelRoot{2}{\beta_2}
\dynkinLabelRoot{3}{\beta_3}
\dynkinLabelRoot{4}{\beta_1}
\dynkinLabelRoot{5}{\beta_4}
\end{tikzpicture}}
& \multirow{4}{*}{$D_5$}
\\ &&\\ && \\ &&\\ \hline
\multirow{4}{*}{$(E_7,\varpi_7)$, $r=2$} &
\multirow{4}{*}{\begin{tikzpicture}
\dynkin{D}{*ooooo}
\dynkinLabelRoot{1}{-\Theta}
\dynkinLabelRoot{2}{\beta_1}
\dynkinLabelRoot{3}{\beta_3}
\dynkinLabelRoot{4}{\beta_4}
\dynkinLabelRoot{5}{\beta_5}
\dynkinLabelRoot{6}{\beta_2}
\end{tikzpicture}}
& \multirow{4}{*}{$D_6$}\\
&&\\ &&\\ &&\\ \hline
\end{array}
$$
\caption{\smaller The Lie algebra $\flb$ and its root system (The cases where $r=1$ or $M$ has
tube type and $r=\rkm$ are not represented.)}
\label{equa:flb}
\end{table}

We consider the root system $R_\flb$ with its basis $(-\Theta,\beta_1,\ldots,\beta_k)$ given by
Lem\-ma~\ref{lemm:flb}. We have

\begin{lemm}
\label{lemm:theta-cominuscule}
The simple root $-\Theta$ in $R(\flb)$ is cominuscule.
\end{lemm}

In other words, the lemma says that any root in $R(\flb)$ has coefficient at most one on the simple
root $-\Theta$. Observe that, as a consequence, $-\Theta$ defines a Hermitian noncompact real
form $\ov{L}_{r,\bR}$ of $\LB$.
The element $y=y_r$ can be considered as an element of $\fl_r^-$.
\begin{proof}
Since any root in $R(\flb)$ can be obtained from $\Theta$ applying a suitable element of the Weyl group of $\flb$, it is enough to
show the implication
$$ |n_{-\Theta}(\alpha)| \leq 1 \Longrightarrow |n_{-\Theta}(s_\gamma(\alpha))| \leq 1,$$
for any simple root $\gamma$ in $R(\flb)$. To prove this, we assume without loss of generality that
$n_{-\Theta}(\alpha) \geq 0$.
If $\gamma$ is one of the simple roots $\beta_j$, then
$n_{-\Theta}(\alpha) = n_{-\Theta}(s_\gamma(\alpha))$, so the implication clearly holds. If $\gamma = -\Theta$, we consider
two cases. If $n_{-\Theta}(\alpha)=1$, then $\alpha$ is a negative root. We have
$0 \leq \scal{\alpha,-\Theta^\vee} \leq 2$, thus $-1 \leq n_{-\Theta}(s_\gamma(\alpha)) \leq 1$.
If $n_{-\Theta}(\alpha)=0$, then $\alpha \neq -\Theta$. Since $\Theta$ is long, it follows that
$|\scal{\alpha,-\Theta^\vee}| \leq 1$. Thus $|n_{-\Theta}(s_\gamma(\alpha))| \leq 1$, as expected.
\end{proof}

As announced, we have the following, which strengthens Lemma~\ref{lemm:dim-Vi}.

\begin{prop}
\label{prop:V-self-dual}
The $\LB$-module generated by $\bE_\varpi$ is $\Vr$. It is self-dual.
\end{prop}
\begin{proof}
First, we have $\flb \subset \fl_r$. This follows from the definition of $\flb$ and the inclusions
$\fl^+_r \subset \fl_r,\, \fl^-_r \subset \fl_r$.

From the inclusion $\flb \subset \fl_r$ and Proposition~\ref{prop:stabV}\eqref{item:parabolic},
it follows that the $\LB$-submodule generated by $\bE_\varpi$ is included in $\Vr$.
By Lemma~\ref{lemm:Viplus1}, $\Vr$ is included in the $\fl_r^-$-submodule generated by $\bE_\varpi$.
Since $\fl_r^- \subset \flb$ by
definition, we deduce that the $\LB$-module generated by $\bE_\varpi$ is exactly $\Vr$.

Now we prove that $\Vr$, as an $\LB$-module, is self-dual. The highest weight of $\Vr$ as an
$\LB$-module is the restriction of $\varpi$, which is the fundamental weight $\varpi_{-\Theta}$
corresponding to the simple root $-\Theta$ in the root system of $\flb$.
The sequence $\alpha_1,\ldots,\alpha_r$ is a sequence of long strongly orthogonal roots. They are all
bigger than $-\Theta$. In fact, if $\alpha_i$ was not bigger than $-\Theta$, then it would be a linear
combination of the roots $\beta_j$, so it would satisfy $n_\zeta(\alpha_i)=0$, contradicting the definition
of the sequence $\alpha_1,\ldots,\alpha_r$.

In view of \eqref{item:varpi-exists} in Proposition~\ref{prop:self-dual},
it is enough to prove that $\sum \alpha_i$ and $2 \varpi_{-\Theta}$ have the same pairing with the
simple roots of $R_\flb$.
For $-\Theta$, we have $\scal{\sum \alpha_i^\vee,-\Theta^\vee}=-2$ and $\scal{2\varpi_{-\Theta},-\Theta^\vee}=-2$.
For $\beta$ a simple root connected to $\Theta$ in $\Pi_{h_r=0}$, we have
$\scal{\sum \alpha_i,\beta^\vee} = 0$ because $\scal{\beta,\sum \alpha_i^\vee}=0$.
We have $\scal{\varpi_{-\Theta},\beta^\vee}=0$.
This proves that $\Vr$ is self-dual.
\end{proof}

From the equivalences in Proposition~\ref{prop:self-dual} and Proposition~\ref{prop:V-self-dual},
and recalling that
$\ovLrR$ is a Hermitian noncompact real form of $\LB$, it follows that the symmetric space
$M_\rgy:=\ovLrR/(\ovLrR\cap K_\bR)$ is a tube type Hermitian
symmetric space of rank $\rgy$.

\begin{rema}\label{rema:boundary-components}
This construction is related, but not identical, to the theory of \emph{boundary components} of
bounded symmetric domains (see
\eg \cites{WolfFineStructure,piatetski-shapiro,Satake,AshMumfordRapoportTai}). The
topological boundary of the symmetric space $M$ inside its compact dual $\check{M}$ is a union
of $r_M$ $G_\bR$-orbits, and each $G_\bR$-orbit is a union of so-called boundary components,
\cf \citestar{AshMumfordRapoportTai}*{Def.\,3.2 p.\,127}. The
components are isomorphic to Hermitian symmetric subspaces $M_\Lambda$ of $M$, where
$\Lambda$ is a subset of the set of noncompact strongly orthogonal roots, see
\citestar{AshMumfordRapoportTai}*{Th.\,3.3 p.\,127} for example. If we take our maximal dominant
orthogonal sequence $(\alpha_1,\ldots,\alpha_\rkm)$ as the set of strongly orthogonal roots, and
$\Lambda=\{\alpha_1,\ldots,\alpha_\rgy\}$ as a subset, boundary components corresponding to this subset
are symmetric spaces of rank $\rgy$, as is the subspace $M_\rgy$,
but they are not of tube type unless $M$ already is. The Lie algebra of the stabilizer of
$M_\Lambda$ in $G_\bR$ is $\fq_\Lambda\cap\fgg_\bR$, where $\fq_\Lambda$ is the parabolic subalgebra of $\fgg$
corresponding to the coweight $h_\rkm-h_\rgy$, whereas we recall that the parabolic subalgebra
$\fq_\rgy$ we are considering here corresponds to the coweight $z-h_\rgy$. This shows that
$\fq_\rgy=\fq_\Lambda$ if and only if $M$ is of tube type.
\end{rema}

\begin{exem}\label{exem:typeA}
We give an example of this construction.
We assume that we are in the first row of Table~\ref{equa:tableau}, namely that $G$ has
type $A_{p+q-1}$ for some positive integers $p\leq q$ and that $\varpi = \varpi_p$. The group $G_\bR$
is therefore $\SU(p,q)$ acting on $\bC^{p+q}$ preserving a Hermitian form $h$ of signature $(p,q)$ given
by the diagonal matrix $\diag(-1,\ldots,-1,1,\ldots,1)$ in the canonical basis
$(e_1,\ldots,e_p,e_{p+1},\ldots,e_{p+q})$.

Let
$y\in\fm^-$ be an element of rank $\rgy$. We have a natural decomposition $\bC^{p+q}=A \oplus N
\oplus B \oplus I$, with $A$, \resp $N,B,I$ generated by $(e_1,\ldots,e_\rgy)$, \resp $(e_{\rgy+1},\ldots,e_p)$, $(e_{p+1},\ldots,e_{p+q-\rgy})$, $(e_{p+q-\rgy+1},\ldots,e_{p+q})$, so
that $y$ is given by the block matrix
$$\arraycolsep3.5pt
\begin{pmatrix}
0 & 0 & 0 & 0 \\
0 & 0 & 0 & 0 \\
0 & 0 & 0 & 0 \\
* & 0 & 0 & 0
\end{pmatrix}.
$$

The element $h_\rgy \in \ft$ acts on the Lie algebra $\fsl(p+q,\bC)$ of $G$, and we represent the weights
of its action as the matrix
$$\arraycolsep3.5pt
\begin{pmatrix}
0 & 1 & 1 & 2 \\
-1 & 0 & 0 & 1 \\
-1 & 0 & 0 & 1 \\
-2 & -1 & -1 & 0
\end{pmatrix},
$$
where the blocks correspond to the decomposition $A \oplus N \oplus B \oplus I$
of $\bC^{p+q}$.
Similarly, the weights of the action of the central element $z$ are given as follows
$$\arraycolsep3.5pt
\begin{pmatrix}
0 & 0 & 2 & 2 \\
0 & 0 & 2 & 2 \\
-2 & -2 & 0 & 0 \\
-2 & -2 & 0 & 0
\end{pmatrix}.
$$
Thus the Lie algebra $\fq_\rgy$, $\fl_\rgy$ and $\fh_\rgy$ are respectively
$$\arraycolsep3.5pt
\begin{pmatrix}
* & 0 & * & * \\
* & * & * & * \\
0 & 0 & * & 0 \\
* & 0 & * & *
\end{pmatrix},\quad
\begin{pmatrix}
* & 0 & 0 & * \\
0 & * & 0 & 0\\
0 & 0 & * & 0 \\
* & 0 & 0 & *
\end{pmatrix}\ \mbox{ and } \
\begin{pmatrix}
* & 0 & 0 & 0 \\
0 & * & 0 & 0\\
0 & 0 & * & 0 \\
0 & 0 & 0 & *
\end{pmatrix}.
$$

We see that $Q_\rgy$ is exactly the stabilizer of the flag $(N \subset N \oplus A \oplus I)$, that
the subgroup $L_\rgy$ of $G$ corresponding to $\fl_\rgy$ is $\rS(\GL(N)\times\GL(A\oplus
I)\times\GL(B))$ and that~$H_\rgy$ is the block diagonal group $\rS(\GL(N) \times \GL(A) \times \GL(I) \times \GL(B))$.
Finally the intersections of $G_\bR=\SU(p,q)$ with these two latter groups are isomorphic to
$\rS(\rU(p-\rgy)\times\rU(\rgy,\rgy)\times\rU(q-\rgy))$ and $\rS(\rU(p-\rgy) \times \rU(\rgy)
\times \rU(\rgy) \times \rU(q-\rgy))$ respectively, so that
$\ov{L}_{\rgy,\bR}\simeq\SU(\rgy,\rgy)$.

On the other hand, we have $\bE=\wedge^p(N \oplus A \oplus I \oplus B)$ and it is easy to check that
$\Vr=\wedge^{p-\rgy}N \otimes \wedge^\rgy(A \oplus I)$, which illustrates that the stabilizer
of $\Vr$ is $Q_\rgy$.

\subsection{\texorpdfstring{The slope of the $H_\rgy$-modules $\Vr_i$}{The slope of the H-r modules V-r-i}}
\label{subsec:slopeVi}

Our results so far allow to compute the \emph{slope} of the $H_\rgy$-modules $\Vr_i$. We first need two
definitions and a lemma. Given an algebraic group $H$, we denote by $X(H)$ its character group
considered as a $\bZ$-module with additive law defined by: $(\chi+\psi)(h)=\chi(h)\psi(h)$.

\begin{defi}
\label{defi:polystable}
Let $H$ be a reductive group and $\bF$ be an $H$-module.
The \emph{slope} of~$\bF$ is the element $\mu_H(\bF)=\sfrac{\det(\bF)}{\dim(\bF)}$ in $X(H) \otimes_\bZ \bQ$.
We say that $\bF$ is a \emph{polystable} $H$-module if $\bF$ is a direct sum of irreducible
$H$-modules of the same slope.
\end{defi}

\begin{lemm}
\label{lemm:polystable}
Let $H$ be a reductive group.
\begin{itemize}
\item Let $\bF$ be a polystable $H$-module. For any $H$-submodule $\bS$ of $\bF$,
$\mu_H(\bS)=\mu_H(\bF)$ and $\mu_H(\bF/\bS)=\mu_H(\bF)$.
\item Let $\bF$, $\bF'$ be polystable $H$-modules. Then $\bF \otimes \bF'$ is a polystable$H$-module,
whose slope is $\mu_H(\bF)+\mu_H(\bF')$.
\item Let $\bF$ be a polystable $H$-module. Then $\bF^\star$ is a polystable
$H$-module, whose slope is $-\mu_H(\bF)$.
\end{itemize}
\end{lemm}

\begin{proof} The first assertion is clear because $H$ is reductive and so any submodule
is a direct sum of irreducible components. We prove the second assertion.
Let $Z$ be the center of $H$. It is known that restriction to $Z$ yields an injection $X(H)
\hto X(Z)$. In~fact, by \citestar{borel}*{Prop.\,14.2},
we have $H=Z \cdot (H,H)$, and any character of $H$ is trivial on $(H,H)$.

Let us first assume that $\bF$ and $\bF'$ are irreducible $H$-modules. By Schur's lemma, there are
characters $\chi,\psi$ of $Z$ such that $\forall g \in Z, \forall w \in \bF, \forall w' \in \bF', g
\cdot w = \chi(g)w$
and $g \cdot w' = \psi(g)w'$. Therefore, in $X(Z) \otimes \bQ$, we have
$\chi = \mu_H(\bF)_{|Z}$ and $\psi = \mu_H(\bF')_{|Z}$.
Let now $\bU \subset \bF \otimes \bF'$ be an irreducible component. We have
$$\forall u \in \bU,\; \forall g \in Z,\quad g\cdot u = \chi(g) \psi(g) u.$$
Therefore $\mu_H(\bU)_{|Z} = \chi + \psi$.
Thus, ${(\mu_H(\bF)+\mu_H(\bF'))}_{|Z} = \mu_H(\bU)_{|Z}$.
Since restriction of characters to $Z$ is injective, we have $\mu_H(\bU) = \mu_H(\bF) + \mu_H(\bF')$.

Let now $\bF$ and $\bF'$ be arbitrary $H$-modules, and write the decomposition into irreducible submodules:
$\bF=\bigoplus \bF_i$ and $\bF'=\bigoplus \bF'_j$. Let $i,j$ be fixed and let $\bU \subset \bF_i
\otimes \bF_j'$ be an
irreducible factor. We have proved that $\mu_H(\bU)=\mu_H(\bF_i)+\mu_H(\bF'_j)$. Since~$\bF$ and~$\bF'$
are polystable,
we have $\mu_H(\bF_i)=\mu_H(\bF)$ and $\mu_H(\bF'_j)=\mu_H(\bF')$. Thus, $\bU$ has slope
$\mu_H(\bF)+\mu_H(\bF')$ and the lemma is proved.

The third item follows from the identity $\det \bF^\star = -\det \bF$ in $X(H)$.
\end{proof}

We apply this lemma to the $H_\rgy$-modules $\Vr_i$:

\begin{prop}
\label{prop:calcul-pente}
We have $\mu_{H_\rgy}(\Vr_i) = \mu_{H_\rgy}(\Vr_0)+i\mu_{H_\rgy}({\Vr_0}^\star \otimes \Vr_1)$.
\end{prop}
\begin{proof}
The $H_\rgy$-modules $\Vr_i$ are irreducible by Lemma~\ref{lemm:Vi-irr}, thus polystable.
The same holds for the $H_\rgy$-module $\fl_\rgy^-\simeq {\Vr_0}^\star \otimes \Vr_1$, and we have
$\mu_{H_\rgy}(\fl_\rgy^-) = \mu_{H_\rgy}({\Vr_0}^\star \otimes \Vr_1)$. Let $i$ be fixed.
By Lemma~\ref{lemm:polystable}, $\Vr_i \otimes \fl_\rgy^-$ is polystable. Since, by
Lemma~\ref{lemm:Viplus1},
$\Vr_{i+1}$ is a quotient of $\Vr_i \otimes \fl_\rgy^-$, it follows that
$\mu_{H_\rgy}(\Vr_{i+1}) = \mu_{H_\rgy}(\Vr_i) + \mu_{H_\rgy}(\fl_\rgy^-)$.
\end{proof}

\Subsection{Interpretation in terms of homogeneous variations of Hodge structure}
\label{subsec:hodge}

As we learned from one of the referees, the cominuscule
representation $\bE$ of $G$ has been interpreted in the language of Hodge structures,
first by Gross~\cite{Gross} in the tube type case, then by Sheng and Zuo~\cite{ShengZuo} in
general. Indeed, the representation~$\bE$ defines a canonical homogeneous complex variation of
Hodge structure ($\bC$-VHS) $E_{CY}:=\bigoplus_{p+q=\rkm} E^{p,q}$ on $M$ or on any quotient
$\Lambda\backslash M$, for $\Lambda$ a discrete
subgroup of $G_\bR$. The Hodge bundle $E^{p,q}$ is just the bundle over $M=G_\bR/K_\bR$ associated
with the $K_\bR$-module $\bE_{\rkm-p}$.
Geometrically, the corresponding filtration is given by the successive osculating tangent spaces
to the cone over the subvariety $M \subset \bP
\bE$.
This $\bC$-VHS is of particular interest because it is of
Calabi-Yau type, meaning that $\dim
E^{\rkm,0}=1$ and $E^{\rkm-1,1}\simeq E^{\rkm,0}\otimes \Omega^1_M$. It is a real variation of
Hodge structure ($\bR$-VHS) of weight $\rkm$ precisely when $M$ is of tube type. This is exactly
the same thing as saying, as we did, that $\bE$ is self-dual in this case.

From this point of view, what we do here can be interpreted in the following way. Let us work
on the projectivized tangent bundle $\bPTM$ of $M$. The notion of rank of an element of $\fm^-$
naturally extends to elements in the tangent bundle of~$M$, and the bundle $\bPTM$ is the union
of $\rkm$ $G_\bR$-orbits ${\Orb}_1,\ldots,{\Orb}_\rkm$, where
\hbox{${\Orb}_r:=\{[\eta]\in\bPTM\mid \rk(\eta)=r\}$}. It is also the quotient
$G_\bR/(K_\bR \cap Z_r)$, where $Z_r$ is the centralizer of a nilpotent element $y$ in $\fm^-$ of rank
$r$.
The sets $\cup_{s\leq r}{\Orb}_s$ are the characteristic varieties of
Mok~\cite{Mok}. Let us fix $1\leq r\leq \rkm$ and call $\pi_M$ the projection $\bPTM\to M$.
The definition of the weight filtration $\Wr_\sbullet$ of $\bE$ associated with a nilpotent
element $y\in\fm^-$ or rank $r$ can also be extended to give an increasing holomorphic weight filtration on the
restriction of the $\bC$-VHS $\pi_M^\star E_{CY}$ to $\Orb_r$, compatible with the decreasing
Hodge filtration
$(\bigoplus_{p\geq k}\pi_M^\star E^{p,q})_k$, and this defines a homogeneous family of
``complex'' mixed Hodge structures on ${\Orb}_r$. When $M$ has tube type, $\pi_M^\star E_{CY}$
is an $\bR$-VHS and we have a real mixed Hodge structure on each fiber of $\pi_M^\star
E_{CY}$ above ${\Orb}_r$.

\begin{exem}
We keep the notations of Example~\ref{exem:typeA}, assuming now that we are in the tube type
case: the group
$G_\bR$ is $\SU(p,p)$ with maximal compact subgroup $K_\bR=\rS(\rU(p)\times\rU(p))$. Let $C$,
\resp $D$, be the subspace of $\bC^{2p}$ generated by $(e_1,\ldots,e_p)$,
\resp $(e_{p+1},\ldots,e_{2p})$. In this
case, the cominuscule representation is
$$\bE=\wedge^p(C\oplus D)=\bigoplus_{m+n=p}\wedge^m C\otimes\wedge^n D.$$
Whenever $V$ is a complex vector space, we denote by $\ov{V}$ the complex vector space
$V$ with the $\bC$-action defined by $\lambda \cdot v = \ov{\lambda}v$. Note that, with this
convention, a real structure on $V$ is the same as an involutory $\bC$-linear isomorphism
$\ov{V} \to V$.
As representations of
$K_\bR$, $\ov{C}\simeq C^\star$, $\ov{D}\simeq D^\star$, and
$\wedge^{2p}(C\oplus D)$ is
the trivial representation.
We get a $K_\bR$-equivariant isomorphism
$$
\ov{\wedge^m C \otimes \wedge^n D} \simeq \wedge^m \ov C \otimes \wedge^n \ov D
\simeq \wedge^m C^\star \otimes \wedge^n D^\star \simeq (\wedge^m C \otimes \wedge^n D)^\star
\simeq \wedge^{p-m} C \otimes \wedge^{p-n} D,
$$
where the last isomorphism is given by the perfect pairing
$$
(\wedge^m C \otimes \wedge^n D) \otimes (\wedge^{p-m}C \otimes \wedge^{p-n}D) \to
\wedge^p C \otimes \wedge^p D \simeq \wedge^{2p} (C \oplus D) \simeq \bC.
$$
Therefore we get a real
structure on $\bE$, invariant by $K_\bR$, using the isomorphism
$$\ov{\bE}=\bigoplus_{m+n=p} \ov{\wedge^m C \otimes\wedge^n D}
\simeq \bigoplus_{m+n=p} \wedge^n C\otimes\wedge^m D = \bE.$$
Since we have tautologically, for this real structure,
$\ov{\wedge^m C \otimes\wedge^n D} = \wedge^n C\otimes\wedge^m D$,
this defines an $\bR$-VHS of weight $p=\rkm$ on the symmetric space
$M=G_\bR/K_\bR$. This is the homogeneous VHS defined in \cite{Gross}.

If now $y$ is an element in $\fm^-$ of rank $\rgy$, we moreover have the decompositions $C=A\oplus N$
and $D=B\oplus I$, as in Example~\ref{exem:typeA}, where the restriction of the action of~$y$ to $A$ gives an isomorphism
$A\simeq I$. The ($p$-shifted) increasing weight filtration $\Wr_\sbullet$ of $\bE$ associated
to the action of $y$ is given by
$$
\Wr_s=\bigoplus_{\substack{i-l\leq s-p\\ i+j+k+l=p}}\wedge^iA\otimes\wedge^jN\otimes\wedge^kB\otimes\wedge^lI,
$$
whereas the decreasing Hodge filtration $\bF^\scbullet$ is given by
$$
\bF^t=\bigoplus_{\substack{i+j\geq t\\ i+j+k+l=p}}\wedge^iA\otimes\wedge^jN\otimes\wedge^kB\otimes\wedge^lI.
$$
Therefore the graded parts are:
$$
\Gr^\bF_t\Gr^{\Wr}_s\bE=
\bigoplus_{\substack{i-l=s-p\\ i+j=t\\ i+j+k+l=p}}\wedge^iA\otimes\wedge^jN\otimes\wedge^kB\otimes\wedge^lI.
$$
The intersection $Z_r \cap K_\bR$ of $K_\bR$ with the centralizer $Z_r$ of $y$ is the group
$$\{(a,b,c)\in\rU(r)\times\rU(p-r)\times\rU(p-r)\mid \det(a^2bc)=1\}.$$
As $(Z_r \cap K_\bR)$-modules, we have:
\begin{align*}
\ov{\wedge^iA\otimes\wedge^jN\otimes\wedge^kB\otimes\wedge^lI}
& \simeq \wedge^iA^\star\otimes\wedge^jN^\star\otimes\wedge^kB^\star\otimes\wedge^lI^\star \\
& \simeq (\wedge^iA^\star\otimes\wedge^jN^\star\otimes\wedge^kB^\star\otimes\wedge^lI^\star)
\otimes\wedge^{2p}(A\oplus N\oplus B\oplus I) \\
& \simeq \wedge^{r-i}A\otimes\wedge^{p-r-j}N\otimes\wedge^{p-r-k}B\otimes\wedge^{r-l}I \\
& \simeq \wedge^{r-l}A\otimes\wedge^{p-r-j}N\otimes\wedge^{p-r-k}B\otimes\wedge^{r-i}I
\\
& \simeq \wedge^{i'}A\otimes\wedge^{j'}N\otimes\wedge^{k'}B\otimes\wedge^{l'}I
\end{align*}
with $i'=r-l$, $j'=p-r-j$, $k'=p-r-k$ and $l'=r-i$
(on the fourth line we used the isomorphism $I \simeq A$).
Therefore, if $i-l=s-p$ and $i+j=t$, then $i'-l'=i-l=s-p$ and $i'+j'=p-j-l=s-(i+j)=s-t$. Hence the real
structure on $\bE$ defined as above with these isomorphisms is such that
$$
\ov{\Gr^\bF_t\Gr^{\Wr}_s\bE}=\Gr^\bF_{s-t}\Gr^{\Wr}_s\bE
$$
and the Hodge filtration induces a real Hodge structure of weight $s$ on $\Gr^{\Wr}_s\bE$, so
that the two filtrations $\Wr_\sbullet$ and $\bF^\scbullet$ indeed define a real mixed Hodge
structure on $\bE$.
\end{exem}

However, again from the Hodge structures point of view, our purpose here is a bit
different. We would like to construct a canonical ``weak sub-$\bR$-VHS'' of weight~$\rgy$ of the
restriction of
the $\bC$-VHS $\pi_M^\star E_{CY}$ to ${\Orb}_\rgy$. This is the role played by the submodule~$\Vr$. The meaning of ``weak $\bR$-VHS'' is that if $V^\rgy$ denotes the equivariant holomorphic
bundle $\Vr\otimes\cO_{\Orb_\rgy}$ on $\Orb_\rgy$ with its Gauss-Manin connection $\nabla$, and if
$F^\scbullet$ denotes the Hodge filtration on $V^r$ defined by the filtration $\bF^\scbullet$ of $\bE$, then
$(V^r,F^\scbullet)$ does have the symmetries of an $\bR$-VHS by Proposition~\ref{prop:V-self-dual},
but does not satisfy
Griffiths transversality $\nabla(F^t)\subset F^{t-1}\otimes\Omega^1_{\Orb_\rgy}$. Rather,
$(V^\rgy,F^\scbullet)$ satisfies a weak transversality condition by
Proposition~\ref{prop:algHiggs}. Indeed we have
$\nabla(F^t\otimes\cO_\bPTM(-1)_{|\Orb_\rgy})\subset F^{t-1}$, where $\cO_\bPTM(-1)$ is the
tautological line bundle on $\bPTM$. This will allow us to construct \emph{leafwise}
Higgs subsheaves (but not plain Higgs subsheaves) in Section~\ref{subsec:HiggsSubsheaf}.
\end{exem}

\section{Higgs bundles on foliated Kähler manifolds}
\label{sec:HiggsFoliated}

\subsection{Harmonic Higgs bundles}
\label{subsec:Higgs}

We keep the notation of the previous sections. In particular, $G_\bR$ is a simple real algebraic
Hermitian group, $K_\bR$ its maximal compact subgroup, $M=G_\bR/K_\bR$ the associated irreducible
Hermitian symmetric space of the noncompact type, $G$ and $K$ are the complexifications of $G_\bR$
and $K_\bR$, and $\fgg=\fk\oplus\fm$ and $\fgg_\bR=\fk_\bR\oplus\fm_\bR$ are the
associated Cartan decompositions of the Lie algebras of $G$ and $G_\bR$.

Let now $(Y,\omega_Y)$ be a compact Kähler manifold, $\Gamma=\pi_1(Y)$ its fundamental group, and $\rho:\Gamma\to
G_\bR$ be a
\emph{reductive} representation (group homomorphism) of $\Gamma$ into $G_\bR$. The assumption that $\rho$
is reductive
means that the real Zariski closure of $\rho(\Gamma)$ in $G_\bR$ is a reductive group.

In this case, by \cite{CorletteFlatGBundles},
there exists a $\rho$-equivariant \emph{harmonic} map $f$ from the universal cover $\tilde Y$ of $Y$ to the
symmetric space $M=G_\bR/K_\bR$ associated with $G_\bR$.
The fact that $Y$ is Kähler and the nonpositivity of the complexified sectional curvature of~$M$
imply by a Bochner formula due to~\cites{Sampson,Siu} that the map $f$ is \emph{pluriharmonic} (\ie its
restriction to 1-dimensional complex submanifolds of $Y$ is still harmonic), and that the image of its
(complexified) differential at every point $y\in Y$ is an abelian subalgebra of $T^\bC_{f(y)} M$
identified with $\fm$.

By the work of Hitchin and Simpson, this
gives a \emph{harmonic $G_\bR$-Higgs principal bundle} $(P_K,\theta)$ on $Y$. We will now briefly describe the
construction and the properties of such a Higgs bundle. Details and proofs can be found in the
original papers~\cites{HitchinSelfDual, HitchinLieTeich, S1, S2}.

There is a flat principal
bundle $P_{G_\bR}\to Y$ of group $G_\bR$ associated with the representation $\rho$. The $\rho$-equivariant map
$f:\tilde Y\to G_\bR/K_\bR$ defines a reduction of its structure group to $K_\bR$, \ie a principal
$K_\bR$ bundle $P_{K_\bR}\subset P_{G_\bR}$. The differential of $f$ can be seen as a 1-form with
values in $P_{K_\bR}(\fm_\bR):=(P_{K_\bR}\times\fm_\bR)/K_\bR$, the vector bundle associated with the adjoint action of $K_\bR$ on
$\fm_\bR$.

If we enlarge the structure group of $P_{K_\bR}$ to $K$,
the pluriharmonicity of $f$ implies that the $K$-principal bundle $P_K\to Y$ is a holomorphic
bundle and that the (1,0)-part $d^{1,0}f:T^{1,0}Y\to T^\bC M$ of the complexified differential
of $f$ defines a holomorphic section $\theta$ of
$P_K(\fm)\otimes\Omega^1_Y$, where $P_K(\fm)$ is the holomorphic vector bundle associated with the
principal bundle $P_K$ and the adjoint representation of $K$ on $\fm$. The section $\theta$ is called
the \emph{Higgs field} and satisfies the integrability
condition $[\theta,\theta]=0$ as a section of $P_K(\fm)\otimes\Omega^2_Y$.
The pair $(P_K,\theta)$ is called a \emph{$G_\bR$-Higgs principal bundle} on $Y$, see
\eg \citestar{BGPG2}*{\textsection 2.2}.

If now $\bE$ is a (complex) representation of $G$ we can construct the associated holomorphic vector bundle
$E:=P_K(\bE)$ over $Y$. Since $\bE$ is a representation of $\fgg$ and not only of $\fk$,
we have a morphism $P_K(\fm) \to P_K(\End(\bE)) = \End(E)$. Thus,
the Higgs field
$\theta$ can be seen as a holomorphic 1-form with values in the endomorphisms of $E$, \ie a
section of $\End(E)\otimes\Omega^1_Y$. The pair $(E,\theta)$
is called a \emph{$G_\bR$-Higgs vector bundle} on $Y$. The harmonic map $f$, seen as a reduction of the
structure group of $P_K$ to the compact subgroup $K_\bR$, together with a $K_\bR$-invariant metric on
$\bE$, gives a Hermitian metric on $(E,\theta)$
called the \emph{harmonic metric}.

The existence of this harmonic metric and the fact that $P_K$ comes from a flat principal $G_\bR$ bundle imply
that for any representation $\bE$ of $G$, the associated Higgs vector bundle $(E,\theta)\to Y$ is Higgs
polystable of degree 0, see~\cite{S1}. To explain
what Higgs polystability means, we first define Higgs subsheaves of the Higgs bundle $(E,\theta)$.
A coherent subsheaf $\cF$ of $\cO_Y(E)$ is a \emph{Higgs subsheaf} if it is invariant by the Higgs
field, \ie if it satisfies
$\theta(\cF)\subset\cF\otimes\Omega^1_Y$. The Higgs vector bundle $(E,\theta)$ is said to be \emph{Higgs stable} if
for any Higgs subsheaf $\cF$ of $(E,\theta)$ such that $0<\rk\cF<\rk E$, we have
$\mu_{\omega_Y}(\cF)<\mu_{\omega_Y}(E)$, where $\mu_{\omega_Y}(\cF)$ is
the slope of $\cF$, \ie its degree (computed w.r.t.\ the Kähler form $\omega_Y$ of $Y$) divided by its rank. The
Higgs bundle $(E,\theta)$ is \emph{Higgs polystable} if it is a direct sum of Higgs stable Higgs vector bundles
of the same slope.
Note that here $E$ is flat as a $C^\infty$ bundle, so that its degree is zero.

\begin{rema}\label{rem:real}
Since moreover we assumed that $G_\bR$ is a Hermitian group, then as a $K$-representation we have
$\fm=\fm^+\oplus\fm^-$
and the Higgs field $\theta$ on the principal bundle~$P_K$ (or on any associated vector bundle $E$) has two
components \hbox{$\beta\in P_K(\fm^-)\otimes\Omega^1_Y$} and $\gamma\in P_K(\fm^+)\otimes\Omega^1_Y$. The
vanishing of $\gamma$, \resp $\beta$, means that the harmonic map $f$ is holomorphic, \resp antiholomorphic. The component $\beta$, \resp $\gamma$, will be called the holomorphic, \resp antiholomorphic, component of the Higgs field $\theta$.
\end{rema}

\subsection{Harmonic Higgs bundles on foliated Kähler manifolds}\label{subsec:HiggsFol}

Assume now that the base Kähler manifold $Y$ of the harmonic $G_\bR$-Higgs vector bundle
$(E,\theta)\to Y$ of degree $0$ admits a smooth holomorphic
foliation by complex curves and that this foliation $\cT$ admits a transverse volume form
$\Omega$. This means that $\Omega$ is a nowhere vanishing semipositive closed $(d-1,d-1)$-form ($d$ is the
dimension of $Y$) such that the interior product of $\Omega$ with any vector tangent to the
foliation $\cT$ is zero.

Our goal in this section is to briefly explain the interplay between the Higgs bundle and the
foliation (equipped with its transverse volume form).
Details can be found in~\citestar{KM}*{\textsection 2.2}.

We first weaken the notion of Higgs subsheaves of $(E,\theta)$ to leafwise Higgs subsheaves by
asking only an
invariance by the Higgs field along the leaves. More \hbox{precisely} we now consider the
Higgs field as a section of $\Hom(E\otimes L,E)$,
where $L$ is the holomorphic
line subbundle of $T_Y$ tangent of the foliation $\cT$.
A \emph{leafwise Higgs subsheaf}~$\cF$ of
$(E,\theta)$ is a subsheaf of $\cO_Y(E)$ such that $\theta(\cF\otimes L)\subset \cF$.

We define the \emph{foliated degree} $\deg_\cT\cF$ and the {\em
foliated slope $\mu_\cT(\cF)$} of a coherent sheaf
$\cF$ on $Y$ by
wedging the first Chern class of $\cF$ with the transverse volume form~$\Omega$:
$$
\deg_\cT\cF:=\int_Y c_1(\cF)\wedge\Omega, \quad\mbox{and}\quad \mu_\cT(\cF)=\frac{\deg_\cT\cF}{\rk \cF}.
$$

The link between the notions of slope (and stability) for modules and for sheaves used in
this paper will be given by Fact~\ref{fact:slopes} below.

In order to state the needed leafwise polystability property of the Higgs bundle $(E,\theta)$ w.r.t.\ the
foliated degree, we first recall a few definitions.

An open subset of $Y$ is called \emph{big} if it is the complement of an analytic set of
codimension at least $2$ in $Y$. Two line bundles on $Y$ which have isomorphic restrictions
to a big open subset of $Y$ are in fact isomorphic over $Y$.

Let $\cF$ be a subsheaf of $\cO_Y(E)$. The complement of the biggest subset of $Y$ where $\cF$ is the
sheaf of sections of a subbundle $F$ of $E$ is called the \emph{singular locus} $\cS(\cF)$ of
$\cF$. Equivalently, $\cS(\cF)$ is the analytic subset of $Y$ where the quotient sheaf $\cO_Y(E)/\cF$ is not locally
free. Recall that a subsheaf $\cF$ of $\cO_Y(E)$ is \emph{saturated} if $\cO_Y(E)/\cF$ is torsion
free. The \emph{saturation} of $\cF$ is the kernel of the morphism $\cO_Y(E) \to
(\cO_Y(E)/\cF)_\rtf$, where
$(\cO_Y(E)/\cF)_\rtf$ denotes the quotient of $\cO_Y(E)/\cF$ by its torsion. If $\cF$ is
saturated, $\cS(\cF)$ has codimension at least $2$ in $Y$, so that the \emph{regular
locus} $Y\backslash\cS(\cF)$ of $\cF$ is a big open subset. Moreover a saturated subsheaf $\cF$ is
normal, so that in particular for all big open subsets $\cU$ of $Y$, the restriction
$\cF(Y)\to\cF(\cU)$ is injective.

A subset of the foliated manifold $(Y,\cT)$ is said to be \emph{saturated under the foliation~$\cT$} if it is a
union of leaves of $\cT$. Obviously if $Z\subset Y$ is saturated under $\cT$, then so is its
complement $ Y\backslash Z$.
(Unfortunately, the word ``saturated'' is used here with two different
meanings but this shouldn't cause any confusion.)

Here is the result (\citestar{KM}*{Prop.\,2.2}).

\begin{prop}\label{prop:polystab}
The harmonic Higgs bundle $(E,\theta)$ on the
foliated Kähler manifold $(Y,\cT,\Omega)$ is \emph{weakly polystable along the leaves} in the following sense:
\begin{enumerate}
\item\label{prop:polystab1} it is \emph{semistable along the leaves} of $\cT$: if $\cF$ is a leafwise Higgs
subsheaf of $(E,\theta)$, then $\deg_\cT\cF\leq 0$.
\item\label{prop:polystab2} if $\cF$ is a saturated leafwise Higgs subsheaf of $(E,\theta)$ such that $\deg_\cT\cF=0$, then the
singular locus $\cS(\cF)$ of $\cF$ is saturated under the foliation $\cT$. Moreover, on $Y\backslash
\cS(\cF)$, and if $F$ denotes the holomorphic subbundle of $E$ such that $\cF=\cO_{Y\backslash\cS(\cF)}(F)$ and $F^\perp$ its $C^\infty$
orthogonal complement w.r.t.\ the harmonic metric, then $\theta(F^\perp\otimes L)\subset F^\perp$. Finally,
for each leaf $\Leaf$ of $\cT$ such that $\Leaf\cap\cS(\cF)=\emptyset$, $F^\perp_{|\Leaf}$ is
holomorphic on $\Leaf$
and $(E,\theta)_{|\Leaf}=(F,\theta_{|F})_{|\Leaf}\oplus(F^\perp,\theta_{|F^\perp})_{|\Leaf}$ is a holomorphic
orthogonal decomposition of Higgs bundles on $\Leaf$.
\end{enumerate}
\end{prop}

\subsection{The tautological foliation on the projectivized tangent bundle of a complex hyperbolic
manifold}\label{subsec:toto}

An $n$-dimensional complex hyperbolic manifold $X$ is a quotient of the complex hyperbolic $n$-space
$\HC=\SU(1,n)/\rU(n)$ by a discrete torsion free subgroup $\Gamma$ of $\SU(1,n)$. It is a Hermitian
locally symmetric space of rank 1.
The complex hyperbolic space $\HC$ can be realized as the subset of
negative lines in $\bC^{n+1}$ for a Hermitian form $h$ of signature $(n,1)$
(recall our convention in Example~\ref{exem:su}), which is an open subset in
the projective space $\bCP^n$.\enlargethispage{.5\baselineskip}%

Intersections of lines
of $\bCP^n$ with $\HC$ are totally geodesic complex subspaces of~$\HC$ isometric to the Poincaré
disc. They are called \emph{complex geodesics}. The space $\cG$ of complex geodesics is the
homogeneous space $\SU(1,n)/\rS(\rU(1,1)\times \rU(n-1))$. It is also a complex manifold, since it
can be realized as the open subset of the Grassmannian manifold $\Gr_2(\bC^{n+1})$ of
$2$-planes in $\bC^{n+1}$ consisting of planes on which the signature of $h$ is $(1,1)$.

The \emph{projectivized tangent bundle} $\PTH$ of $\HC$ is the space of lines in the holomorphic
tangent bundle $\THC$ of $\HC$. As a homogeneous space it can be identified with
$\SU(1,n)/\rS(\rU(1)\times\rU(1)\times \rU(n-1))$. Again, its also a complex manifold (in fact a Kähler one)
since it identifies with an open subset of the manifold $\rF_{1,2}(\bC^{n+1})$
of incomplete flags $(\ell\subset\Pi\subset\bC^{n+1})$ with
$\dim \ell=1$, $ \dim\Pi=2$. A flag $(\ell\subset\Pi\subset\bC^{n+1})$ belong to $\PTH$ if the line
$\ell$ is negative and the plane $\Pi$ has signature $(1,1)$ for the Hermitian form $h$.

The natural $\SU(1,n)$-equivariant fibration $\pi_{\cG}:\PTH\to\cG$ which associates
to a tangent line to $\HC$ the complex geodesic it defines is a
disc bundle over $\cG$.

By $\SU(1,n)$-equivariance, this fibration endows the projectivized tangent bundle
$\PTX=\Gamma\backslash\PTH$ of $X=\Gamma\backslash\HC$ with a smooth holomorphic foliation $\cT$ by
complex curves, whose
leaves are given by the tangent spaces of the (immersed) complex geodesics in $X$. This foliation is
called the \emph{tautological foliation} of $\PTX$ because the tangent line bundle $L$ to the leaves is
naturally isomorphic to the tautological line bundle $\OX$ on $\PTX$.

The space $\cG$ of complex geodesics of $\HC$ is a pseudo-Kähler manifold: it admits an
$\SU(1,n)$-invariant Kähler form $\omega_{\cG}$, which is closed, of type $(1,1)$, non degenerate,
but not positive definite. The form $\omega_{\cG}$ is moreover unique up to scaling. This form
defines a transverse volume form $\Omega_{\cG}$
for the tautological foliation $\cT$ on $\PTX$, and the associated notion of foliated degree
$\deg_\cT$ for sheaves on $\PTX$ has the following fundamental property (after a suitable
normalization of the involved forms)
\citestar{KM}*{Prop.\,3.1}.

\begin{prop}\label{prop:degT}
Assume that $X=\Gamma\backslash\HC$ is compact and let $\pi:\PTX\to X$ be the projectivized tangent
bundle of $X$. If
$\cF$ is a coherent $\cO_X$-sheaf, then $\deg_\cT(\pi^*\cF)=\deg_\omega\cF$, where $\deg_\omega\cF$ is the usual
degree of $\cF$ computed w.r.t.\ the Kähler form $\omega$ on $X$.
\end{prop}

\vspace*{-\baselineskip}\vskip0pt
\section{The Milnor-Wood inequality}
\label{sec:MW}

Let $\Gamma$ be a uniform lattice in $\SU(1,n)$, $X$ be the compact quotient $\Gamma\backslash\HC$, and let~$\rho$
be a representation of $\Gamma$ in a simple real algebraic Hermitian Lie group $G_\bR$, whose associated
symmetric space is denoted by $M$. In this section we use
the material developed or recalled in Sections~\ref{sec:subnil} and~\ref{sec:HiggsFoliated}
to prove the Milnor-Wood inequality
$$
|\tau(\rho)|\leq \rkm\vol(X).
$$

Before going into the proof, we remark that if the lattice $\Gamma$ has torsion, the measure defined by
the Kähler form $\omega$ on $\HC$ descends to $X$ so that $\vol(X)$ is well-defined.
There are several ways to define the Toledo invariant of the representation $\rho$ of $\Gamma$ in
this case. Since this is how we will handle lattices with torsion, we observe that by
Selberg's Lemma (see \citestar{ratcliffe}*{p.\,327} or \cite{selberg}) there exists a finite index
torsion free subgroup~$\Gamma'$ of~$\Gamma$. We then define the Toledo invariant of $\rho:\Gamma\to
G_\bR$ by \hbox{$\tau(\rho):= \sfrac{\tau(\rho')}{|\Gamma/\Gamma'|}$}, where $\rho'=\rho_{|\Gamma'}$ is the restriction of $\rho$ to
$\Gamma'$ and $\tau(\rho')$ is defined as in the introduction. One checks easily that this does not depend
on the choice of~$\Gamma'$. Moreover, if
$X'=\Gamma'\backslash\HC$, then $\vol(X)=\sfrac{\vol(X')}{|\Gamma/\Gamma'|}$ so that proving the Milnor-Wood
inequality for $\rho:\Gamma\to G_\bR$ is equivalent to proving the Milnor-Wood inequality for $\rho':\Gamma'\to
G_\bR$.

In addition, the representation $\rho$ can always be deformed to a
reductive representation with the same Toledo invariant, see \eg \citestar{KM}*{Lem.\,4.11}. Moreover,
we may change the complex structure of $M$ to its opposite (\ie replace $z\in\ft$ by $-z\in\ft$, or
equivalently $\fm^\pm$ by $\fm^{\mp}$), thereby changing the sign of the Toledo
invariant of $\rho$.

Therefore we may assume without loss of generality that
\begin{itemize}
\item the lattice $\Gamma$ is torsion free, so that $X=\Gamma\backslash\HC$ is a compact
complex hyperbolic manifold of (complex) dimension $n$;
\item the representation $\rho$ is reductive, so that the results discussed in
Section~\ref{sec:HiggsFoliated} are available;
\item the Toledo invariant $\tau(\rho)$ is positive, so we only have to prove
\hbox{$\tau(\rho)\!\leq\!\rkm\vol(X)$}. (Observe that this implies that the harmonic map $f$ is not
antiholomorphic.)
\end{itemize}
We also assume as in Section~\ref{sec:subnil} that
\begin{itemize}
\item the complexification $G$ of $G_\bR$ is simply connected.
\end{itemize}
See Remarks~\ref{rema:1-connected} and~\ref{rema:1-connected-bis} where we explain why this is also
not a loss of generality.

\subsection{Setup}\label{subsec:setup}

Consider the representation $\bE$ of $G$ defined in Section~\ref{subsec:minuscule-rep}. As \hbox{explained}
in Section~\ref{subsec:Higgs}, this gives rise to a flat
harmonic $G_\bR$-Higgs vector bundle $(\ovE,\ovt)$ over $X$, associated to the $G_\bR$-Higgs principal
bundle $(\bar{P}_K,\ovt)$.
As a representation of $K$, $\bE=\bigoplus_{i=0}^\rkm\bE_{i}$, where $\rkm$ is the real rank of $G_\bR$. This
means that the holomorphic bundle $\ovE$ admits the holomorphic decomposition
$\ovE=\bigoplus_{i=0}^\rkm \ovE_{i}$, which is orthogonal for the harmonic metric. Moreover,
the components $\ovb\in P_K(\fm^-){\otimes\Omega^1_X}$ and $\ovg\in P_K(\fm^+){\otimes\Omega^1_X}$ of
the Higgs field $\ovt\in \Hom(\ovE,\ovE)\otimes\Omega^1_X$ (see Remark~\ref{rem:real}) satisfy
$$
\ovb\in\bigoplus_{i=0}^{\rkm-1}\Hom(\ovE_{i},\ovE_{i+1})\otimes \Omega^1_X\quad\mbox{and}\quad
\ovg\in\bigoplus_{i=0}^{\rkm-1}\Hom(\ovE_{i+1},\ovE_{i})\otimes \Omega^1_X.
$$

We pull-back the harmonic Higgs bundle $(\ovE,\ovt)\to X$ to the projectivized tangent bundle
$\PTX$ of $X$ to obtain a harmonic Higgs bundle that
we denote by $(E,\theta)\to\PTX$. We restrict the Higgs field $\theta$ to the tangent space $L$ of the tautological
foliation on
$\PTX$, so that its components $\beta$ and $\gamma$ satisfy
$$
\beta\in\bigoplus_{i=0}^{\rkm-1}\Hom(E_{i}\otimes L, E_{i+1})\quad\mbox{and}\quad
\gamma\in\bigoplus_{i=0}^{\rkm-1}\Hom(E_{i+1}\otimes L, E_{i}).
$$
We also call $P_K$ the pull-back of the principal bundle $\bar P_K$ to $\PTX$.

\begin{nota}
A tangent vector $\xi\neq 0$ to $X$ at a point $x\in X$ defines a point~$[\xi]$ in the
projectivized tangent bundle $\PTX$. It also defines an element
in the fiber ${\OX}_{[\xi]}=\{v\in T_{X,x}\mid v\in[\xi]\}$ of the tautological line bundle $\OX$
at~$[\xi]$, hence an element in the fiber $L_{[\xi]}$ of the tangent line bundle $L$ to the
tautological foliation $\cT$ at
$[\xi]$, which will also be denoted by $\xi$.
\end{nota}

\begin{defi}
\label{defi:rg-beta}
For $\xi\neq 0$ a vector tangent to $X$ and $[\xi]$ the corresponding point in $\PTX$,
the \emph{rank $\rk\beta_{[\xi]}$ of $\beta_{[\xi]}$} is the rank of the
corresponding element $\beta(\xi)$
in $\fm^-$, as defined in Definition \ref{defi:rank}.
By Proposition \ref{prop:rang-y}, it is also the largest value of $k$
such that $(\beta_{[\xi]})^k:E_0\otimes L^{k}\to E_{k}$ is not zero.

\begin{itemize}
\item
The \emph{generic rank $\rk\beta$
of $\beta$} is the maximum of the ranks of $\beta_{[\xi]}$ for $[\xi]\in\PTX$.

\item
The \emph{singular locus of $\beta$} is the following subset of $\PTX$:
\begin{align*}
\cS(\beta): &=\{[\xi]\in\PTX\mid
\rk\beta_{[\xi]}<\rk\beta\}\\
&=\{[\xi]\in\PTX\mid (\beta_{[\xi]})^{\rk\beta}: E_0\otimes L^{\rk\beta}\rightarrow E_{\rk\beta}\mbox{ vanishes}\}.
\end{align*}

\item
The \emph{regular locus of $\beta$} is $\cR(\beta):=\PTX\backslash\cS(\beta)$.

\item
The \emph{singular locus $\cS(\ovb)$ of $\ovb$} is the projection to $X$ of $\cS(\beta)$.

The \emph{regular locus $\cR(\ovb)$ of $\ovb$} is $X\backslash\cS(\ovb)$ (it is not the
projection to $X$ of ${\cR}(\beta)$).

\end{itemize}
\end{defi}

Our assumption that $\tau(\rho)\!>\!0$ implies that $\ovb\!\neq\!0$ and $\beta\!\neq\!0$, therefore
\hbox{$1\leq \rk\beta\leq\rkm$}.
Observe that while $\cS(\beta)$
is an analytic subset of $\PTX$ of codimension at least 1,
its projection $\cS(\ovb)$, although
it is an analytic subset of $X$ (because
$\pi:\PTX\to X$ is a proper map), might be equal to the whole base $X$. A very important
consequence of the
weak polystability along the
leaves (Proposition~\ref{prop:polystab} (2)), is that this does not happen when there is equality
in the refined Milnor-Wood inequality of Theorem~\ref{theo:MW}, see
Proposition~\ref{prop:equal}. This will indeed be a crucial point when dealing with maximal
representations.

\subsection{Rewording of the inequality}\label{subsec:reword}

Since the Hermitian symmetric space $M$ asso\-ciated with
$G_\bR$ is a Kähler-Einstein manifold, the first Chern
form $c_1(T_M)$ of its tangent bundle is a constant multiple of the $G_\bR$-invariant Kähler form $\omega_M$:
$c_1(T_M)=- \frac 1{4\pi}\,\cst \,\omega_M$ for some positive constant $\cst$ ($\cst=c_1(\check{M})$ is the first
Chern number of the compact dual $\check{M}$ of $M$). On the other hand, the line bundle $\cL$
associated with the $K$-representation $\bE_0$ is
a generator of the Picard group of the compact dual $\check{M}$ of $M$ and it can be checked that the
canonical bundle
$K_{\check{M}}$ of $\check{M}$ is precisely given by ${\cL}^{-\cst}$, see
\eg \citestar{KMgeneraltype}*{\textsection 2}.
Therefore the pull-back
$f^\star\omega_M$ is $4\pi$ times the first Chern form of the line bundle $f^\star\cL=\ovE_0$, so
that the Toledo invariant of $\rho$ is
\begin{equation}
\label{equa:toledo}
\tau(\rho)=4\pi\deg (\ovE_0)=4\pi\deg_\cT(E_0),
\end{equation}
where the last equality
follows from Proposition~\ref{prop:degT}. Similarly, we get that
\begin{equation}
\label{equa:deg-canonique}
\deg(K_X)=\frac{n+1}{4\pi}\vol (X).
\end{equation}

Recall that $L$ is the tangential line bundle to the tautological foliation $\cT$ on
the projectivized tangent bundle $\PTX$, and let $L^\star$ be its dual line bundle. One can compute
as explained in~\citestar{KM}*{\textsection 4.3.1} that
\begin{equation}
\label{equa:deg-L}
\deg_\cT(L^\star) = \frac{1}{2\pi} \vol(X).
\end{equation}

Therefore, the Milnor-Wood inequality can be rephrased as an inequality between the foliated degrees
of the line bundles $E_0$ and $L^\star$ on $\PTX$ and we need to prove:
\begin{equation}
\label{equa:rewording}
\deg_\cT(E_0) \leq \frac \rkm{2} \deg_\cT(L^\star),
\end{equation}
where $\rkm$ is the rank of the symmetric space $M$.

\subsection{Leafwise Higgs subsheaves associated with the holomorphic component of the Higgs field}
\label{subsec:HiggsSubsheaf}

We now define a subsheaf
$\cV$ of $\cE:=\cO(E)$ associated with~$\beta$ (\hbox{recall} that $\beta\neq 0$) in the same way we
defined the submodule $\Vr$ of $\bE$ associated with the nilpotent element $y\in\fm^-$ of rank $\rgy$
in Definition~\ref{defi:V} (for all $\xi\in L$, \hbox{$\beta(\xi)\in P_K(\fm^-)$} is a nilpotent endomorphism
of the bundle $E$). This subsheaf will be shown to be a leafwise Higgs subsheaf of the Higgs bundle
$(E,\theta)$ on $\PTX$. In Section~\ref{subsec:proofMW} this will be used to prove the Milnor-Wood
inequality.

More precisely, mimicking Definition~\ref{defi:V} and
for $-\rgb\leq k\leq \rgb$, we define the following subsheaves of $\cE$:

$$
\cW_k:=\sum_{\substack{\ell \geq 0 \\ k+\ell+1 \geq 0}}
\Ker \beta^{k+\ell+1} \cap \Im \beta^\ell,
$$
where in order to define $\Ker \beta^j$ we consider $\beta^j$ as a sheaf morphism from $\cE$
to $\cE\otimes {(L^\star)}^j$ and to define $\Im \beta^j$ we consider $\beta^j$ as a sheaf
morphism from $\cE\otimes L^j$ to $\cE$.

For $k=0,1,\ldots,\rgb$, let $\cV_k$ be the saturation in $\cE_k:=\cO(E_k)$ of the
subsheaf $\cE_k\cap \cW_{\rgb-2k}$ and set $$\cV=\bigoplus_{0\leq k\leq \rgb} \cV_k.$$

Since the sheaves $\cV_k$ are saturated subsheaves of $\cO(E_k)$, there exist a big open
subset $\cU$ of $\PTX$
and subbundles $V_k$ of $ E_k$ defined on $\cU$ such that the
restriction of the $\cV_k$'s to $\cU$ are the sheaves of
sections of the $V_k$'s. On $\cU$
we let $V$ be the subbundle $\oplus_{0\leq k\leq r} V_k$, so that $\cV_{|\cU}={\cO}_{\cU}(V)$.

Observe that on the regular locus
$\cR(\beta)$ of $\beta$, the rank of $\beta^k$, as a vector bundle morphism from $E\otimes
L^k$ to $E$, is constant for all $k$. Indeed, by Definition \ref{defi:rank}, the rank of
$\beta(\xi)$ viewed
as an element of $\fm^-$ determines its $K$-orbit.
Hence on this open subset the formulas used above to define the
subsheaves $\cW_k$ of $\cE$ in fact define subbundles~$W_k$ of $E$ such that ${\cW_k}_{|{\cR}(\beta)}=
\cO_{\cR(\beta)}(W_k)$. Therefore, on $\cR(\beta)$, the
subbundles $V_k$ such that ${\cV_k}_{|{\cR}(\beta)}=\cO(V_k)$ are given by $V_k=E_k\cap W_{\rgb-2k}$.
Thus, we may assume that $\cR(\beta)$ is contained (and therefore dense) in $\cU$.

We view an element $p$ of the $K$-principal bundle $P_K \To{\pi_K}\PTX$ above
$\xi\in\PTX$ as an isomorphism between the fiber $ E_\xi$ of $ E=P_K(\bE)$ and the model space $\bE$.
The component $\beta$ of the Higgs field is a section of $P_K(\fm^-)\otimes L^\star\subset
P_K(\End(\bE))\otimes L^\star$.

Let $y=y_\rgb$ and $Q_{\rgb}$ be defined as in Section~\ref{subsec:higgs-submodule}.

\begin{lemm}\label{lemm:reduction}
On the big open set $\cU$, the subsheaf $\cV$ defines a reduction $P_{K\cap Q_\rgb}$ of the
structure group of $P_K$ to the subgroup $K\cap Q_\rgb\subset K$.
\end{lemm}

\begin{proof}
We begin by working on $\cR(\beta)\subset\cU$.
Since for
all $\xi\in\cR(\beta)$ and all \hbox{$\eta\in L_\xi$}, $\eta\neq 0$, we have that $\beta_\xi(\eta)$ has rank
$\rgb$, there exists $p\in (P_K)_\xi$ such that \hbox{$p\circ\beta_\xi(\eta)\circ p^{-1}=y\in\fm^-\subset\End
(\bE)$}, so
that $p (V_\xi)=\Vb$. Therefore, on
$\cR(\beta)$, by Proposition~\ref{prop:stabV}\,\eqref{item:stabilizer-V},
the subbundle $V$ of $ E$ defines a (holomorphic) reduction $P_{K\cap Q_\rgb}$ of the
structure group of $P_K$ to the subgroup $K\cap Q_\rgb$ of $K$ ($Q_\rgb$ is the normalizer in $G$ of the
parabolic subalgebra $\fq_\rgb$, see Definition~\ref{defi:qlh}). Explicitly $P_{K\cap Q_\rgb}=\{p\in P_K\mid
p(V_{\pi_K(p)})=\Vb \}$.

We now work on $\cU$. Enlarge the structure group of $P_K$ to $\GL(\bE)$. The subbundle
$V=\bigoplus_{k=0}^\rgb V_k$ of $ E$ defines a reduction $P_S$ of the structure group of $P_{\GL(\bE)} \To{\pi_{\GL(\bE)}}\PTX$ to the
stabilizer $S$ of $\Vb$ in $\GL(\bE)$ by setting \hbox{$P_S\!=\!\{p\!\in\!P_{\GL(\bE)}\mid p (V_{\pi_{\GL(\bE)}(p)})\!=\!\Vb\}$}.

Let $B\subset\cU$ be an open ball on which $P_K$ is trivial. Then the reductions $P_{K\cap Q_\rgb}$ of
$P_K$ on $B\cap\cR(\beta)$
and $P_S$ of $P_{\GL(\bE)}$ on $B$ are respectively given by holomorphic maps $\sigma:B\cap\cR(\beta)\to
K/(K\cap Q_\rgb)$ and $s:B\to \GL(\bE)/S$. Moreover, if $\iota$ denotes the natural map $K/(K\cap Q_\rgb)\to
\GL(\bE)/S$, which is injective, then we have $s=\iota\circ\sigma$ on $B\cap\cR(\beta)$. Since
$K/(K\cap Q_\rgb)$ is compact, its image by $\iota$ is closed in $\GL(\bE)/S$. Therefore, since $B\cap\cR(\beta)$ is dense in $B$, $s$ maps $B$ to $\iota(K/(K\cap Q_\rgb))$. This means that the reduction
$P_{K\cap Q_\rgb}$ initially defined on $\cR(\beta)$ extends to $\cU$.
\end{proof}

We deduce the following result.

\begin{prop}\label{prop:VHiggs}
The subsheaf $\cV$ is a leafwise Higgs subsheaf of the Higgs bundle $(E,\theta)$ on $\PTX$.
\end{prop}

\begin{proof}
By Proposition~\ref{prop:algHiggs}, we know that $y$ and $\fm^+$ stabilize
$\Vb$. Therefore, on $\cR(\beta)$, the two components $\beta$ and $\gamma$ of the Higgs field
stabilize the subsheaf $\cV$ since it is the
sheaf of sections of the subbundle $V=P_{K\cap Q_\rgb}(\Vb)$
of $ E=P_{K\cap Q_\rgb}(\bE)$. By continuity, this still holds on $\cU$ since on this big open set
$\cV$ is also the sheaf of section of $V=P_{K\cap Q_\rgb}(\Vb)$.
Now, $\cV$ is a saturated, hence normal, subsheaf of
$\cO(E)$ by definition. Hence the restriction map $\cV(\PTX)\to \cV(\cU)$ is an isomorphism
since $\cU$ is big. Therefore $\cV$ is indeed a leafwise Higgs subsheaf of
$(E,\theta)$ on $\PTX$.
\end{proof}

\subsection{Proof of the Milnor-Wood inequality}\label{subsec:proofMW}

Recall that the slope of an $H$-module was introduced in Section \ref{subsec:slopeVi} and is
an element in $X(H) \otimes \bQ$. To relate it to the slope of the
associated vector bundle, we use the following fact which follows immediately from the definitions.

\begin{fact}
\label{fact:slopes}
Let $H$ be a complex reductive group and let $P_H \to Z$ a holomorphic $H$-principal bundle over a
complex manifold $Z$.
Let $\bV_i,i \in \{1,2\}$, be $H$-modules and let $V_i=P_H(\bV_i)$ be the associated vector bundles.
Assume that $\mu_H(\bV_1)=\mu_H(\bV_2)$. Then
$$ \frac{\det V_1}{\rk V_1} = \frac{\det V_2}{\rk V_2 }, $$
as elements of $\Pic(Z) \otimes_\bZ \bQ$. In other words, the line bundles
$(\det V_1)^{\otimes\rk V_2}$ and $(\det V_2)^{\otimes\rk V_1}$ are isomorphic.
\end{fact}

Therefore, if one is given a reasonable way to compute slopes of vector
bundles on
$Z$, \eg the usual one if $Z$ is $X$ with its Kähler form $\omega$,
or the foliated one if $Z$ is $\PTX$ with its tautological foliation $\cT$ and transverse
volume form
$\Omega_{\cG}$, then using the notation of Fact~\ref{fact:slopes}, $\mu_H(\bV_1)=\mu_H(\bV_2)$ implies
that the slopes of $V_1$ and $V_2$
are equal.

This observation, together with the computation of the slopes of the $H_\rgb$-submodules $\Vb_k$ of
$\bE$ and the construction of the leafwise Higgs subsheaf $\cV$ of $(E,\theta)$,
gives the Milnor Wood inequality.

\begin{theo}
\label{theo:MW} We have the inequalities
$$\deg_\cT(E_0)+\frac{\rk\beta}{2}\deg_\cT(L) \leq\mu_\cT(\cV)\leq 0.$$
Therefore the
Milnor-Wood inequality
$\deg_\cT(E_0)\leq \Psfrac \rkm{2}\deg_\cT(L^\star)$ holds.
\end{theo}

\begin{proof}
Let $n_k:=\dim \Vb_k$.
First, recall that Proposition~\ref{prop:stabV}\,\eqref{item:reductive-Q} states that the unipotent
radical of $Q_\rgb$ acts trivially on $\Vb$.
Hence so does the unipotent radical of $K \cap Q_\rgb$. Thus, in fact, $\Vb$ is a $(K \cap Q_\rgb)/R_u(K
\cap Q_\rgb)$-module, and
\hbox{$(K \cap Q_\rgb)/R_u(K \cap Q_\rgb) \simeq H_\rgb$} is reductive.

Therefore, by Lemma~\ref{lemm:reduction}, on the big open set $\cU$, the vector bundles $V_k$ such that
$\cV_k=\cO(V_k)$ are given by $V_k=P_{K\cap Q_\rgb}(\Vb_k)$.
By Proposition~\ref{prop:calcul-pente} and Fact \ref{fact:slopes} applied to
$\Vb_k$ and $\bV'_k:=\Vb_0 \otimes (\Vb_1 \otimes {\Vb_0}^\star)^{\otimes k}$, we have on $\cU$
$$
(\det V_k)^{\rk V'_k} \simeq (\det V'_k)^{\rk V_k},
$$
where $V'_k=P_{K\cap Q_\rgb}(\bV'_k)$. Letting $\cV'_k:=\cV_0 \otimes (\cV_1 \otimes
{\cV_0}^\star)^{\otimes k}$, we deduce that the line bundles $(\det \cV_k)^{\rk \cV_k'}$ and
$(\det \cV'_k)^{\rk \cV_k}$ on $\PTX$ are isomorphic on $\cU$, hence on $\PTX$, because $\cU$ is a
big open set of $\PTX$. Therefore
$$
\mu_\cT(\cV_k)=\mu_\cT(\cV'_k)=\deg_\cT(\cV_0)+k\,\mu_\cT(\cV_0^\star\otimes\cV_1).
$$

Since $\beta^\rgb:\cV_0\otimes L^\rgb\to\cV_\rgb$ is not zero and $n_0=n_\rgb=1$, we also have $\mu_\cT(\cV_\rgb)\geq \mu_\cT(\cV_0)+\rgb\,\mu_\cT(L)$ so that $\mu_\cT(\cV_0^\star\otimes\cV_1)\geq \deg_\cT (L)$.

Finally, remembering that $n_k=n_{\rgb-k}$ by Proposition~\ref{prop:self-dual} and that $\cV_0=E_0$, we~get
\begin{align*}
2\deg_\cT(\cV) & = \sum_{k=0}^\rgb \deg_\cT(\cV_k)+\sum_{k=0}^\rgb \deg_\cT(\cV_{\rgb-k}) \\
& = \sum_{k=0}^\rgb \left({ n_k\,\mu_\cT(\cV_k)+n_{\rgb-k}\, \mu_\cT(\cV_{\rgb-k})}\right) \\
& \geq \sum_{k=0}^\rgb n_k \left(\deg_\cT(\cV_0)+k\deg_\cT(L)+ \deg_\cT(\cV_0)+(\rgb-k)\deg_\cT(L)\right) \\
& = (\dim\Vb)\,(2\deg_\cT(E_0)+\rgb\deg_\cT(L)).
\end{align*}

The inequality $\mu_\cT(\cV)\leq 0$ follows from Propositions~\ref{prop:VHiggs}
and~\ref{prop:polystab},
and the
conclusion from the fact that $\rk\beta\leq \rkm$.
\end{proof}

In case of equality in Theorem~\ref{theo:MW}, we have

\begin{prop}\label{prop:equal}
Assume that $\deg_\cT(E_0)+ \frac{\rk\beta}{2}\deg_\cT(L) =0$. Then
\begin{enumerate}
\item on the regular locus $\cR(\beta)=\PTX\backslash\cS(\beta)$ of $\beta$, the orthogonal complement
$V^\perp=(\bigoplus V_i)^\perp$ of the subbundle $V=\bigoplus V_i$ of $E$ w.r.t.\ the harmonic metric is
stable under the Higgs field $\theta:E\otimes L\to E$;
\item the regular locus $\cR(\ovb)\subset X$ of $\ovb$ is (open and) dense in $X$.
\end{enumerate}
\end{prop}

\begin{proof}
The first point follows from the discussion after the definition of the subsheaves $\cV_k$ and the
polystability property \eqref{prop:polystab2} in Proposition~\ref{prop:polystab}, since our hypothesis implies that
$\deg_\cT \cV=0$ by Theorem~\ref{theo:MW}.

Proposition~\ref{prop:polystab}\,\eqref{prop:polystab2} also implies that the singular locus $\cS(\beta)$ of $\beta$, which is a
closed proper subset of $\PTX$, is saturated
under the tautological foliation $\cT$, see the proof of
\citestar{KM}*{Lem.\,4.5}. Since the leaves of the tautological foliation on $\PTX$ are
projections of $\SU(1,1)$-homogeneous subsets of $\Gamma\backslash\SU(1,n)$, M.\,Ratner's
results on unipotent flows and the fact that $\HC=\SU(1,n)/\rU(n)$ is a rank $1$ symmetric
space, then imply by~\citestar{KM}*{Prop.\,3.6} that
the singular locus $\cS(\ovb):=\pi(\cS(\beta))$ of $\ovb$ is a proper analytic subset of $X=\Gamma\backslash\HC$,
hence the second point of the proposition.
\end{proof}

\section{Maximal representations}\label{sec:max}

Maximal representations $\rho:\Gamma\to G_\bR$, where $\Gamma$ is a uniform lattice in $\SU(1,n)$ with
$n\geq 2$ and $G_\bR$ is a \emph{classical} Hermitian group,
were classified in \cite{KM}. Therefore we focus here on \emph{exceptional} targets, namely $G_\bR$ is either
$\mathrm{E}_{6(-14)}$, which is not of tube type, or $\mathrm{E}_{7(-25)}$, which is.

In Section~\ref{subsec:tube} we exclude the possibility of maximal representations in $\mathrm{E}_{7(-25)}$.
In fact, our uniform approach allows to easily prove that maximal
representations in tube type target groups $G_\bR$ do not exist (Proposition~\ref{mainprop}). The
case of $\mathrm{E}_{6(-14)}$ is
treated in Section~\ref{subsec:E6}, Theorem~\ref{maintheo} and Corollary~\ref{maincoro} are
established in Section~\ref{subsec:proof}.

\subsection{Tube type targets}\label{subsec:tube}

We prove Proposition~\ref{mainprop}, namely that whenever $G_\bR$ has tube type and $n\geq 2$,
the Toledo invariant of a representation
from $\Gamma$ to
$G_\bR$ satisfies an inequality stronger than the Milnor-Wood inequality, preventing a
representation in such a group to be maximal.

\begin{prop}
\label{prop:tube}
Let $\Gamma$ be a uniform lattice in $\SU(1,n)$, and let $X=\Gamma\backslash\HC$.
Assume that the simple real algebraic Hermitian Lie group $G_\bR$ has tube type and let $\rkm$ be the real rank of $G_\bR$. Let $\rho$ be a
representation $\Gamma\to G_\bR$. Then
$$
|\tau(\rho)|\,\leq\, \max\Bigl\{\rkm-1,\frac \rkm{2}\cdot \frac{n+1}{n}\Bigr\}\vol(X). $$
\end{prop}

\begin{rema}
Observe that $\max\left\{{\rkm-1,\frac \rkm{2}\cdot \frac{n+1}{n}}\right\}=\rkm-1$ unless
$\rkm\leq 2$ or $\rkm=3$ and $n=2$.
\end{rema}

\begin{proof}
We may assume that $\Gamma$ is torsion free, that $\tau(\rho)>0$, that $\rho$
is reductive (see the beginning of Section~\ref{sec:MW}), and that $G$ is simply connected (see
Remarks~\ref{rema:1-connected} and~\ref{rema:1-connected-bis}). Then, the constructions of
Sections~\ref{sec:subnil}, \ref{sec:HiggsFoliated} and~\ref{sec:MW}
are valid and the inequality of the proposition is equivalent to the inequality
$$
\deg_\cT(E_0) \,\leq\,
\max\Bigl({\frac{\rkm-1}{2},\frac {\rkm}{2} \cdot \frac{n+1}{2n}} \Bigr) \deg_\cT(L^\star).
$$
We use freely the notation of Section~\ref{sec:MW}. If the generic rank of $\beta$ on the
projectivized tangent bundle $\PTX$
of $X$ is $\leq \rkm-1$ then we are done by Theorem~\ref{theo:MW}. Therefore we may assume that the
generic rank of $\beta$ on $\PTX$ is $\rkm$.

We come back to the Higgs bundle $(\ovE,\ovt)$ on $X$ and we consider \hbox{$\ovb:\ovE\otimes T_X\to\ovE$}. The fact that $\rk\beta=\rkm$ implies that $\ovb^\rkm$, seen as a morphism from
$\ovE_0\otimes\nobreak\ovE_\rkm^\star$ to the $\rkm$-th symmetric power $S^\rkm\Omega^1_X$ of $\Omega^1_X$, is not
zero. Since $\Omega^1_X$ is a semi\-stable bundle
over $X$ ($X$ is Kähler-Einstein), so is $S^\rkm\Omega^1_X$. On the other hand,
$\ovE_0\otimes\ovE_\rkm^\star$ is also semistable
because it is a line bundle by Section~\ref{subsubsec:tube}.
Therefore \hbox{$\mu_\omega(\ovE_0\otimes\ovE_\rkm^\star)\leq \mu_\omega(S^\rkm\Omega^1_X)$}
(recall that $\omega$ is our Kähler form on $X$),
so that $\deg_\omega \ovE_0-\deg_\omega \ovE_\rkm\leq \rkm\,\mu_\omega(\Omega^1_X)$. Now, as explained in
Section~\ref{subsubsec:tube}, the $K$\nobreakdash-mod\-ules~$\bE_\rkm$ and $\bE_0^\star$ are isomorphic, so that $\deg_\omega \ovE_\rkm=-\deg_\omega \ovE_0$
by Fact~\ref{fact:slopes}.
We~get the result, since by Equations~\eqref{equa:deg-canonique} and~\eqref{equa:deg-L}
in Section \ref{subsec:reword}
and the isomorphism $K_X \simeq \det \Omega^1_X$, we have
$\deg_\omega(\Omega^1_X)=\frac{n+1}{2}\deg_\cT(L^\star)$.
\end{proof}

\Subsection{Target group $\mathrm{E}_{6(-14)}$}\label{subsec:E6}

\subsubsection{Algebraic preliminaries}

In the case $G_\bR=\mathrm{E}_{6(-14)}$, the cominuscule
representation of $G=\mathrm{E}_6$ is the defining
representation of ${\mathrm{E}_6}$ on the 27-dimensional complex exceptional Jordan algebra $\bE=\fJ_\bC$.
The real rank of $\mathrm{E}_{6(-14)}$ is $2$ and $\bE=\bE_0\oplus\bE_1\oplus\bE_2$ with $\bE_0$, $\bE_1$,
and $\bE_2$ of dimension 1, 16 and 10 respectively. The semi-simple part of $K$ is isomorphic to
$\Spin_{10}$ and as $\Spin_{10}$-representations, $\bE_0$, $\bE_1$, and $\bE_2$ are respectively a
trivial, half-spin, and vector representation.

Given $x \in \fm^-$ (\resp $y \in \fm^+$), $x$ (\resp $y$) defines linear maps $\bE_0 \to \bE_1$ and $\bE_1 \to \bE_2$ (\resp $\bE_1 \to \bE_0$ and $\bE_2 \to \bE_1$). We denote these maps by $\lambda_1(x),\lambda_2(x)$
\resp $\mu_1(y),\mu_2(y)$. We thus have maps
$\lambda_1(x):\bE_0 \to \bE_1$, $\lambda_2(x):\bE_1 \to \bE_2$ and
\hbox{$\mu_1(y):\bE_1 \to \bE_0$}, $\mu_2(y):\bE_2 \to \bE_1$.

We start with a description of the $\Spin_{10}$-representations $\bE_1$ and $\bE_2$ in terms of octonions.
More precisely, by Proposition~\ref{prop:Ei}\eqref{prop:Eid}, there is a $\Spin_{10}$-equivariant isomorphism
$\bE_1 \simeq \bE_0 \otimes \fm^-$.
Choosing a non-zero vector in $\bE_0$, this yields an isomorphism $\alpha: \bE_1 \to \fm^-$.
We consider the quadratic map $\kappa: \bE_1 \to \bE_2$ defined by
$\kappa(x)=\lambda_2(\alpha(x)) \cdot x$. This is a $\Spin_{10}$-equivariant quadratic map $\bE_1 \to \bE_2$.

As the next proposition shows, there is, up to a scale, only one such map, and this map
can be described in terms of octonions. We denote by $\bO$ the (non associative) algebra
of octonions (see \citestar{chevalley}*{\textsection 4.5} or \cite{springer}).
We denote by $N: \bO \to \bC$ the norm, a quadratic map such that
$N(z_1z_2)=N(z_1)N(z_2)$ for all octonions $z_1,z_2$.
The following is certainly well-known to specialists, however we could not find an adequate reference.

\begin{prop}
\label{prop:spin-10}
There are $(\Spin_8 \times \Spin_2)$-equivariant identifications
of $\bE_1$ with $\bO \oplus \bO$ and $\bE_2$ with $\bC \oplus \bO \oplus \bC$
such that $\kappa(u,v)=(N(u),uv,N(v))$.
\end{prop}
\begin{proof}
We consider the $\Spin_{10}$ half-spin representation $\bE_1$. According to \cite{chevalley}, when we
restrict to $\Spin_8$, this representation splits as $\cS^+ \oplus \cS^-$, where $\cS^\pm$ denote the two
half-spin representations of $\Spin_8$. Similarly, the $\Spin_{10}$ vector representation~$\bE_2$ splits as
$\bC \oplus \bV \oplus \bC$, where $\bV$ denotes the $8$-dimensional vector representation.

Now, the quadratic map $\kappa$ is given by a $\Spin_{10}$-equivariant linear morphism
\hbox{$S^2 \bE_1 \to \bE_2$} which is itself induced by an
injection
$$\bE_2 \subset \bE_1 \otimes \bE_1 =
\cS^+ \otimes \cS^+ \oplus \cS^+ \otimes \cS^- \oplus \cS^- \otimes \cS^+ \oplus \cS^- \otimes \cS^-.$$
Since there are $\Spin_8$-equivariant maps $\cS^+ \otimes \cS^- \to \bV$,
$\cS^+ \otimes \cS^+ \to \bC$ and $\cS^- \otimes \cS^- \to \bC$, and no $\Spin_8$-equivariant maps from other
factors in the tensor product $\bE_1 \otimes \bE_1$ to $\bE_2$, $\kappa$ is of the form
$\kappa(s_+,s_-) = (\psi_+(s_+), \varphi(s_+,s_-), \psi_-(s_-))$, for some equivariant
quadratic maps $\psi_\pm: \cS^\pm \to \bC$, and for a bilinear map $\varphi:\cS^+ \times \cS^- \to \bV$.
None of these maps can vanish, otherwise the image of $\kappa$
would be degenerate. Moreover, there are, up to scale, only one such map, as it follows from
the decomposition of the tensor product of two spin representations
\citestar{chevalley}*{\textsection 3.3}.
It is given,
once $\cS^+,\cS^-$ and $\bV$ are identified with the space of octonions $\bO$, by the formulas:
$\psi_+(s_+)=N(s_+),\varphi(s_+,s_-)=s_+s_-$ and $\psi_-(s_-)=N(s_-)$
(see again \citestar{chevalley}*{\textsection 4.5}). The proposition follows.
\end{proof}

We can deduce from the explicit formula above some information about maps~$\lambda_2(x)$.

\begin{prop}
\label{prop:fu-}
Let $x,y \in \bE_1 \simeq \fm^-$.
\begin{enumeratea}
\item\label{prop:fu-a} $x$ has rank one \iff $x \neq 0$ and $\kappa(x) = 0$.
\item\label{prop:fu-b} $x$ has rank one \iff $\dim \left (\im \lambda_2(x) \right) =5$.
\item\label{prop:fu-c} $x$ has rank two \iff $\dim \left (\im \lambda_2(x) \right) =9$ \iff $\kappa(x) \neq 0$.
\item\label{prop:fu-d} Assume that all non trivial linear combinations of $x$ and $y$ have rank $2$ (in~particular,
$x$ and $y$ are not colinear). Then
$\dim (\Ker \lambda_2(x) \cap \Ker \lambda_2(y)) \leq 3$.
\item\label{prop:fu-e} Assume that $x$ and $y$ have rank $1$ and $\dim (\im \lambda_2(x) \cap \im \lambda_2(y)) \geq 4$.
Then $x$ and $y$ are proportional.
\end{enumeratea}
\end{prop}
\begin{proof}
We use the above isomorphism $\bE_1 \simeq \bO \oplus \bO$. According to \citestar{igusa}*{Prop.\,2}, there are exactly
$3$ orbits in $\bE_1$ under $\Spin_{10} \times \bC^*$.
Let $u \in \bO$ such that $N(u)=0$. We have
$\kappa(u,0)=(0,0,0)$ and $\kappa(1,0)=(1,0,0)$. Thus, $(u,0)$ and $(1,0)$ cannot be in the same orbit.
It follows that $(u,0)$ has rank $1$ and $(1,0)$ has rank $2$ and statement~\eqref{prop:fu-a} of the proposition is
proved.

Let $\tilde \kappa: \bE_1 \times \bE_1 \to \bE_2$ be the polarization of $\kappa$, namely, the unique
symmetric bilinear map such that $\tilde \kappa (x, x) = \kappa (x)$ for all $x$ in $\bE_1$.
We have $\lambda_2(x)=\tilde \kappa (x, \cdot)$. Thus the image of $\lambda_2(u,0)$ is the set of triples
$(t,z,0)$ with $t \in \bC$ arbitrary and $z$ of the form $uw$ for some $w$ in $\bO$: this space
has dimension $5$. On the other hand, the image of $\lambda_2(1,0)$ is the set of triples
$(t,z,0)$ with $t$ and $z$ arbitrary. It has dimension $9$. Points~\eqref{prop:fu-b} and \eqref{prop:fu-c} are
proved.

For point \eqref{prop:fu-d}, let us assume that any non trivial linear combination of $x$ and $y$ has rank $2$. Thanks to the
result of Igusa, we may assume that $x=(1,0)$. Writing $y=(a,b)$, the assumption implies that $b \neq 0$
(in fact, if $y=(a,0)$, then some linear combination of $x$ and $y$ will be of the form $(u,0)$ with
$N(u)=0$). The kernel of $\lambda_2(x)$ is the space of elements of the form $(u,0)$ with $\scal{u,1}=0$
(here, $\scal{\cdot\,, \cdot}$ denotes the symmetric bilinear form associated with $N$).
If such an element is in the kernel of $\lambda_2(y)$ then $bu=0$ and so $N(u)=0$.
In fact, the octonion algebra is \emph{alternative}, meaning that for all elements $a,b$ one has
$(ab)b=a(b^2)$. Since $\ov b$ is a linear combination of $1$ and $b$, we also have
$(ab)\ov b=a(b \ov b)$.
From the equation $0=bu$ it thus follows $0=(bu)\ov u=b(u \ov u)=bN(u)$ so
$N(u)=0$.
Thus, the intersection of
the kernels of $\lambda_2(x)$ and $\lambda_2(y)$ is isomorphic to an isotropic subspace of the space of
octonions $u$ with $\scal{u,1}=0$.
Such an isotropic subspace can have at most dimension $3$.

Finally, let us assume that $x$ and $y$ have rank $1$ and that $$\dim (\im \lambda_2(x) \cap \im \lambda_2(y)) \geq 4.$$
Then we may assume that $x=(u,0)$ with $N(u)=0$ as
above. The image of $\lambda_2(x)$ is then the set of triples $(t,z,0)$ with $t$ arbitrary and $z$ of the
form $uw$ for some octonion~$w$. Thus, this space is an isotropic subspace of $\bE_2$ of maximal dimension $5$.
Using the $\Spin_{10}$-action, it follows that for any $x \in \bE_1$ of rank $1$, the image of $\lambda_2(x)$
is an isotropic subspace of dimension $5$. Since two maximal isotropic subspaces
in the same family can intersect only
in odd dimension, it follows from the hypothesis on $x$ and $y$
that the images of $\lambda_2(x)$ and $\lambda_2(y)$ are equal. Thus
$x$ and $y$ are pure spinors representing the same maximal isotropic subspace, so they are
proportional by \citestar{chevalley}*{III.1.4}. It is also possible to draw this conclusion
by a direct computation, using the fact
that for $u,v \in \bO$ such that $N(u)=N(v)=0$, the equality
\hbox{$\{uz\mid z \in \bO\} = \{vz\mid z \in \bO\}$} is equivalent to the fact that $u$ and $v$ are proportional.	
\end{proof}

Recall that we constructed a quadratic $\Spin_{10}$-equivariant map $\kappa: \bE_1 \to \bE_2$ by setting
$\kappa(x)=\lambda_2(\alpha(x)) \cdot x$, where $\alpha: \bE_1 \to \fm^-$ is a $\Spin_{10}$-equivariant isomorphism.
Similarly, let $\beta:\bE_1^\star \to \fm^+$ be a $\Spin_{10}$-equivariant isomorphism.
Let $\iota: \bE^\star_1 \to \bE_2^\star$ be the quadratic equivariant map obtained setting
$\iota(w)={}^t \mu_2(\beta(w)) \cdot w$.
With the same proof, we get information about $\fm^+$ and the linear maps $\mu_2(w)$.

\begin{prop} \label{prop:fuplus}
Let $w,z \in \bE_1^* \simeq \fm^+$.
\begin{enumeratea}
\item $w$ has rank one \iff $w \neq 0$ and $\iota(w) = 0$.
\item $w$ has rank one \iff $\dim (\im \mu_2(w))=5$.
\item $w$ has rank two \iff $\dim (\im \mu_2(w))=9$ \iff $\iota(w) \neq 0$.
\item Assume that all non trivial linear combinations of $w$ and $z$ have rank $2$ (in particular,
$w$ and $z$ are not colinear). Then
$\dim (\im \mu_2(w) \cap \im \mu_2(z)) \leq 3$.
\item Assume that $w$ and $z$ have rank $1$ and $\dim (\Ker \mu_2(w) \cap \Ker \mu_2(z)) \geq 4$.
Then~$w$ and $z$ are proportional.
\end{enumeratea}
\end{prop}

\subsubsection{Maximal representations}

Let $\Gamma$ be a torsion free uniform lattice in $\SU(1,n)$ and $\rho:\Gamma\to \mathrm{E}_{6(-14)}$ be a
reductive representation.
We may therefore
consider the the Higgs bundle $(\ovE,\ovt)$ on $X$ and its pull-back $(E,\theta)$ on $\PTX$
associated with $\rho$ and the representation of $\mathrm{E}_{6}$ on
$\bE=\bE_0\oplus\bE_1\oplus\bE_2$ as in Section~\ref{sec:MW}.

Recall that the components of the Higgs field $\ovt$ are
$$
P_K(\fm^-)\ni\ovb=:(\ovb_1,\ovb_2)\in\bigl(\Hom(\ovE_0,\ovE_1)\otimes\Omega^1_X \bigr)\oplus \bigl(\Hom(\ovE_1,\ovE_2)\otimes\Omega^1_X \bigr)
$$
and
$$
P_K(\fm^+)\ni\ovg=:(\ovg_1,\ovg_2)\in \bigl(\Hom(\ovE_1,\ovE_0)\otimes\Omega^1_X \bigr)\oplus \bigl(\Hom(\ovE_2,\ovE_1)\otimes\Omega^1_X \bigr).
$$

To lighten the notation, the fibers of the bundles $\ovE$, $\ovE_0$, $\ovE_1$ and
$\ovE_2$ above some $x\in X$ will also be denoted by $\ovE$, $\ovE_0$, $\ovE_1$ and
$\ovE_2$.

Propositions~\ref{prop:fu-} and~\ref{prop:fuplus} immediately imply the following.

\begin{coro}\label{lem:ranks} Let $x\in X$ and $\xi$, $\eta$ be tangent vectors at $x$.

As an element of $\Hom(\ovE_1,\ovE_2)$, $\ovb_2(\xi)$ has rank
$0$, $5$ or $9$. Moreover:
\begin{enumeratea}
\item If $\ovb(\xi)$ has rank 1, \ie $\ovb_1(\xi)\neq 0$ but $\ovb_2(\xi)\ovb_1(\xi)=0$, then $\ovb_2(\xi):\ovE_1\to \ovE_2$ has rank $5$;

\item If $\ovb(\xi)$ has rank $2$, \ie if $\ovb_2(\xi)\ovb_1(\xi)\neq 0$, then $\ovb_2(\xi):\ovE_1\to \ovE_2$ has rank $9$;

\item If all non trivial linear combinations of $\ovb(\xi)$ and $\ovb(\eta)$ have rank
$2$ (in particular, $\xi$ and $\eta$ are not colinear), then we have
$\dim(\Ker\ovb_2(\xi)\cap\Ker\ovb_2(\eta))\leq 3$.
\end{enumeratea}

Similarly, as an element of $\Hom(\ovE_2,\ovE_1)$, $\ovg_2(\xi)$ has rank
$0$, $5$ or $9$. Moreover:
\begin{enumerateap}
\item If $\ovg(\xi)$ has rank $1$, \ie $\ovg_2(\xi)\neq 0$ but $\ovg_1(\xi)\ovg_2(\xi)=0$, then $\ovg_2(\xi):\ovE_2\to \ovE_1$ has rank $5$;

\item If $\ovg(\xi)$ has rank $2$, \ie if $\ovg_1(\xi)\ovg_2(\xi)\neq 0$, then $\ovg_2(\xi):\ovE_2\to \ovE_1$ has rank $9$;

\item If $\ovg(\xi)$ and $\ovg(\eta)$ have rank $1$ and $\dim (\Ker \ovg_2(\xi) \cap \Ker \ovg_2(\eta)) \geq 4$, then
$\ovg(\xi)$ and $\ovg(\eta)$ are colinear.
\end{enumerateap}
\end{coro}

Thanks to this corollary, in case of equality in the Milnor-Wood inequality, we may prove

\begin{prop}\label{prop:vanishing} Let $(\bar E, (\bar\beta,\bar\gamma))$ be the harmonic Higgs bundle on $X$
associated with the reductive representation $\rho:\Gamma\to \mathrm{E}_{6(-14)}$ of the torsion free
uniform lattice
$\Gamma\subset\SU(1,n)$ and assume
$n\geq 2$.

If $\deg_\cT(E_0)=\deg_\cT(L^\star)$ and $x\in\cR(\ovb)$, then for all $\xi\in T_{X,x}$, $\ovg(\xi)=0$.
\end{prop}

\begin{proof}
The
letters $\xi$ and $\eta$ will denote tangent
vectors at $x$.

The equality $\deg_\cT(E_0)=\deg_\cT(L^\star)$ and Theorem~\ref{theo:MW} imply that the generic rank
of $\beta$ on $\PTX$ is 2. Therefore, since
$x$ belongs to the regular locus $\cR(\ovb)$ of $\ovb$, for all $\xi\neq 0$ in $T_{X,x}$, the
rank of $\ovb(\xi)$ is 2. This implies by
Corollary~\ref{lem:ranks}~(b) that the rank of $\ovb_2(\xi)$ is 9 for all $\xi\neq 0$, and by
Corollary~\ref{lem:ranks}~(c) and linearity of $\bar\beta: T_{X,x}\to \Hom(\bar E,\bar E)$, that if $\xi$
and $\eta$ are not colinear, $\dim(\Ker\ovb_2(\xi)\cap\Ker\ovb_2(\eta))\leq 3$.

We will make a crucial use of the
integrability relation $[\ovt,\ovt]=0$ of the Higgs field~$\ovt$. This relation is equivalent to the
following three conditions:
$$
\begin{cases}
\ovg_1(\xi)\ovb_1(\eta) = \ovg_1(\eta)\ovb_1(\xi) & \mbox{in } \End(\ovE_0)\\
\ovb_1(\xi)\ovg_1(\eta)+\ovg_2(\xi)\ovb_2(\eta) = \ovb_1(\eta)\ovg_1(\xi)+\ovg_2(\eta)\ovb_2(\xi) & \mbox{in } \End(\ovE_1) \\
\ovb_2(\xi)\ovg_2(\eta) = \ovb_2(\eta)\ovg_2(\xi) & \mbox{in } \End(\ovE_2)
\end{cases}
$$
which hold for all $\xi,\eta$.

Suppose first that there exists $\xi$ such that $\ovg_2(\xi):\ovE_2\to \ovE_1$ has rank 9. Consider
the subspace
$W:=\Ker\ovg_1(\xi)\cap\Ker \ovb_2(\xi)\subset \ovE_1$. Since $\dim \ovE_1=16$, $\dim
\Ker\ovg_1(\xi)=15$ and
$\dim\Ker\ovb_2(\xi)=7$, we have $\dim W\geq 6$. On this subspace, the second integrability
condition
reads $\ovb_1(\xi)\ovg_1(\eta)+\ovg_2(\xi)\ovb_2(\eta) =0$ for all $\eta$. Therefore
\hbox{$\ovg_2(\xi)\ovb_2(\eta)(W)\subset \ovE_1$} is
$1$-dimensional. Because of our assumption on the rank of $\ovg_2(\xi)$, $\ovb_2(\eta)(W)$ is of
dimension at most 2, and this implies that $\dim W\cap\Ker \ovb_2(\eta)\geq 4$, hence that $\dim
\Ker\ovb_2(\xi)\cap\Ker\ovb_2(\eta)\geq 4$. Since $n\geq 2$, we may choose $\xi$ and $\eta$ to be
linearly independent and we get a contradiction with Corollary~\ref{lem:ranks}~(c).

Suppose now that for all $\xi\neq 0$, $\ovg_2(\xi)$ has rank 5. Fix $\xi\neq 0$, and let
$[\xi]$ be the class of $\xi$ in the fiber of $\PTX$ above $x$. Let $V(\xi)=V_0(\xi)\oplus
V_1(\xi)\oplus V_2(\xi)$ be the fiber above~$[\xi]$ of the subbundle $V$ of the Higgs bundle
$(E,\theta)$ on $\PTX$. We have
$$
\begin{cases}
V_0(\xi)=E_0=\ovE_0 \\
V_1(\xi)=E_1\cap F_0=E_1\cap(\Ker\beta_{[\xi]}\oplus\Ker\beta_{[\xi]}^2\cap\Im\beta_{[\xi]})=\Ker\ovb_2(\xi)\oplus\Im\ovb_1(\xi)\\
V_2(\xi)=E_2\cap F_{-2}= E_2\cap(\Ker\beta_{[\xi]}\cap\Im\beta^2_{[\xi]})=\Im\ovb_2(\xi)\ovb_1(\xi)
\end{cases}
$$
and we know by Proposition~\ref{prop:equal}~(1) that the orthogonal complement $V_1(\xi)^\perp\oplus
V_2(\xi)^\perp$ of $V_0(\xi)\oplus V_1(\xi)\oplus V_2(\xi)$ is invariant by $\ovt(\xi)$, in particular that
$\ovg_2(\xi)$ maps $V_2(\xi)^\perp$ to $V_1(\xi)^\perp$. Here and after, for $i\in\{1,2\}$, we
denote by $V_i(\xi)^\perp$ the
orthogonal complement of $V_i(\xi)$ in $E_i$ with respect to the harmonic metric.

By the third integrability condition, $\ovg_2(\xi)$ maps
$\Ker\ovg_2(\eta)$ in $\Ker\ovb_2(\eta)$. Hence $\ovg_2(\xi)$ maps
$V_2(\xi)^\perp\cap\Ker\ovg_2(\eta)$ to
$\Ker\ovb_2(\eta)\cap V_1(\xi)^\perp$.

But $\ovb_2(\xi)$ is injective on $V_1(\xi)^\perp$ because $\Ker\ovb_2(\xi)\subset
V_1(\xi)$. Hence for $\eta$ close to $\xi$, $\ovb_2(\eta)$ is also injective on $V_1(\xi)^\perp$, so that
$\Ker\ovb_2(\eta)\cap V_1(\xi)^\perp=0$ and hence
$V_2(\xi)^\perp\cap\Ker\ovg_2(\eta)\subset\Ker\ovg_2(\xi)$ by the previous paragraph. Now, \hbox{$\dim
V_2(\xi)^\perp=9$} and
$\rk\ovg_2(\eta)=5$, thus $V_2(\xi)^\perp\cap\Ker\ovg_2(\eta)$ is at least 4-dimensional, and so is \hbox{$\Ker\ovg_2(\xi)\cap\Ker\ovg_2(\eta)$}. This implies
by Corollary~\ref{lem:ranks}~($c'$) that $\ovg_2(\xi)$ and
$\ovg_2(\eta)$ are colinear, a contradiction since $n\geq 2$ and $\bar\gamma_2$ is injective by our
assumption that all the $\ovg_2(\zeta)$, $\zeta\neq 0$, have rank 5.

We conclude that there exists $\xi\neq 0$ such that $\ovg_2(\xi)=0$. Then also $\ovg_1(\xi)=0$ and by
the second integrability condition, for all $\eta$, $\ovb_1(\xi)\ovg_1(\eta) =
\ovg_2(\eta)\ovb_2(\xi)$. Therefore
$\ovg_2(\eta)$ has rank at most $1$ on $\Im\ovb_2(\xi)$ which is 9-dimensional in $\ovE_2$, so that
$\ovg_2(\eta)$ has
rank at most 2, hence vanishes. Therefore $\ovg_2=0$ and $\ovg=0$ identically on $T_{X,x}$.
\end{proof}

\begin{theo}\label{thm:E6}
Let $\Gamma$ be a torsion free uniform lattice in $\SU(1,n)$ with $n\geq 2$ and $\rho$ be a reductive maximal representation of
$\Gamma$ in $\mathrm{E}_{6(-14)}$. Then there exists a holomorphic or antiholomorphic $\rho$-equivariant
map from $\bH^n_\bC$ to the symmetric space $M$
associated with $\mathrm{E}_{6(-14)}$.
\end{theo}

\begin{proof}
As explained at the beginning of Section~\ref{sec:MW}, we may and do assume that
$\tau(\rho)>0$.
By Proposition~\ref{prop:vanishing}, $\ovg$ vanishes on the regular locus ${\cR}(\ovb)$ of $\ovb$. By
Proposition~\ref{prop:equal}~(2), ${\cR}(\ovb)$ is dense in $X$, so that $\ovg$ vanishes
identically on $X$. This means that the $\rho$-equivariant harmonic map $f:\HC\to M$ used to
define the Higgs bundle $(\ovE,\ovt)$ is holomorphic.
\end{proof}

\subsection{Proof of the main results}
\label{subsec:proof}

In this subsection, we give detailed proofs of Theorem~\ref{maintheo} and Corollary~\ref{maincoro}
stated in the introduction,
although some of the arguments might be well-known.

So let $n\geq 2$, $\Gamma\subset\SU(1,n)$ be a uniform lattice and $\rho$ a maximal
representation of~$\Gamma$ in a real algebraic Hermitian Lie group $G_\bR$.
Recall that $G_\bR$ is semisimple
with no compact factors and that $G_\bR$ is the connected component
$\bfG(\bR)^\circ$ of the group of real point of a real algebraic group $\bfG$.

At first we will assume that $\Gamma$ is torsion free and, as before, that $G = \bfG(\bC)$ is
simply connected. Remarks~\ref{rema:1-connected} and~\ref{rema:1-connected-bis} explain how to deal
with non simply connected $G$'s and Remark~\ref{rema:torsion} with the presence of torsion.

We may moreover assume that $G_\bR$ is simple,
because a representation in a product of simple groups
is maximal if and only if the induced representations in the simple factors are maximal.

Furthermore, we may and will assume that $\rho$ is
reductive, since by~\citestar{BIW}*{Cor.\,4}, maximal representations are always
reductive. Alternatively, one may proceed as in the following, but first with
reductive maximal representations and prove Theorem~\ref{maintheo} and Corollary~\ref{maincoro} for
them. Once the compactness of the centralizer $Z_\bR$ for reductive maximal representations has been
proved, one can prove that all maximal representations are reductive as
in~\citestar{KM}*{\textsection 4.5}.

Finally, as explained at the beginning of Section~\ref{sec:MW}, we may and do assume that $\tau(\rho)>0$.

Let then $f: \HC \to M =G_\bR/K_\bR$ be a harmonic $\rho$-equivariant map
(such a map exists by \cite{CorletteFlatGBundles}).
By Proposition~\ref{prop:tube}, $G_\bR$ is not of tube type, hence $G_\bR$ is isogenous to
either $\SU(p,q)$ with $p<q$, $\SO^\star(2m)$ with $m\geq 5$ odd, or $\mathrm{E}_{6(-14)}$ (whose
real ranks are respectively $p$, $\psfrac{m-1}2$, and $2$). Then
\citestar{KM}*{Th.\,4.1 \& Prop.\,5.6} and Theorem~\ref{thm:E6} imply that $f$ is holomorphic.
By the Ahlfors-Schwarz lemma (\cf \cite{Royden}), since the holomorphic sectional
curvature is $-1$ on $\HC$ and bounded above by $-\sfrac 1\rkm$ on $M$, we have the pointwise
inequality $f^\star \omega_M\leq \rkm \,\omega$. The maximality of $\rho$ then implies that
$f^\star\omega_M=\rkm\,\omega$.
Since there is equality in the Ahlfors-Schwarz lemma, $f$ is totally geodesic (see \eg \cite{Royden}).

Our results so far and \citestar{KM}*{Prop.\,4.10 \& 5.7} also imply that~$G_\bR$ is
either $\SU(p,q)$ with $q\geq pn$ (in which maximal representations exist indeed) or $\mathrm{E}_{6(-14)}$,
in particular~ $G_\bR$ is not isogenous to $\SO^\star(2m)$ with $m$ odd. In the $\SU(p,q)$ case,
we also know from~\citestar{KMrank2}*{\textsection 3.1} (see also~\citestar{KM}*{Prop.\,4.10})
that a holomorphic map $f:\HC\to M$ such that $f^\star\omega_M=\rkm\,\omega$ is unique up to composition by
an element of $\SU(p,q)$ and equivariant w.r.t a morphism
$\varphi:\SU(1,n)\to\SU(p,q)$.

We need another argument to handle the $\mathrm{E}_{6(-14)}$ case.
Following \citestar{BIW}*{p.\,690} and \citestar{Hamlet}*{p.640}, a totally geodesic map $g:\HC \to M$
is called \emph{tight} if
$\sup_{\Delta \subset \HC} \int_\Delta g^* \omega_M= \sup_{\Delta_M \subset M} \int_{\Delta_M}
\omega_M$, where $\Delta$, \resp $\Delta_M$, denote a geodesic triangle in $\HC$, \resp $M$
(recall that $\omega_M$ is the Kähler form on $M$).
It follows from the fact that our map $f$ is holomorphic and such that $f^\star\omega_M=\rkm\,\omega$
that it is tight.
Tight holomorphic maps between Hermitian
symmetric spaces were classified in~\cite{Hamlet} using \cite{ihara}.
As we assumed that $G_\bR$ is simple, $M$ is irreducible. In this case, and since $n\geq 2$, tight
holomorphic totally geodesic maps $\HC\to M$ exist if and only
$G_\bR=\SU(p,q)$ with $q\geq pn$ or $G_\bR=\mathrm{E}_{6(-14)}$ if $n=2$.
Moreover if $G_\bR=\mathrm{E}_{6(-14)}$ they are also
deduced one from another by composition by an element of $G_\bR$.

\begin{rema}
There is a small inaccuracy in~\cite{Hamlet}, where it is said that there are
two ``tight regular'' (in the terminology of this paper) maximal noncompact subalgebras of
${\fe}_{6(-14)}$ in which $\fsu(1,2)$ can be tightly embedded. In fact, only $\fsu(2,4)\subset
{\fe}_{6(-14)}$ is a maximal subalgebra for
these properties. This was confirmed to us by the
author.
\end{rema}

We assume that $G_\bR \simeq \mathrm{E}_{6(-14)}$ and
we now show that there is a morphism of groups $\varphi:\SU(1,2) \to G_\bR$ such that $f$ is
$\varphi$-equivariant. The application $f$ is given by a morphism
of Lie algebras $\fsu(1,2) \to \fsu(2,4) \to \fe_{6(-14)}$.
As explained in \citestar{ihara}*{\textsection 4.5},
the inclusion $\fsu(2,4) \to \fe_{6(-14)}$ is the
inclusion of the Lie algebra $\fu$ whose root system is generated by
$\beta_1,\beta_3,\beta_4,\beta_2,\gamma$, where the $\beta_i$'s are the simple roots
(we use Bourbaki's notations for the simple roots,
whereas Ihara has different notations) and
$\gamma=\beta_2+\beta_3+2\beta_4+2\beta_5+\beta_6$.

To show the existence of our morphism $\varphi$, it is enough to show that the
subgroup $U_\bR \subset G_\bR$ whose Lie algebra is $\fu$ is isomorphic to $\SU(2,4)$.
Since, by \citestar{ihara}*{\textsection 4.5}, $\fu \simeq \fsu(2,4)$, it is enough to show that
the complexification $U$ of $U_\bR$ is simply connected. To prove this, we show that for $T_U \subset U$
a maximal torus, the characters of $T_U$ are all the weights of the root system of $U$.

Since $G$ is simply connected, all the fundamental weights of its root system define characters of $T$, and
those characters restrict to some characters of $T_U$. It is thus enough to show that the fundamental weights
for $U$ can be expressed as linear combinations of the restrictions of the fundamental weights for $G$.
The fundamental weights corresponding to the roots $\beta_1,\beta_3,\beta_4,\beta_2,\gamma$ are the restriction
of the weights $\varpi_1,\varpi_3-2\varpi_6,\varpi_4-3\varpi_6,\varpi_2-2\varpi_6,\varpi_6$
respectively. This proves that $U_\bR \simeq \SU(2,4)$.

This proves all the assertions of Theorem~\ref{maintheo}
except the uniqueness of the harmonic map $\HC\to M$ that is
$\rho$-equivariant. To prove
it, we need to have a closer look at $f$. We know that
$f$ is equivariant with respect to a morphism of
Lie groups $\varphi: \SU(1,n) \to G_\bR$ and that up to conjugacy of $\rho$, we may assume that $f$
and $\varphi$ are as follows:
\begin{itemize}
\item for $G_\bR=\SU(p,q)$ with $q\geq pn$, by~\citestar{KMrank2}*{\textsection 3.1} (see
also~\citestar{KM}*{Prop.\,4.10}, or~\cite{Hamlet}), $\varphi$ is the composition
$$
\SU(1,n)\Hto{\diag}
\SU(1,n)^p
\hto\SU(p,pn)\hto
\SU(p,q)\,;
$$

\item for $G_\bR=\mathrm{E}_{6(-14)}$ and $n=2$, by~\cite{Hamlet} and the discussion above,
$\varphi$ is the composition
$$
\SU(1,2) \Hto{\diag}
\SU(1,2)^2
\hto\SU(2,4)\hto \mathrm{E}_{6(-14)}.
$$
\end{itemize}

In both cases, the image $N$ of $f$ in $M=G_\bR/K_\bR$ is the orbit of $o=K_\bR$ under
$H_\bR:=\varphi(\SU(1,n))\subset G_\bR$.

We now describe the centralizer $Z_\bR$ of $H_\bR$ in $G_\bR$.
In case $G_\bR = \SU(p,q)$, let $C_\bR$ denote the group $\rU(p) \times \rU(q-pn)$ and let
$\chi: C_\bR \to \rU(1)$ be the character defined by $\chi(x,y)=\det(x)^{n+1} \cdot \det(y)$.
In case $G_\bR = \mathrm{E}_{6(-14)}$, let $C_\bR = \rU(2) \times \rU(2)$ and let $\chi: C_\bR \to \rU(1)$
be the character defined by $\chi(x,y)=\det(x)^{21} \cdot \det(y)^6$. Then:

\begin{lemm}
\label{lemm:Z}
The centralizer $Z_\bR$ of $H_\bR$ in $G_\bR$ is a subgroup of $K_\bR$ (hence it is compact). It is
isomorphic to the kernel of $\chi$ in $C_\bR$.
\end{lemm}
\begin{proof}
In the case of $\SU(p,q)$, the description of $\varphi$ given above shows that the standard
representation $\bC^{p+q}$ of $\SU(p,q)$, when seen as a representation of $\SU(1,n)$ via
$\varphi$, splits as
$$
\bC^{p+q}=\bC^{p+pn}\oplus\bC^r=\bC^{1+n}\otimes\bC^p\oplus\bC^r,
$$
where $\bC^{1+n}$ is the standard representation of $\SU(1,n)$, $r=q-pn$, and
$\bC^p$ and $\bC^r$ are trivial representations of $\SU(1,n)$.
To conclude, we argue as follows. Let \hbox{$g \in Z_\bR$}. Then~$g$ yields an endomorphism
of the $H_\bR$-module $\bC^{1+n}\otimes\bC^p\oplus\bC^r$. Since by Schur's lemma
such an endomorphism will
preserve isotypic factors, we see that
$g$ must preserve the factors $\bC^{1+n}\otimes\bC^p$ and $\bC^r$. Moreover its
restriction to $\bC^{1+n}\otimes\bC^p$ belongs to
$\rU(p,pn)$ and so by Schur lemma again
it must act by an element of $\rU(p)$, so that it belongs to $C_\bR$. Now, the determinant of an element
$(x,y) \in \rU(p) \times \rU(q-pn)$ considered as an element in $\rU(\bC^{1+n}\otimes\bC^p\oplus\bC^r)$ is
$\det(x)^{n+1} \cdot \det(y)$, which implies that $\chi(x,y)=1$.

In the case of $\mathrm{E}_{6(-14)}$, we use a model given
by Manivel in \citestar{manivel}*{Ex.\,3 p.\,464} of the $27$-dimensional representation $\bE$.
There is a subgroup in $\mathrm{E}_{6(-14)}$ isomorphic to $\SU(2,4) \times \SU(2)$ and $\bE$ splits as
$\wedge^2 \bU \oplus \bU \otimes \bA$, where
$\bU$ (\resp $\bA$) is the natural representations of $\SU(2,4)$ (\resp $\SU(2)$) of complex dimension $6$ (\resp $2$).
Here we restrict further to $\SU(1,2) \times \SU(2)$, where the first factor $\SU(1,2)$ is
diagonally embedded in $\SU(2,4)$, meaning that the representation
$\bU$ splits as $\bV \otimes \bB$, with $\bV$ the natural $3$-dimensional representation of
$\SU(1,2)$ and $\bB$
the trivial $2$-dimensional representation. We get
$$
\bE \simeq \wedge^2(\bV \otimes \bB) \oplus \bV \otimes \bA \otimes \bB
\simeq \wedge^2 \bV \otimes S^2 \bB \oplus S^2 \bV \otimes \wedge^2 \bB \oplus \bV \otimes \bA \otimes \bB.
$$
As in the case of $\SU(p,q)$, an element $g$ in the centralizer of $H_\bR$ will yield an
$H_\bR$\nobreakdash-equivariant endomorphism, and will preserve each of these factors, so that it belongs to
$\GL(\bA) \times \GL(\bB)$.
Since it is an element of the group $E_{6(-14)}$, one sees that it must be given by an
element in $\rU(\bA) \times \rU(\bB)$.

The computation of the character $\chi$ is done as follows. If $f=(x,y) \in \rU(\bA) \times \rU(\bB)$,
then the determinant of the action of $f$ on
$\bE$ is the product of the determinants of the actions of $f$ on $\wedge^2(\bV \otimes \bB)$ and on
$\bV \otimes \bA \otimes \bB$.
The action on $\bV \otimes \bA \otimes \bB$ has determinant $\det(x)^6\det(y)^6$, and the action on
$\wedge^2(\bV \otimes \bB)$
has determinant $\det(y)^{15}$.
\end{proof}

\begin{rema}
\label{rema:compactness}
The compactness of $Z_\bR$ is proved in greater generality in \citestar{BIW}*{Cor.\,4}.
\end{rema}

\begin{lemm} \label{lemm:fixator}
The pointwise stabilizer of $N=f(\HC)$ in $G_\bR$ is exactly $Z_\bR$. The stabilizer of $N$ in $G_\bR$ is the
almost direct product $H_\bR\cdot Z_\bR$.
\end{lemm}
\begin{proof}
Let $o=K_\bR \in N$ be the base point of $M$.
Let us denote by $\Fix(N) \subset G_\bR$ the subgroup of elements which fix all the elements in $N$.
We want to prove that $\Fix(N)=Z_\bR$. We have an inclusion $Z_\bR \subset \Fix(N)$. Indeed, if $h \in H_\bR$ and
$z \in Z_\bR$, then $z\cdot o = o$ since $Z_\bR\subset K_\bR$. Thus, since $z$ and $h$ commute, $z
\cdot (h \cdot o) = h \cdot (z \cdot o) = h \cdot o$.

The subgroup $H_\bR$ may be defined referring only to $N$ as follows.
Let $\fgg_\bR = \fk_\bR \oplus \fm_\bR$ be the Cartan decomposition of $\fgg_\bR$. The tangent space
$T_oN$ identifies with a subspace of $\fm_\bR$ that we denote by $\fn_\bR$.
The space $\fn_\bR$ defines a Lie triple system, so that $\fh_\bR:= [\fn_\bR,\fn_\bR] \oplus \fn_\bR
\subset \fgg_\bR$ is a Lie subalgebra. Then, $H_\bR$ is the connected Lie group of $G_\bR$
with Lie algebra $\fh_\bR$.

For the reverse inclusion we need to prove that $\Fix(N) \subset Z_\bR$.
It follows from the given description of $H_\bR$ that $H_\bR$ is normalized by $\Fix(N)$.
Let $g \in \Fix(N)$ and $h \in H_\bR$. Then the commutator
$ghg^{-1}h^{-1}$ belongs to $H_\bR$ and acts trivially on $N$. Thus, it belongs to the center of
$H_\bR$. Since this center is finite, the connectedness of~$H_\bR$ implies that $ghg^{-1}h^{-1}$ is the neutral element. Therefore $g$ centralizes $H_\bR$, so
it belongs to $Z_\bR$.

Since the automorphism group of $N$ is isogenous to $H_\bR$, the second assertion of the Lemma
follows from the first.
\end{proof}

\begin{rema}
Lemmas~\ref{lemm:Z} and~\ref{lemm:fixator} fix a mistake in the description given
in~\citestar{KMrank2}*{\textsection 3.1.2} of the stabilizer of $N$ in the $\SU(p,q)$ case.
\end{rema}

\begin{proof}[Proof of Corollary~\ref{maincoro}]
The facts that $\rho$ is discrete and
faithful and that $\rho(\Gamma)$ acts cocompactly on $N$ follow from the $\rho$-equivariance of the totally geodesic
embedding~$f$. The reductivity of $\rho$ has been already asserted and the compactness of $Z_\bR$
was established in Lemma~\ref{lemm:Z}.
Now, given $\gamma \in \Gamma$, the equivariance of $f$ w.r.t. $\rho$ and $\varphi$ means that
$\rho(\gamma)$ and $\varphi(\gamma)$ have the same action on $N$.
We let $\rhoc(\gamma)=\rho(\gamma)\varphi(\gamma)^{-1}$. This is an element of the pointwise stabilizer of $N$,
which is equal to the centralizer $Z_\bR$ of $H_\bR$ by Lemma~\ref{lemm:fixator}. Since
$\varphi(\gamma) \in H_\bR$ by definition of $\varphi$, the elements $\varphi(\gamma)$ and
$\rhoc(\gamma)$ commute.
It follows that $\varphi(\gamma)$ and $\rho(\gamma)$ commute, and that $\rhoc$ is a morphism of groups.
\end{proof}

\begin{proof}[Proof of the uniqueness of $f$]
By the uniqueness statement for tight holomorphic
\hbox{totally} geodesic maps $\HC\to M$, we know that if $f':\HC\to M$ is
another $\rho$-equivariant harmonic map, then there exists $g\in G_\bR$ such that $f'=g\circ f$. By
$\rho$-equivariance of $f$ and $f'$, we have that
$$
\rho(\gamma)\circ g(x)=g\circ\rho(\gamma)(x),\quad\forall \gamma\in\Gamma\mbox{ and }\forall x\in N.
$$
It follows that $g \cdot N$ is $\rho(\Gamma)$-stable.
Thus the map
\begin{align*}
d_{g\cdot N}:N&\to\bR\\ x&\mto \inf_{y \in g \cdot N} d(x,y),
\end{align*}
where
$d$ denotes the distance in
$M$, is invariant under the cocompact action of $\rho(\Gamma)$ on~$N$. It is therefore
bounded. Since it is moreover convex (\citestar{BH}*{p.\,178}), it is constant, equal to $a$, say. In the same way, the map
\begin{align*}
d_N:g\cdot N&\to\bR\\
x&\mto \inf_{y \in N} d(x,y)
\end{align*}
is also constant equal to $a$.

If $a>0$ it follows from the sandwich
lemma (\citestar{BH}*{p.\,182}) that the convex hull of $N\cup g\cdot N$ in $M$ is isometric to the product $N\times
[0,a]$. This implies that there exists a tangent vector $v\in T_oM\simeq \fm_\bR$, orthogonal to
$T_oN\simeq \fn_\bR$ such that $[v,u]=0$ for all $u\in\fn_\bR$. Indeed the
norm (for the Killing form) of $[v,u]\in \fgg_\bR$ is up to a constant the sectional curvature of the
plane generated by the tangent vectors $u$ and $ v$, which is $0$ since they belong to different factors
of a Riemannian product.
In this case the 1-parameter group of transvections along
the geodesic defined by $v$ is included in the centralizer $Z_\bR$ of $H_\bR$, a contradiction since
$Z_\bR$ is compact.

Hence $a=0$ and $g\cdot N=N$. Therefore there exist $h\in H_\bR$ and $z\in Z_\bR$ such
that $g=hz=zh$. The above commutation relation between $\rho(\gamma)=\varphi(\gamma)\rhoc(\gamma)$ and~$g$ on $N$
means that $\rho(\gamma)g\rho(\gamma)^{-1}g^{-1}$ fixes $N$ pointwise and hence belongs to $Z_\bR$ by
Lemma~\ref{lemm:fixator}. Hence for
all $\gamma\in\Gamma$ we obtain that $\varphi(\gamma)h\varphi(\gamma)^{-1}h^{-1}$ belongs to $Z_\bR\cap H_\bR$
(recall that $\rhoc(\gamma) \in Z_\bR$). Now $\Gamma$
is Zariski dense in $\SU(1,n)$ by the Borel density theorem and we deduce that
$\varphi(x)h\varphi(x)^{-1}h^{-1}\in Z_\bR\cap H_\bR$ for all $x\in\SU(1,n)$.
Since $Z_\bR\cap H_\bR$ is finite and
$\SU(1,n)$ is connected, $h\in Z_\bR$. Therefore $g\in Z_\bR$ and $f'=f$.
\end{proof}

\begin{proof}[Proof of the uniqueness of $\varphi$]
Since $f$ is well-defined, it makes sense to claim that there is a unique morphism $\varphi:\SU(1,n) \to G_\bR$
such that $f$ is equivariant with respect to $\varphi$. We prove that this statement holds: let
$\varphi_1,\varphi_2$ be two
such morphisms. We get by equivariance, for $x \in \HC$ and $g \in \SU(1,n)$:
$$
f(g \cdot x) = \varphi_1(g) \cdot f(x) = \varphi_2(g) \cdot f(x).
$$
Thus, setting $\psi(g)=\varphi_2(g)^{-1}\varphi_1(g)$, $\psi(g)$ is in the pointwise stabilizer of $N$. It
follows from Lemma~\ref{lemm:fixator} that $\psi(g)$ commutes with $\varphi_1(g')$
and $\varphi_2(g')$, for all $g'\in\SU(1,n)$. Thus, $\psi:\SU(1,n) \to G_\bR$ is a morphism of groups, and
$\psi(g_1)\psi(g_2)=\psi(g_2)\psi(g_1)$ for all $g_1,g_2$ in $\SU(1,n)$. It follows from the fact
that $\SU(1,n)$ is simple that $\psi$ is trivial.
\end{proof}

\begin{rema}
\label{rema:1-connected}
If we drop the assumption that the group of complex points \hbox{$G=\bfG(\bC)$} of the algebraic
group $\bfG$ is simply connected, then $\bE$ might no longer be a
representation of $G$ and our constructions cannot be made. However, in this case, letting $\tilde
G$ be the simply connected cover of $G$ and $\bE$ the cominuscule representation
of~$\tilde G$ that we have been considering, there is an integer $k$ such that $\bE^{\otimes k}$ is
a representation of $G$. The arguments given in the article can be adapted with the representation~$\bE^{\otimes k}$ instead of $\bE$, and the main results (Theorem~\ref{theo:MW}, Theorem~\ref{maintheo},
Corollary~\ref{maincoro} and Proposition~\ref{mainprop})
remain true without the simple connectedness assumption.
\end{rema}

\begin{proof}[Proof of Remark~\ref{rema:1-connected}]
We denote by $\Zmax$ the maximum value $\scal{\Lambda,z}$ for a weight $\Lambda$ in $X(\Ek)$.
Since such a weight is of the form $\chi_1 + \cdots + \chi_k$ for some weights $\chi_i \in X(\bE)$, we have
$\Zmax = k \zmax$. Similarly the highest value of $\scal{\Lambda,h}$ for $\Lambda \in X(\Ek)$ is $kr$.
We define as in Definition~\ref{defi:Ei} and~\ref{defi:V} the subspaces $\Ek_i$ and
$\Vkr_i$ by the relations
$$
\Ek_i = \bigoplus_{\Lambda: \scal{\Lambda,z} = k\zmax-2i} \Ek_\Lambda \quad\text{and}\quad
\Vkr_i = \bigoplus_{\substack{\Lambda \in X(\Ek_i):\\ \scal{\Lambda,h}=kr-2i}} \bE_\chi.$$
We set $\Vkr = \bigoplus_i \Vkr_i$.
If $\Lambda = \chi_1 + \cdots + \chi_k \in X(\Ek_i)$, then for all $l$ we have
$\chi_l \in \bE_{j_l}$ with integers $j_l$ such that $i=j_1 + \cdots + j_k$. We have the inequality
$\scal{\chi_l,h} \geq r-2j_l$ by Proposition~\ref{prop:chi-Ei}. Thus, if $\Ek_\Lambda \subset \Vkr_i$, this implies
that for each $l$, $\bE_{\chi_l} \subset \Vr_l$. We deduce that
$\Vkr = (\Vr)^{\otimes k}$. It follows that at the level of sheaves, we will be able to
define a subsheaf of $\cEk$ by the same trick as in~\eqref{equa:filtration-y}.
Since passing from $\bE$ to $\Ek$ and from $\Vr$ to $(\Vr)^{\otimes k}$ just multiplies the slopes by $k$,
we get the same Milnor-Wood inequality. In the case of equality, the arguments of
Section~\ref{subsec:E6} are still valid and we get that $\gamma=0$,
proving that the map $f$ is holomorphic.
\end{proof}

\begin{rema}\label{rema:1-connected-bis}
Alternatively, it is possible to deal with non simply connected groups~$G$ as follows. Let
$\widetilde \bfG$
be the universal cover of $\bfG$, $\widetilde G_\bR=\widetilde \bfG(\bR)^\circ$, and
$\widetilde G=\widetilde \bfG(\bC)$. Then $\widetilde G$ is simply connected ($\widetilde G_\bR$
might not be). By Selberg's lemma,
there is a finite index sublattice $\Gamma' \subset \Gamma$ such that
the representation $\rho:\Gamma \to G_\bR$ restricted to $\Gamma'$ lifts to $\widetilde G_\bR$, see \eg \citestar{KMgeneraltype}*{Lem.\,2.2}: we have
$\rho': \Gamma' \to \widetilde G_\bR$. Since $\tau(\rho') = |\Gamma/\Gamma'|\, \tau(\rho)$ and
$\vol(\Gamma'\backslash\HC)=|\Gamma/\Gamma'|\vol(\Gamma\backslash\HC)$, we may use our results
for $\rho'$ and deduce the Milnor-Wood inequality for $\rho$, which is maximal if and only if $\rho'$
is. Moreover, in the maximal case, we have by \cite{CorletteFlatGBundles}
a $\rho$-equivariant harmonic map $f:\HC \to M$ and, by our results, a unique
$\rho'$-equivariant harmonic map $f':\HC \to M$. Since $f$ is also $\rho'$-equivariant, by uniqueness, we have $f=f'$
and our results follow.
\end{rema}

\begin{rema}\label{rema:torsion}
If $\Gamma$ has torsion, then by
Selberg's Lemma (see \citestar{ratcliffe}*{p.\,327} or \cite{selberg}) there exists a finite index
torsion free normal subgroup $\Gamma'$ of $\Gamma$. Let $\rho'$ denotes the
restriction of $\rho$ to $\Gamma'$. Since $\tau(\rho') = |\Gamma/\Gamma'|\, \tau(\rho)$ and
$\vol(\Gamma'\backslash\HC)=|\Gamma/\Gamma'|\vol(\Gamma\backslash\HC)$, $\rho:\Gamma\to G_\bR$ is
maximal if and only if $\rho':\Gamma'\to G_\bR$ is. Then, in view of the foregoing,
Theorem~\ref{maintheo} and Corollary~\ref{maincoro} hold for $\rho'$. Therefore there exist a unique
$\rho'$-equivariant map $f':\HC\to M$. Now if $\gamma\in\Gamma$, the map $f'_\gamma:=\rho(\gamma)^{-1}\circ f'\circ\gamma$ is
also harmonic and the fact that $\Gamma'$ is normal in $\Gamma$ implies that it is also
$\rho'$-equivariant.
Hence $f'=f'_\gamma$ for all $\gamma\in\Gamma$, \ie $f'$ is in fact
$\rho$-equivariant and Theorem~\ref{maintheo} and Corollary~\ref{maincoro} also hold for~$\rho$.
\end{rema}
\egroup

\backmatter
\bibliographystyle{jepalpha}
\bibliography{chaput-maubon}
\end{document}
